% Options for packages loaded elsewhere
\PassOptionsToPackage{unicode}{hyperref}
\PassOptionsToPackage{hyphens}{url}
\PassOptionsToPackage{dvipsnames,svgnames,x11names}{xcolor}
%
\documentclass[
  11pt,
]{article}
\usepackage{amsmath,amssymb}
\usepackage{iftex}
\ifPDFTeX
  \usepackage[T1]{fontenc}
  \usepackage[utf8]{inputenc}
  \usepackage{textcomp} % provide euro and other symbols
\else % if luatex or xetex
  \usepackage{unicode-math} % this also loads fontspec
  \defaultfontfeatures{Scale=MatchLowercase}
  \defaultfontfeatures[\rmfamily]{Ligatures=TeX,Scale=1}
\fi
\usepackage{lmodern}
\ifPDFTeX\else
  % xetex/luatex font selection
  \setmainfont[]{Latin Modern Roman}
  \setsansfont[]{Latin Modern Sans}
  \setmonofont[]{DejaVu Sans Mono}
\fi
% Use upquote if available, for straight quotes in verbatim environments
\IfFileExists{upquote.sty}{\usepackage{upquote}}{}
\IfFileExists{microtype.sty}{% use microtype if available
  \usepackage[]{microtype}
  \UseMicrotypeSet[protrusion]{basicmath} % disable protrusion for tt fonts
}{}
\makeatletter
\@ifundefined{KOMAClassName}{% if non-KOMA class
  \IfFileExists{parskip.sty}{%
    \usepackage{parskip}
  }{% else
    \setlength{\parindent}{0pt}
    \setlength{\parskip}{6pt plus 2pt minus 1pt}}
}{% if KOMA class
  \KOMAoptions{parskip=half}}
\makeatother
\usepackage{xcolor}
\usepackage[margin=24mm]{geometry}
\usepackage{longtable,booktabs,array}
\usepackage{calc} % for calculating minipage widths
% Correct order of tables after \paragraph or \subparagraph
\usepackage{etoolbox}
\makeatletter
\patchcmd\longtable{\par}{\if@noskipsec\mbox{}\fi\par}{}{}
\makeatother
% Allow footnotes in longtable head/foot
\IfFileExists{footnotehyper.sty}{\usepackage{footnotehyper}}{\usepackage{footnote}}
\makesavenoteenv{longtable}
\usepackage{graphicx}
\makeatletter
\def\maxwidth{\ifdim\Gin@nat@width>\linewidth\linewidth\else\Gin@nat@width\fi}
\def\maxheight{\ifdim\Gin@nat@height>\textheight\textheight\else\Gin@nat@height\fi}
\makeatother
% Scale images if necessary, so that they will not overflow the page
% margins by default, and it is still possible to overwrite the defaults
% using explicit options in \includegraphics[width, height, ...]{}
\setkeys{Gin}{width=\maxwidth,height=\maxheight,keepaspectratio}
% Set default figure placement to htbp
\makeatletter
\def\fps@figure{htbp}
\makeatother
\setlength{\emergencystretch}{3em} % prevent overfull lines
\providecommand{\tightlist}{%
  \setlength{\itemsep}{0pt}\setlength{\parskip}{0pt}}
\setcounter{secnumdepth}{5}
% Chinese names/titles are short and explicitly wrapped by the build script.
% This avoids requiring the optional xeCJK package in a minimal TeX install.
\newfontfamily\cjkfont{Noto Serif CJK SC}
\newcommand{\CJKtext}[1]{{\cjkfont #1}}
\usepackage{amssymb,booktabs,longtable,array}
\usepackage{microtype}
\usepackage{fancyhdr}
\pagestyle{fancy}
\fancyhf{}
\fancyhead[L]{\small Algebra as a structural toolkit}
\fancyhead[R]{\small\thepage}
\renewcommand{\headrulewidth}{0.3pt}
\setlength{\headheight}{15pt}
\setlength{\emergencystretch}{3em}
\setlength{\tabcolsep}{4pt}
\AtBeginEnvironment{longtable}{\small}
\makeatletter
\renewcommand*\l@subsection{\@dottedtocline{2}{1.5em}{3.4em}}
\makeatother
\let\originaltableofcontents\tableofcontents
\renewcommand{\tableofcontents}{\originaltableofcontents\clearpage}
\ifLuaTeX
  \usepackage{selnolig}  % disable illegal ligatures
\fi
\usepackage{bookmark}
\IfFileExists{xurl.sty}{\usepackage{xurl}}{} % add URL line breaks if available
\urlstyle{same}
\hypersetup{
  pdftitle={Algebra as a structural toolkit},
  pdfauthor={Research notes for returning to algebra},
  colorlinks=true,
  linkcolor={blue!45!black},
  filecolor={Maroon},
  citecolor={Blue},
  urlcolor={blue!55!black},
  pdfcreator={LaTeX via pandoc}}

\title{Algebra as a structural toolkit}
\usepackage{etoolbox}
\makeatletter
\providecommand{\subtitle}[1]{% add subtitle to \maketitle
  \apptocmd{\@title}{\par {\large #1 \par}}{}{}
}
\makeatother
\subtitle{Foundations, splitting obstructions, integral reconstruction,
and proof comparisons}
\author{Research notes for returning to algebra}
\date{20 October 2026}

\begin{document}
\maketitle

{
\hypersetup{linkcolor=}
\setcounter{tocdepth}{2}
\tableofcontents
}
\section*{Returning to algebra: a compact working
guide}\label{returning-to-algebra-a-compact-working-guide}
\addcontentsline{toc}{section}{Returning to algebra: a compact working
guide}

For the continuation, use the
\href{RECONSTRUCTION-MAP.md}{reconstructive map} after these
foundational exercises. It connects invariant-complement obstructions,
averaging, integral spectral reconstruction, and common sign lifts, with
explicit hypotheses and stopping tests. The two expanded proof
comparisons identify what can actually be shortened in the inspected
parts of Xi, Hungerford, and Grillet.

\subsection*{The essential memory
spine}\label{the-essential-memory-spine}
\addcontentsline{toc}{subsection}{The essential memory spine}

Retain the following statements with their hypotheses and the type of
every map. The rest of the toolkit should be reconstructible from them.

\begin{enumerate}
\def\labelenumi{\arabic{enumi}.}
\tightlist
\item
  A map \(V/U\to X\) is a map \(V\to X\) that kills \(U\).
\item
  Dualization is precomposition; evaluation gives \(V\cong V^{**}\) in
  finite dimension. An identification with \(V^*\) needs extra data.
\item
  A direct sum is encoded by projections, and \(P^2=P\) gives
  \(V=\operatorname{im}P\oplus\ker P\).
\item
  A split exact sequence has a section. Sections differ by a map from
  the quotient to the kernel; existence does not imply a preferred
  choice.
\item
  An operator is a \(k[t]\)-module. The PID structure theorem classifies
  finite-dimensional operators through invariant factors; Jordan form
  needs splitting of the minimal polynomial.
\item
  Coprime annihilator factors produce polynomial idempotents by Bézout.
\item
  Tensor products represent bilinear or balanced maps. A formula on pure
  tensors must respect the defining relations.
\item
  \(W\otimes V^*\) identifies canonically with the finite-rank maps
  \(V\to W\), hence with all of \(\operatorname{Hom}(V,W)\) if either
  \(V\) or \(W\) is finite-dimensional. Trace is evaluation in finite
  dimension.
\item
  Tensor, symmetric and exterior algebras are free constructions with
  specified relations. \(\Lambda^r\) transports alternating data; the
  top exterior power gives a determinant line.
\item
  Forms transform by congruence, operators by similarity. A positive
  metric supplies adjoints, orthogonal projections, and the normal
  spectral theorem.
\item
  Exactness is image equals kernel. Over a field short exact sequences
  split and positive Ext/Tor vanish; homology of an arbitrary complex
  need not.
\item
  A representation is a module over a group algebra. Maschke needs
  invertible group order; Schur's scalar conclusion needs finite
  dimension and an algebraically closed field.
\end{enumerate}

For Hadamard/Williamson add five facts: Gram orthogonality is an
integral identity; tensors multiply valid orders; cyclic Fourier
coordinates are CRT coordinates; compression loses lifting information;
and equality of Boolean phases must be distinguished from independence
of their signs.

\subsection*{A selective route, with stopping
tests}\label{a-selective-route-with-stopping-tests}
\addcontentsline{toc}{subsection}{A selective route, with stopping
tests}

\textbf{First return: regain the language.} Read the
quotient/duality/splitting chapter, with Xi II §§1.4--1.6 and 2.4 beside
it. Stop when you can construct the quotient map and its dual, describe
the affine space of sections, and explain why an invariant complement is
an additional demand. Xi I can be a reference for matrix maps and
determinant characterizations rather than a sequential restart.

\textbf{Second return: recover operator theory from one theorem.} Choose
Grillet VIII.6--7 or Hungerford IV.6 then VII.4--5 as the principal
module route. Read Xi II §2.5's actual Bézout and quotient-induction
proofs as a comparison, not a second complete linear algebra course.
Stop when you can obtain primary projections, invariant factors, Jordan
blocks under the field hypothesis, and the block counts from kernel
growth.

\textbf{Third return: turn multilinearity into maps.} Read Xi II §6.2,
especially Theorem 6.15 and the Hom/scalar-extension examples. Then use
Grillet XI.5--6 and XIII.2--4 for balanced products, quotient
constructions, and characteristic issues. Read Mac Lane III.1--2 and
IV.1 when the universal examples are already familiar. Stop when you can
define tensor-induced maps without a basis and derive Cauchy--Binet from
\(\Lambda^r(gf)\).

\textbf{Keep a metric branch active.} Read Xi II §§1.7, 3.3--3.5 and 3.8
as needed for forms, adjoints, normal operators, and projections. This
is central for Hadamard Gram matrices. Defer the rest of
affine/projective geometry only according to current goals, without
concluding that the content was replaced.

\textbf{Introduce homology at its first real obstruction.} Read the
exactness chapter here, then Hilton--Stammbach III.1--4 and IV.1--3 if
you want the general construction. Grillet X.1, X.3 and XII are an
alternative. Stop the first pass after computing the \(k[t]/(t^2)\)
extension and explaining nonzero field-complex homology. Spectral
sequences and derived categories are not prerequisites for the current
Hadamard bridge.

\textbf{For the confirmed problem direction.} Read the Hadamard bridge
in order: Gram matrices and tensor closure; Xi III §5.6 finite fields;
the character-sum proof and Paley blocks; cyclic group algebra and
Fourier CRT; then the pinned third-round proofs with their color and
Boolean equations. Use Xi III Chapter 3 or Grillet IX.7 for group
algebras and representations, keeping Grillet's initial
indecomposable/irreducible convention visible. Read the prior
limitations before proposing a larger search. The essential stopping
test is being able to say exactly which implication is necessary,
sufficient, or only a relaxation.

Hall's further Lie theory, Lang's local arithmetic, Hartshorne's
scheme/sheaf theory, and Hirsch's transversality/bundle development can
then be pursued as separate branches. Their prerequisite bridges are
recorded, not placed in front of every use of linear algebra.

\subsection*{Reconstruction exercises}\label{reconstruction-exercises}
\addcontentsline{toc}{subsection}{Reconstruction exercises}

These are deliberately small. Their purpose is to test whether a tool
can be rebuilt, not to begin the later difficult problem-set campaign.

\begin{enumerate}
\def\labelenumi{\arabic{enumi}.}
\tightlist
\item
  \textbf{Quotient versus complement.} For \(U\subseteq V\), construct
  \((V/U)^*\cong U^\circ\) and show why no complement is needed. Then
  describe all sections of \(V\to V/U\) once one has been chosen.
\item
  \textbf{One invariant line.} In \(k[t]/(t^2)\), list the subspace
  killed by \(t\). Prove it has no \(t\)-invariant complement, and
  identify the associated nonzero Ext class.
\item
  \textbf{Recover a projection.} Assuming \(T^2(T-I)^2=0\), derive the
  projector \(3T^2-2T^3\) from polynomial congruences. Check the
  argument in characteristics \(2\) and \(3\) without dividing by those
  integers.
\item
  \textbf{Recover Jordan data.} Let \(T:V\to V\) with \(\dim_kV=10\) and
  \(\lambda\in k\). The dimensions of \(\ker(T-\lambda1_V)^j\) for
  \(j=0,1,\ldots,5\) are \(0,4,7,9,10,10\). Find all its block sizes,
  minimal polynomial, and the dimension of its commutant.
\item
  \textbf{Nonsplitting field.} Classify multiplication by \(t\) on
  \(\mathbb R[t]/(t^2+1)\) as a real module and after complexification.
  Explain which use of ``semisimple'' is valid before extension.
\item
  \textbf{Tensor universality.} Construct the currying bijection and the
  canonical map \(W\otimes V^*\to\operatorname{Hom}(V,W)\). Identify the
  image of the latter and prove injectivity. Explain why finite
  dimension of either \(V\) or \(W\) suffices for surjectivity, and give
  a failure when both are infinite-dimensional.
\item
  \textbf{Determinants from wedges.} Derive the minors formula for
  \(\Lambda^2f\) and the \(2\times2\) Cauchy--Binet identity for a
  \(2\times3\) matrix times a \(3\times2\) matrix.
\item
  \textbf{Forms versus operators.} Give two congruent real symmetric
  matrices that are not similar. Derive \(H^{-1}A^*H\) from the defining
  adjoint identity, stating the Hermitian linearity convention.
\item
  \textbf{Homology survives.} For \(d:k^3\to k^2\), \(d(x,y,z)=(x,y)\),
  compute the two homology groups of the two-term complex, and write an
  explicit contractible/homology decomposition.
\item
  \textbf{The right group action.} Verify the inverse in the dual
  action, and differentiate it to obtain the minus sign in the Lie
  algebra dual action.
\item
  \textbf{Character sum.} Reproduce the count of \(y^2=x(x+a)\) over an
  odd finite field. Identify every use of \(a\ne0\) and \(2\ne0\), then
  check one Paley block multiplication.
\item
  \textbf{Compression and phase dependence.} Check the core-nine
  compressed autocorrelation example. Explain why failure of its chosen
  lifts alone cannot prove absence of all lifts. Separately give a
  numerical assignment to three formally independent signs that violates
  \(\xi_{p+q}=\xi_p\xi_q\).
\end{enumerate}

\subsection*{Brief answer checks}\label{brief-answer-checks}
\addcontentsline{toc}{subsection}{Brief answer checks}

For Exercise 1, pull back along the quotient, and use \(s'=s+i h\). For
2, the killed subspace is \(kt\) and the class is \(1\in k\) under
\(\operatorname{Ext}^1_{k[t]}(k,k)\cong k\). For 3, require a polynomial
to be zero modulo \(t^2\) and one modulo \((t-1)^2\); the integer
identity remains valid after any reduction of characteristic.

For 4, the differences are \(4,3,2,1,0\), so the sizes are \(4,3,2,1\).
Since the kernels reach \(\dim_kV=10\), there is no other primary
summand. The minimal polynomial is \((t-\lambda)^4\) and the commutant
dimension is \(\sum_{i,j}\min(i,j)=30\) for \(i,j=1,2,3,4\). For 5, the
real module is simple; after complexification it is the direct sum of
the \(i\) and \(-i\) eigenspaces. It is semisimple over \(\mathbb R[t]\)
but not real diagonalizable.

For 6, currying needs no finite dimension. The canonical Hom/tensor map
is always injective, with image exactly the finite-rank maps. Thus it is
onto precisely when every map \(V\to W\) has finite rank; finite
dimension of either space suffices. For \(V=W\) infinite-dimensional,
the identity is outside its image. For 7, sum the three intermediate
\(2\)-subsets. For 8, \(I\) and \(2I\) suffice, and
\(h(x,y)=\bar x^{\mathsf T}Hy\) is the convention used here. For 9,
\(H_1\) is the \(z\)-axis and \(H_0=0\); the \((x,y)\) coordinates map
isomorphically to \(k^2\). For 10, transpose reverses products and
differentiation of an inverse at the identity contributes a minus sign.
For 11, the count is \(q-1\). For 12, choose all three phase signs
negative: the product of the first two is positive, contradicting the
third. These checks indicate the target, while the chapters provide the
proofs to reconstruct.

\subsection*{What can be deferred}\label{what-can-be-deferred}
\addcontentsline{toc}{subsection}{What can be deferred}

Defer a formula once its map or universal property can be reconstructed.
Defer a subject when it has no role in the current destination. Do not
defer the hypotheses, the equivalence relation, or a known failure case.
For this collection, an effective memory sheet contains definitions and
four or five counterexamples, with links back to proofs, rather than a
long inventory of named theorems.

\section*{A reconstructive map: from decompositions to their
obstructions}\label{a-reconstructive-map-from-decompositions-to-their-obstructions}
\addcontentsline{toc}{section}{A reconstructive map: from decompositions
to their obstructions}

The basic study explains quotients, duality, operators as modules,
tensor universality, exterior algebra, forms, exactness, and
representations. The continuation asks a sharper question: \textbf{which
extra structure survives a decomposition, and how can its failure be
computed?} Four concrete outputs organize the new material:

\begin{itemize}
\tightlist
\item
  an invariant complement, or a reverse-intertwiner witness that it
  cannot exist;
\item
  an equivariant section obtained by a justified averaging operation;
\item
  an integral coefficient vector reconstructed from cyclotomic data, or
  a congruence obstruction;
\item
  a common sign lift of compressed rows, or a residual correlation
  defect.
\end{itemize}

These outputs require different hypotheses. A vector-space splitting is
not automatically invariant; a rational splitting is not automatically
integral; an integral vector is not automatically a sign vector; four
individually possible lifts need not satisfy their equations together.

\subsection*{Dependencies and what each step
adds}\label{dependencies-and-what-each-step-adds}
\addcontentsline{toc}{subsection}{Dependencies and what each step adds}

\begin{figure}
\centering
\includegraphics{.build/reconstruction-dependencies.png}
\caption{Dependencies for reconstruction. Arrows indicate prerequisite
routes; each theorem retains the hypotheses stated below.}
\end{figure}

Begin with quotients and the factorization rule: a map descends
precisely when it kills the relations. This one rule constructs tensor
products, module presentations, quotient operators, and cyclic
group-algebra compression. It tells us how to form the new object, not
whether a section exists.

Next add polynomial evaluation and Bézout. Coprime annihilators give
idempotent projectors. In the operator setting, the same identity gives
the inverse of a Sylvester operator. With an invertible derivative it
gives a separability tensor and a universal averaging operation. Over
the integers, the denominators that appear are arithmetic information:
they lead to resultants, finite cokernels, and Smith divisibility.

Add finite-dimensional duality to turn a failed linear equation into a
witness. The trace pairing identifies the annihilator of the Sylvester
image with reverse intertwiners. Add module presentations to compute the
same obstruction as an extension class, and Smith form to determine how
that class changes the middle module. The witness and the module type
are different outputs.

Finally add a coefficient lattice, an involution, and real positivity.
Cyclotomic evaluation preserves multiplication and involution and
identifies the integral coefficient lattice with a finite-index
sublattice of the product of cyclotomic integer rings. Allowing
arbitrary tuples in that larger product forgets the integral
compatibility. Trace reconstruction describes its exact image.
Membership in the affine lattice of odd coefficient vectors, together
with the exact total energy, forces signs. Compression discards a
complementary ideal, so it retains only part of the Gram equation; the
explicit core-nine defect lies in the discarded part.

No step requires importing schemes or differential topology. Homological
language is useful here because its first extension group has an
explicit matrix presentation; the Hadamard calculations themselves do
not depend on a general course in derived functors.

\subsection*{Eight principles to reconstruct rather than
memorize}\label{eight-principles-to-reconstruct-rather-than-memorize}
\addcontentsline{toc}{subsection}{Eight principles to reconstruct rather
than memorize}

\subsubsection*{1. Quotients encode
relations}\label{quotients-encode-relations}
\addcontentsline{toc}{subsubsection}{1. Quotients encode relations}

\textbf{Minimum data.} A module \(F\), a submodule of relations \(N\),
and a module map \(f:F\to Q\) with \(N\subseteq\ker f\).

\textbf{Construction.} Define \(\bar f(x+N)=f(x)\). The kernel condition
is exactly what makes the answer independent of the representative, and
surjectivity of \(F\to F/N\) proves uniqueness.

\textbf{Families recovered.} Balanced tensor maps, cokernel maps,
cyclic-module presentations, induced quotient operators, and the map
\(\mathbb Z[C_m]\to\mathbb Z[C_d]\).

\textbf{Failure to detect.} Factorization says nothing about a section.
The quotient \(\mathbb Z\to\mathbb Z/2\) has no additive section; an
invariant quotient of a Jordan block need not have an invariant
complement.

\subsubsection*{2. Bézout gives projectors and
inverses}\label{buxe9zout-gives-projectors-and-inverses}
\addcontentsline{toc}{subsubsection}{2. Bézout gives projectors and
inverses}

\textbf{Minimum data.} An endomorphism \(T\) of a module over a
commutative ring, with \(a(T)b(T)=0\) and \(ua+vb=1\) in that ring's
polynomial algebra.

\textbf{Construction.} The operators \(v(T)b(T)\) and \(u(T)a(T)\) add
to the identity. The first acts as the identity on \(\ker a(T)\) and
vanishes on \(\ker b(T)\); the second does the reverse. Their images lie
in those respective kernels, and their cross product is zero, so they
are complementary idempotents.

\textbf{Families recovered.} Primary decomposition, CRT coordinates,
polynomial spectral projectors, and the divided-difference inverse of
\(h\mapsto Ah-hB\) when \(a(A)=b(B)=0\).

\textbf{Failure to detect.} Coprimality over \(\mathbb Q[t]\) does not
imply an integral Bézout identity. The factors \(t\) and \(t-2\) produce
a rational graph with slope \(1/2\) and an integral obstruction of order
two.

\subsubsection*{3. A dual kernel certifies a missing
image}\label{a-dual-kernel-certifies-a-missing-image}
\addcontentsline{toc}{subsubsection}{3. A dual kernel certifies a
missing image}

\textbf{Minimum data.} A linear map \(D:H\to H'\) of finite-dimensional
vector spaces and perfect pairings with their dual spaces.

\textbf{Construction.} \(\operatorname{im}D=(\ker D^*)^\perp\). One
direction is evaluation; the converse extends a basis of the image to a
basis of the codomain and constructs a separating linear functional.

\textbf{Families recovered.} Solvability certificates, annihilator
identities, and the trace criterion \(C\in\operatorname{im}(Ah-hB)\)
exactly when \(\operatorname{tr}(gC)=0\) for every \(gA=Bg\).

\textbf{Failure to detect.} A field-valued witness cannot see a torsion
cokernel that disappears over that field. Integral divisibility requires
more than rational trace equations.

\subsubsection*{4. A presentation records both a module and how it is
attached}\label{a-presentation-records-both-a-module-and-how-it-is-attached}
\addcontentsline{toc}{subsubsection}{4. A presentation records both a
module and how it is attached}

\textbf{Minimum data.} Relations \(fu=0\), \(ge=cu\) over a PID \(R\),
with \(f,g\ne0\) and specified endpoint identifications with
\(R/(f),R/(g)\).

\textbf{Construction.} The relation matrix is
\(\left(\begin{smallmatrix}f&-c\\0&g\end{smallmatrix}\right)\). Changing
the chosen lift changes \(c\) by \(gh\) modulo \(f\). The extension
class is \(c\bmod(f,g)\); the middle module has Smith factors
\(\gcd(f,g,c)\) and \(fg/\gcd(f,g,c)\).

\textbf{Families recovered.} Cyclic Ext computations, Jordan-block
gluing, integer-lattice quotients, and the difference between
classification with fixed endpoint maps and classification after
forgetting them.

\textbf{Failure to detect.} Equal middle-module types need not mean
equivalent extensions. The residues \(1\) and \(1+t\) modulo \(t^2\)
give distinct endpoint-preserving classes but the same middle module
\(k[t]/(t^4)\).

\subsubsection*{5. A separability tensor repairs
maps}\label{a-separability-tensor-repairs-maps}
\addcontentsline{toc}{subsubsection}{5. A separability tensor repairs
maps}

\textbf{Minimum data.} An associative unital algebra \(S\) over a field
\(k\), with a tensor \(E=\sum_i a_i\otimes_k b_i\) satisfying
\(\sum_i a_ib_i=1\) and
\(\sum_i sa_i\otimes b_i=\sum_i a_i\otimes b_i s\).

\textbf{Construction.} For left \(S\)-modules, send a \(k\)-linear map
\(f\) to \(v\mapsto\sum_i a_i f(b_iv)\). The two tensor identities prove
\(S\)-linearity and show that this operation fixes every \(S\)-linear
map.

\textbf{Families recovered.} Maschke averaging and polynomial spectral
averaging. Sections and retractions are repaired when the surrounding
maps are already module maps.

\textbf{Failure to detect.} Averaging an arbitrary idempotent map need
not preserve idempotence. Endpoints annihilated by \(m\) do not imply
that the middle module of an extension is annihilated by \(m\). Over an
imperfect field, a field extension can be semisimple as a ring without
being separable over the ground field.

\subsubsection*{6. An injective rational decomposition can have an
integral
defect}\label{an-injective-rational-decomposition-can-have-an-integral-defect}
\addcontentsline{toc}{subsubsection}{6. An injective rational
decomposition can have an integral defect}

\textbf{Minimum data.} Monic \(f,g\in\mathbb Z[t]\) coprime over
\(\mathbb Q[t]\).

\textbf{Construction.} The integral CRT map has cokernel
\(\mathbb Z[t]/(f,g)\) as an additive group, of order
\(|\operatorname{Res}(f,g)|\). Smith form refines its order into exact
congruences. In a product of cyclotomic factors, trace inversion tests
the whole coefficient lattice.

\textbf{Families recovered.} Prime gluing, prime-power overlap jets, the
cyclic integral Fourier lattice, and the resultant formula for the
determinant of a Sylvester operator.

\textbf{Failure to detect.} Pairwise congruences among several factors
can miss a congruence against the whole previously glued block. At order
four, the tuple \((0,0,1+i)\) passes every pairwise overlap but
reconstructs half-integer coefficients.

\subsubsection*{7. Parity plus a sharp energy bound forces
signs}\label{parity-plus-a-sharp-energy-bound-forces-signs}
\addcontentsline{toc}{subsubsection}{7. Parity plus a sharp energy bound
forces signs}

\textbf{Minimum data.} A list \(x_1,\ldots,x_N\) of odd integers with
\(\sum_jx_j^2=N\).

\textbf{Construction.} Each squared coefficient is at least one.
Equality in the sum forces every square to be one.

\textbf{Families recovered.} The all-order affine-lattice sign
certificate: under the exact four-row spectral norm equations,
membership of each tuple in \(J+2L\) forces every reconstructed
coefficient to be a sign. Here \(L\) is the integral Fourier image and
\(J=(n,0,\ldots,0)\) is the image of \(1+t+\cdots+t^{n-1}\).

\textbf{Failure to detect.} Integral coefficients and the correct energy
alone allow zeros and larger entries. At odd order, component parity can
be combined with integral gluing because the CRT defect has no
two-torsion. At even order this simplification fails; one must test the
half-shifted tuple in \(L\).

\subsubsection*{8. An idempotent splits equations but not discrete
choices}\label{an-idempotent-splits-equations-but-not-discrete-choices}
\addcontentsline{toc}{subsubsection}{8. An idempotent splits equations
but not discrete choices}

\textbf{Minimum data.} A self-adjoint central idempotent \(e\) in a real
algebra with involution.

\textbf{Construction.} Write \(X=eX+(1-e)X\). Cross products vanish, so
a Gram equation splits into its two ideal components. For cyclic
compression, \(e\) is averaging over the kernel subgroup.

\textbf{Families recovered.} An exact retained-energy calculation, the
residual Gram equation, individual fiber lift counts, and the core-nine
defect \(8(S_p-S_k)\).

\textbf{Failure to detect.} A split real equation need not admit a
solution in the prescribed finite fiber vertex sets. Exact relaxation is
still relaxation when it forgets part of the domain.

\subsection*{A shorter route through the
books}\label{a-shorter-route-through-the-books}
\addcontentsline{toc}{subsection}{A shorter route through the books}

There are two coherent routes to canonical forms. A reader staying with
undergraduate linear algebra can use Xi's Cayley--Hamilton proof,
primary splitting, quotient-lift Jordan construction, and image--kernel
chain algorithm, assuming the minimal polynomial splits over the working
field. A reader who needs integral presentations and nonsplit operators
can instead establish polynomial torsion by linear dependence among
powers, prove one PID structure theorem, and derive rational and Jordan
forms. Using Jordan as the sole proof of Cayley--Hamilton while also
using Cayley--Hamilton to establish that Jordan route would be circular.

The source comparisons give precise inspected pages and stopping tests.
They support compressing repeated theorem coverage, not discarding all
alternate proofs. Xi's quotient-lift and filtration constructions teach
different skills; Hungerford's bounded-primary splitting lemma has
additional scope; Grillet's relative-basis proof is particularly useful
for the integral chapters. Xi already proves tensor universality and two
forms of Maschke's theorem. The graduate texts supply broader module
settings, exactness arguments, and different dependencies.

For the present project, learn one proof of each common structural
theorem, then reconstruct the new chapters in this order:

\begin{enumerate}
\def\labelenumi{\arabic{enumi}.}
\tightlist
\item
  Extract \(C\) from an invariant-subspace problem, solve \(Ah-hB=-C\),
  and give a trace witness for a nonsplit example.
\item
  Compute the same class through a cyclic presentation and determine its
  middle-module Smith factors. Check a characteristic-two example.
\item
  Construct the polynomial averaging tensor and the group averaging
  tensor, then explain precisely which sections each repairs.
\item
  Derive the prime integral CRT image and trace inverse. Find the
  missing congruence in the order-four pairwise counterexample.
\item
  Prove the affine-lattice sign theorem from odd coefficient energy and
  explain the odd/even distinction.
\item
  Count normalized symmetric lifts and reproduce the core-nine residual
  polynomial without enumerating Gram matrices by hand.
\end{enumerate}

These are performance tests, not a page quota. Passing them means the
reader can reconstruct a decomposition, name its domain of validity, and
exhibit a concrete obstruction to the tempting stronger assertion.

\section*{Architecture: structures, maps, and
obstructions}\label{architecture-structures-maps-and-obstructions}
\addcontentsline{toc}{section}{Architecture: structures, maps, and
obstructions}

\subsection*{The organizing claim}\label{the-organizing-claim}
\addcontentsline{toc}{subsection}{The organizing claim}

Much ordinary linear algebra is the study of finite-dimensional vector
spaces with additional structure. The most reusable move is to identify
that structure, declare which maps preserve it, and then ask which
familiar constructions remain available. An abstract tool earns its
place when it preserves information or exposes an obstruction that a
coordinate calculation conceals.

Four layers organize the collection:

\begin{enumerate}
\def\labelenumi{\arabic{enumi}.}
\tightlist
\item
  \textbf{Linear structure.} Kernels, images, quotients, duals, direct
  sums, and tensor products describe how linear information is
  discarded, tested, assembled, and transported.
\item
  \textbf{An action or a multiplication.} An endomorphism makes a vector
  space a \(k[t]\)-module. A group action makes it a \(k[G]\)-module. A
  bilinear multiplication is a map \(V\otimes V\to V\) satisfying
  specified identities.
\item
  \textbf{Decomposition and its obstructions.} Idempotents encode direct
  sums; coprime annihilators produce idempotents; semisimplicity
  guarantees invariant complements. Extensions measure ways in which
  complements fail to exist.
\item
  \textbf{Variation and locality.} Families of vector spaces introduce
  bundles and sheaves; changing coefficients introduces localization,
  integrality, and arithmetic. A theorem about each individual vector
  space need not control a family or an integral lattice.
\end{enumerate}

These layers are connected, but none erases the others. A classification
up to arbitrary linear isomorphism does not solve a problem about
lengths, positive forms, an integral lattice, or smooth variation.

\subsection*{The prerequisite map}\label{the-prerequisite-map}
\addcontentsline{toc}{subsection}{The prerequisite map}

The following arrows mean a useful prerequisite for the stated
destination, not a claim that every possible approach must take this
route.

\begin{longtable}[]{@{}
  >{\raggedright\arraybackslash}p{(\columnwidth - 4\tabcolsep) * \real{0.3333}}
  >{\raggedright\arraybackslash}p{(\columnwidth - 4\tabcolsep) * \real{0.3333}}
  >{\raggedright\arraybackslash}p{(\columnwidth - 4\tabcolsep) * \real{0.3333}}@{}}
\toprule\noalign{}
\begin{minipage}[b]{\linewidth}\raggedright
Foundation
\end{minipage} & \begin{minipage}[b]{\linewidth}\raggedright
Construction or theorem it supports
\end{minipage} & \begin{minipage}[b]{\linewidth}\raggedright
Destination and extra bridge
\end{minipage} \\
\midrule\noalign{}
\endhead
\bottomrule\noalign{}
\endlastfoot
Equivalence relations, linear maps, subspaces & Quotients, kernels,
first isomorphism theorem & Exact sequences; then modules and abelian
categories if needed \\
Bases, linear functionals, quotients & Dual exact sequences,
annihilators, evaluation & Adjoints after choosing a form; trace duality
after separability \\
Direct sums, polynomial division, Bézout & Coprime primary decomposition
& \(k[t]\)-modules, rational/Jordan forms, commuting operators \\
PID ideals and presentations & Smith normal form, cyclic module
decomposition & Integral abelian groups and operator classification;
higher-dimensional polynomial rings need different tools \\
Bilinear maps and quotients & Tensor products and Hom--tensor adjunction
& Scalar extension, representations, multilinear algebra \\
Tensor algebra and relations & Symmetric and exterior algebras &
Polynomial coordinate rings, determinants, differential forms after
smoothness \\
Splitting and projective presentations & \(\operatorname{Ext}^1\),
\(\operatorname{Tor}_1\) & Extension obstructions; derived functors
after complexes and homotopies \\
Groups, invariant subspaces, duals, tensors & Group representations,
Maschke, Schur & Matrix Lie groups after topology and the exponential
map; Lie algebra representations after differentiation \\
Finite field extensions and multiplication maps & Norm, trace, separable
trace pairing & Lang: integral closure, Dedekind ideals, valuations,
completions, ramification \\
Commutative rings, modules, localization & Affine coordinate rings,
local rings, tangent spaces & Hartshorne: sheaves, schemes, locally free
sheaves, cohomology \\
Real linear maps, inner products, multivariable calculus & Tangent
spaces, normal spaces, quadratic forms & Hirsch: manifolds,
function-space topology, approximation, transversality, bundles \\
\end{longtable}

An economical order for this study is

\[
\text{maps and quotients}\longrightarrow
\begin{cases}
\text{duality and splitting},\\
\text{polynomials and modules},\\
\text{bilinear maps and tensors}.
\end{cases}
\]

These three paths rejoin at invariant decompositions, determinant lines,
forms, and representations. Exactness becomes useful as soon as a
construction stops preserving injections or an invariant complement
fails. There is no need to put spectral sequences in front of that first
obstruction.

\subsection*{What the collection actually
contributes}\label{what-the-collection-actually-contributes}
\addcontentsline{toc}{subsection}{What the collection actually
contributes}

Xi's \CJKtext{基础代数} I develops the concrete languages of equations,
matrix maps, determinants, polynomial division, roots, symmetric
polynomials, and resultants. Volume II develops vector spaces and
operators, but also the metric and geometric meaning of those objects.
Volume III adds transformation groups, representations, modules, Lie
algebra examples, and Galois theory. This is a connected course, not
merely three fragments awaiting replacement by a graduate book.

Grillet is particularly well aligned with the present structural route:
VIII.6--7 connects PID modules to Jordan form, X handles splitting and
projectivity, XI develops Hom and tensors, XII introduces Ext/Tor, and
XIII builds free algebras and their quotients. Hungerford independently
recommends PID torsion modules before his linear algebra chapter
(printed p.~xviii, PDF p.~15). This is strong evidence that the module
viewpoint is a sound organizational choice, rather than a gratuitous
modernization of the user's plan.

Mac Lane supplies a language for recognizing the same construction in
different categories. His universal arrows and representability are
useful after several concrete examples have been worked through. They do
not by themselves provide Jordan bases, a spectral theorem, or a
treatment of Euclidean geometry.

Hilton--Stammbach supplies the next layer when exactness fails:
projective presentations, extension classes, complexes, and derived
functors. Its route through extensions before general derived functors
is especially compatible with a reader coming from linear algebra. Hall
turns familiar matrix groups and the exponential into systematic Lie
theory and develops representation theory beyond the initial
\(\mathfrak{sl}_2\) example.

Hartshorne and Lang are destinations requiring new bridges. Neither
should be mistaken for an interchangeable abstract-algebra text. Hirsch
is a different destination again: his preface (printed p.~vii, PDF p.~8)
deliberately avoids tensors and differential forms, and largely avoids
algebraic topology. The relevant foundations for his approach are
geometry, analysis, and topology as well as linear algebra.

\subsection*{The memory hierarchy}\label{the-memory-hierarchy}
\addcontentsline{toc}{subsection}{The memory hierarchy}

There are three different things to retain.

\textbf{Definitions that reconstruct maps.} Remember the quotient
factorization property, tensor universality, evaluation, and the
equations defining a split sequence or an idempotent. Each tells you how
a new map must be defined and what must be checked for it to exist.

\textbf{Theorems whose hypotheses control classification.} Remember
finite generation and the PID hypothesis in the module structure
theorem; splitting of the minimal polynomial for Jordan form; positivity
and the base field in the spectral theorem; invertibility of \(|G|\) in
Maschke's theorem; separability in trace duality. The missing hypothesis
often explains the entire difficulty.

\textbf{Geometric or computational interpretations.} A quotient is not a
preferred complement. A dual is not an inner-product identification. A
tensor is not usually a pure tensor. A canonical form is not a canonical
basis. A fiber is not a smooth bundle, and a finite-dimensional vector
space over a number field is not its integral lattice. These
distinctions protect the structural tools from misuse.

\subsection*{What this architecture does not
promise}\label{what-this-architecture-does-not-promise}
\addcontentsline{toc}{subsection}{What this architecture does not
promise}

Two commuting operators give a module over \(k[s,t]\), which is
generally not a PID. The one-operator classification therefore does not
automatically extend to pairs. Likewise, finite-dimensionality over
\(k\) does not make a \(k[G]\)-module semisimple, and a bundle of rank
\(n\) does not automatically possess a global frame. These are
productive boundaries, not reasons to abandon abstraction.

The existing algebra direction has been confirmed as the
Hadamard/Williamson study. Its dedicated bridge locates its Gram
matrices, integral cyclic group algebras, finite-field characters,
symmetry quotients, and Boolean constraints within this map. The other
subject bridges retain their independent value; they are not presented
as prerequisites for that particular problem.

\section{Quotients, duality, and
splitting}\label{quotients-duality-and-splitting}

\subsection{Conventions and the type of a
map}\label{conventions-and-the-type-of-a-map}

Throughout the core, \(k\) is a field and vector spaces are
finite-dimensional unless a paragraph explicitly removes that
hypothesis. Rings are associative with identity, modules are unital, and
modules are left modules unless sides are specified. A commutative ring
is stated when tensor products are themselves used as modules over that
ring. The zero vector space is allowed.

For a linear map \(f:V\to W\), \(f\) is the map and \(f(v)\) is its
value at \(v\in V\). Write \(\operatorname{im}f\subseteq W\) and
\(\ker f\subseteq V\). Composition \(g\circ f\) first applies \(f\). An
equality between maps includes their domains and codomains; when a
subspace is identified with its image, the inclusion will be stated.
This matters especially when comparing an operator on a quotient with a
restriction to a complement.

\subsection{Quotients: impose relations
once}\label{quotients-impose-relations-once}

Let \(U\subseteq V\) be a subspace. Define \(v\sim w\) if \(v-w\in U\)
and \(V/U=\{v+U:v\in V\}\). Addition and scalar multiplication of cosets
are well defined because replacing representatives changes the result by
an element of \(U\). Let \(q:V\to V/U\), \(q(v)=v+U\).

\textbf{Quotient principle.} For any vector space \(X\) and linear map
\(f:V\to X\), there is a unique linear map \(\bar f:V/U\to X\) with
\(\bar f\circ q=f\) if and only if \(U\subseteq\ker f\).

\textbf{Proof.} Necessity follows from \(q(u)=0\) for \(u\in U\).
Conversely set \(\bar f(v+U)=f(v)\). If \(v-w\in U\), then
\(f(v)=f(w)\); hence this is well defined. It is linear by the coset
operations. Surjectivity of \(q\) forces uniqueness. \(\square\)

The construction has a concrete description, but its reusable content is
the factorization property. When constructing a map from a quotient,
specify its value on a representative and prove it annihilates the
relations. When proving two maps from the quotient equal, precompose
with \(q\).

\textbf{First isomorphism theorem.} If \(f:V\to W\) is linear, then

\[V/\ker f\xrightarrow{\ \widetilde f\ }\operatorname{im}f,
\qquad \widetilde f(v+\ker f)=f(v)
\]

is a canonical isomorphism. It is surjective by definition of the image;
its kernel is zero by definition of \(\ker f\). Thus rank--nullity
follows from dimension of a quotient, rather than being an isolated
formula.

\textbf{Why universal objects are unique.} If \((Q,q)\) and \((Q',q')\)
have the same quotient property for \(U\subseteq V\), there are unique
maps \(a:Q\to Q'\) and \(b:Q'\to Q\) satisfying \(a q=q'\) and
\(b q'=q\). Then \(baq=q\), so uniqueness gives \(ba=1_Q\); similarly
\(ab=1_{Q'}\). Uniqueness means a unique isomorphism preserving the
specified structure, not a unique bare vector-space isomorphism.

Two useful reconstructions are

\[
U/(U\cap W)\cong (U+W)/W,
\qquad (V/U)/(W/U)\cong V/W\quad(U\subseteq W\subseteq V).
\]

For the first, restrict the quotient map \(V\to V/W\) to \(U\) and
compute its kernel. For the second, send \((v+U)+(W/U)\) to \(v+W\) and
apply the quotient principle twice. These proofs explain the maps as
well as the dimensions.

\textbf{Recognition signal.} Whenever a problem says ``modulo,''
identifies choices, or asserts that a quantity depends only on an
equivalence class, a quotient factorization is present. If an operation
is not compatible with the relation, the quotient does not inherit that
operation. For example, \(T:V\to V\) induces \(\bar T:V/U\to V/U\)
precisely when \(T(U)\subseteq U\).

\subsection{Duality: test vectors by scalar-valued
maps}\label{duality-test-vectors-by-scalar-valued-maps}

Define \(V^*=\operatorname{Hom}_k(V,k)\). For \(f:V\to W\) define

\[f^*:W^*\to V^*,\qquad f^*(\lambda)=\lambda\circ f.\]

Then \((g\circ f)^*=f^*\circ g^*\) and \(1_V^*=1_{V^*}\). The reversal
is forced by the domains. In dual bases the matrix of \(f^*\) is the
transpose of the matrix of \(f\); complex conjugation is absent here.

The evaluation map is

\[\epsilon_V:V\to V^{**},\qquad
\epsilon_V(v)(\lambda)=\lambda(v).\]

If \(e_1,\ldots,e_n\) is a basis with dual basis \(e^1,\ldots,e^n\),
then \(\epsilon_V(e_i)\) is the corresponding dual basis of \(V^{**}\).
Thus \(\epsilon_V\) is an isomorphism. It is natural: for \(f:V\to W\),

\[f^{**}\circ\epsilon_V=\epsilon_W\circ f.\]

Indeed, evaluate either side at \(v\in V\) and \(\lambda\in W^*\); both
give \(\lambda(f(v))\). The proof uses a basis to establish
invertibility, but the map itself uses no chosen basis.

There is generally no such preferred identification \(V\cong V^*\). A
basis or a nondegenerate bilinear form provides one, and different
choices provide different maps. Naturality of the double dual should not
be confused with a metric on \(V\).

For \(U\subseteq V\) define its annihilator

\[U^\circ=\{\lambda\in V^*: \lambda(u)=0\text{ for every }u\in U\}.\]

The quotient principle identifies \((V/U)^*\) canonically with
\(U^\circ\): pull back along \(q\). Restriction to \(U\) gives an exact
sequence

\[0\longrightarrow(V/U)^*\xrightarrow{q^*}V^*
\xrightarrow{\mathrm{res}}U^*\longrightarrow0.\]

To prove surjectivity of restriction, extend a basis of \(U\) to one of
\(V\), and prescribe zero on the added vectors. The extension exists; it
is usually not unique. The sequence is canonical even though a chosen
splitting is not.

It follows that \(\dim U^\circ=\dim V-\dim U\). Under \(V\cong V^{**}\),
\((U^\circ)^\circ=U\): inclusion is immediate, and the dimensions agree.
For \(f:V\to W\),

\[\ker f^*=(\operatorname{im}f)^\circ,
\qquad \operatorname{im}f^*=(\ker f)^\circ.\]

The second identity follows by factoring a functional vanishing on
\(\ker f\) through \(\operatorname{im}f\) and extending it to \(W\).
Consequently \(\operatorname{rank}f^*=\operatorname{rank}f\) and row
rank equals column rank. These are manifestations of testing a map by
functionals.

\textbf{Boundary.} For infinite-dimensional algebraic vector spaces,
dualization of short exact sequences is still exact when bases and their
usual choice principles are available. But \(V\to V^{**}\) need not be
surjective and \(V^*\otimes W\to\operatorname{Hom}(V,W)\) generally has
only finite-rank maps in its image. Continuous duals in analysis form a
different theory; the algebraic extension argument need not preserve
continuity.

\subsection{Direct sums and
idempotents}\label{direct-sums-and-idempotents}

The external direct sum \(A\oplus C\) has inclusions \(i_A,i_C\) and
projections \(p_A,p_C\) satisfying

\[p_Ai_A=1_A,\quad p_Ci_C=1_C,\quad p_Ai_C=p_Ci_A=0,
\quad i_Ap_A+i_Cp_C=1_{A\oplus C}.\]

These equations are more useful than the ordered-pair notation: they
allow block matrices and maps to be reconstructed without coordinates. A
finite direct sum is both a product and a coproduct. For infinite
families, the direct sum consists of finitely supported tuples, whereas
the product includes all tuples; they are different constructions.

\textbf{Idempotent principle.} For an endomorphism \(P:V\to V\), the
equation \(P^2=P\) is equivalent to a decomposition

\[V=\operatorname{im}P\oplus\ker P\]

in which \(P\) is projection onto the first summand.

\textbf{Proof.} Write \(v=P(v)+(v-P(v))\); the second term lies in
\(\ker P\). If \(x=P(y)\) and \(P(x)=0\), then \(x=P^2(y)=P(y)=0\).
Conversely projection in a direct sum plainly satisfies \(P^2=P\).
\(\square\)

This works for modules over any ring. If \(P\) commutes with an operator
\(T\), both summands are \(T\)-invariant. Conversely a decomposition
into \(T\)-invariant summands gives a projection commuting with \(T\).
Thus finding a decomposition can be translated into finding idempotents
in a commutant.

\textbf{Failure of a tempting test.} Pairwise zero intersection of three
subspaces does not imply their sum is direct. In \(k^2\), the lines
spanned by \(e_1\), \(e_2\), and \(e_1+e_2\) have pairwise intersection
zero, but satisfy a nontrivial three-term relation. The right test is
that the addition map \(\bigoplus_i U_i\to V\) be injective, or that
\(U_i\cap\sum_{j\ne i}U_j=0\) for each \(i\).

\subsection{Splitting: existence versus
choice}\label{splitting-existence-versus-choice}

A short exact sequence

\[0\longrightarrow A\xrightarrow{i}B\xrightarrow{p}C\longrightarrow0\]

means that \(i\) is injective, \(p\) is surjective, and
\(\ker p=\operatorname{im}i\). It splits if there exists a linear
section \(s:C\to B\) with \(p\circ s=1_C\). In that case

\[F:A\oplus C\to B,\qquad F(a,c)=i(a)+s(c)\]

is an isomorphism preserving the inclusion and quotient maps.

\textbf{Proof.} If \(F(a,c)=0\), apply \(p\) to get \(c=0\), then
injectivity of \(i\) gives \(a=0\). For \(b\in B\), the element
\(b-s(p(b))\) lies in \(\ker p\), so is \(i(a)\) for a unique \(a\).
This also describes \(F^{-1}\). A retraction \(r:B\to A\), \(r i=1_A\),
is equivalent: its kernel is a complement to \(i(A)\) and restricts
under \(p\) isomorphically to \(C\). \(\square\)

Over a field, lift a basis of \(C\) to \(B\) to obtain \(s\). Over a
general ring, the lifts may violate relations in \(C\); that is the
first source of extension obstructions.

The set of sections, if nonempty, is an affine space under
\(\operatorname{Hom}(C,A)\). For sections \(s,s'\), \(p(s'-s)=0\), so
uniquely \(s'-s=i\circ h\) for \(h:C\to A\). Conversely \(s+i h\) is a
section for any \(h\). This statement holds in module categories with
the indicated Hom group. It precisely locates the ambiguity.

There is no natural complement in general. In the split model
\(B=A\oplus C\), the automorphism

\[u_h(a,c)=(a+h(c),c)\]

preserves \(i\) and \(p\) but sends the section \(c\mapsto(0,c)\) to
\(c\mapsto(h(c),c)\). Whenever \(h\ne0\), no distinguished section can
be invariant under all such structure-preserving automorphisms. A chosen
positive inner product changes the problem: the orthogonal complement
then becomes canonical relative to that extra data.

\subsection{What to reconstruct and
when}\label{what-to-reconstruct-and-when}

For an ordinary rank problem, first factor through the image. For
constraints on functionals, take annihilators. For a proposed
decomposition, construct its projections and check their equations. For
an invariant subspace, ask explicitly whether the desired complement
must preserve an operator or a group action. The underlying vector-space
splitting proves less than an equivariant one.

The main source anchors are Xi II §§1.4--1.6 and 2.4 (printed
pp.~14--24, 67--71; PDF pp.~20--30, 73--77), Grillet VIII.2--5 and X.1
(especially X.1.5, printed p.~395, PDF p.~407), and Mac Lane III.1--2
for the broader universal-property language. Proofs above are written
out independently rather than inferred from those references.

\section{Operators as modules and polynomial
projections}\label{operators-as-modules-and-polynomial-projections}

\subsection{Put the operator into the scalar
ring}\label{put-the-operator-into-the-scalar-ring}

Let \(T:V\to V\) be \(k\)-linear. Define an action of \(R=k[t]\) on
\(V\) by

\[p\cdot v=p(T)(v),\qquad
p(T)=a_0 1_V+a_1T+\cdots+a_dT^d
\quad\text{if }p(t)=\sum_{j=0}^d a_jt^j.\]

Polynomial evaluation is a unital \(k\)-algebra homomorphism
\(k[t]\to\operatorname{End}_k(V)\), because the powers of a single
operator commute. Denote the resulting module by \(V_T\).

For another operator \(S:W\to W\), a map \(f:V\to W\) is \(k[t]\)-linear
exactly when it is \(k\)-linear and

\[f\circ T=S\circ f.\]

The forward implication follows by applying \(f\) to \(t\cdot v\); the
reverse follows first for powers, then for polynomials. Thus operators
and intertwining maps form the same category as \(k[t]\)-modules whose
underlying \(k\)-spaces lie in the chosen dimension class. In
particular, finite-dimensional operators correspond to modules
finite-dimensional over \(k\), a smaller class than all finitely
generated \(k[t]\)-modules. Invariant subspaces are submodules,
invariant quotient operators are quotient modules, and similarity means
module isomorphism. This translation preserves the maps, not merely the
list of isomorphism classes.

If \(n=\dim_k V<\infty\), a \(k\)-basis generates \(V_T\) over \(R\).
Moreover, \(1_V,T,\ldots,T^{n^2}\) are linearly dependent in
\(\operatorname{End}_k(V)\), so some nonzero polynomial annihilates
\(V_T\). Its annihilator ideal has a unique monic generator \(m_T\), the
minimal polynomial. For \(V=0\), set \(m_T=1\). This argument avoids
using Cayley--Hamilton before proving it.

\subsection{The PID structure theorem and its
mechanism}\label{the-pid-structure-theorem-and-its-mechanism}

\textbf{Structure theorem.} If \(R\) is a principal ideal domain and
\(M\) is a finitely generated \(R\)-module, then

\[M\cong R^r\oplus\bigoplus_{i=1}^s R/(d_i),
\qquad d_1\mid d_2\mid\cdots\mid d_s,\]

where the \(d_i\) are nonzero nonunits. The free rank and the
invariant-factor ideals \((d_i)\) are unique. Over \(k[t]\), choose the
\(d_i\) monic. The isomorphism and its cyclic generators are generally
not canonical.

Here is the construction in the Euclidean rings relevant to ordinary
linear algebra, \(k[t]\) and \(\mathbb Z\). Write \(M\) as a cokernel of
a matrix \(A:R^a\to R^b\). Such a finite presentation exists: a
submodule of a finite free module over a PID is finite free. Invertible
row and column operations change the bases of the target and source and
preserve its cokernel.

If \(A\ne0\), move a nonzero entry into the first position. Divide
entries in its column and row by this pivot. A nonzero remainder can be
moved to the pivot and has smaller Euclidean size. Repetition ends with
a pivot dividing its entire row and column, which can then be cleared.
If the pivot fails to divide an entry of the remaining block, add the
corresponding row to the first row; a subsequent column division
produces a smaller nonzero pivot. Thus this stage also terminates. The
pivot now divides the remaining block. Repeat on that block to obtain a
diagonal matrix with \(d_i\mid d_{i+1}\), including units and zero
entries as appropriate. Taking cokernels gives the theorem's existence
statement in these Euclidean cases.

For uniqueness, let \(I_j(A)\) be the ideal generated by all
\(j\times j\) minors. These ideals are unchanged by invertible row or
column operations; multilinearity of minors gives containment in one
direction and the inverse operation gives the other. For Smith diagonal
form,

\[I_j(A)=(d_1d_2\cdots d_j)\]

for the nonzero diagonal positions, including unit positions. Thus the
successive products, and hence the factors up to units, are determined.
For arbitrary module presentations, invariance of the corresponding
presentation ideals is a further theorem; alternatively uniqueness of
the abstract decomposition follows from primary filtrations below. We
are not identifying matrices of different presentation sizes without
that extra argument.

For a general PID there need not be a Euclidean norm, so the specific
descent algorithm above is not a proof in that generality. Grillet
VIII.6.1--6.3 (printed pp.~336--342, PDF pp.~348--354) gives the general
proof using principal ideals and a basis theorem. This is a cited
extension of the construction just explained.

\textbf{Why finite-dimensional operator modules have no free summand.}
The \(k\)-dimension of \(k[t]\) is infinite, whereas \(\dim_k V\) is
finite. Hence

\[V_T\cong\bigoplus_{i=1}^s k[t]/(d_i),\qquad d_1\mid\cdots\mid d_s.\]

Every block \(k[t]/(d)\) has basis \(1,t,\ldots,t^{d_0-1}\) if
\(d_0=\deg d>0\). Multiplication by \(t\) has companion matrix

\[C(d)=\begin{pmatrix}
0&0&\cdots&0&-a_0\\
1&0&\cdots&0&-a_1\\
0&1&\cdots&0&-a_2\\
\vdots&\vdots&\ddots&\vdots&\vdots\\
0&0&\cdots&1&-a_{d_0-1}
\end{pmatrix},\quad
d(t)=t^{d_0}+a_{d_0-1}t^{d_0-1}+\cdots+a_0.
\]

The subdiagonal entries are all \(1\) and unspecified entries are \(0\);
for \(d_0=1\) this is the one-entry matrix \((-a_0)\). The defining
relation \(d(t)=0\) explains the last column. Expansion of
\(\det(xI-C(d))\) along the last column gives \(d(x)\); equivalently the
cofactor of its \(j\)th last-column entry is the appropriate power of
\(x\). The invariant-factor blocks give the rational canonical form over
any field. Consequently

\[\chi_T=\prod_i d_i,\qquad m_T=d_s\quad(V\ne0).\]

In particular \(m_T\mid\chi_T\), so \(\chi_T(T)=0\). This derives
Cayley--Hamilton from the structure theorem, rather than assuming it in
the proof that the module is torsion. For \(V=0\) the empty product is
\(1\).

There is also a useful presentation that makes contact with matrices:

\[V_T\cong\operatorname{coker}\bigl(tI-A:R^n\to R^n\bigr),\]

where \(A\) is the matrix of \(T\) in a chosen basis. Send a vector
polynomial \(p(t)=\sum_j t^jv_j\) to \(\sum_j A^jv_j\). The relation
\(t v\equiv Av\) reduces every class modulo \((tI-A)R^n\) to a constant.
It also shows that a constant class maps to itself in \(k^n\), so no
nonzero constant is killed. Thus the induced map is bijective. Computing
Smith form of \(tI-A\) therefore computes similarity invariants of
\(A\); this is not the same as applying arbitrary left and right
equivalence to \(A\) over \(k\).

\subsection{Bézout and Chinese remainder
projections}\label{buxe9zout-and-chinese-remainder-projections}

Suppose \(f_1,\ldots,f_r\in k[t]\) are pairwise coprime and
\(m=f_1\cdots f_r\) annihilates \(T\). Set \(g_i=m/f_i\). Bézout gives
\(a_i g_i\equiv1\pmod{f_i}\). Let \(e_i=a_i g_i\) in \(k[t]/(m)\). Its
residues are \(1\) modulo \(f_i\) and \(0\) modulo every other factor.
The Chinese remainder theorem, or divisibility by each coprime factor,
gives

\[e_i^2=e_i,\qquad e_ie_j=0\ (i\ne j),\qquad\sum_i e_i=1
\quad\pmod m.\]

Evaluate at \(T\) to obtain projections \(P_i=e_i(T)\) satisfying these
same identities in \(\operatorname{End}_k(V)\). Their images are exactly
\(\ker f_i(T)\): first \(f_i e_i\) is divisible by \(m\), so the image
lies there; conversely \(e_i-1\) is divisible by \(f_i\), so \(e_i(T)\)
fixes that kernel. Hence

\[V=\bigoplus_i\ker f_i(T).\]

The polynomials \(a_i\) are not unique, but the projections are uniquely
determined by the specified factors and the action. If \(m\) is the
minimal polynomial and the factors are its primary powers, these are the
canonical primary summands. Any operator commuting with \(T\) commutes
with all the \(P_i\), and thus preserves these summands.

\textbf{Small example.} If \(T^2(T-1_V)=0\), take \(P_1=T^2\) and
\(P_0=1_V-T^2\). Indeed \(T^3=T^2\) implies \(T^4=T^2\), so
\(P_1^2=P_1\), and

\[V=\ker T^2\oplus\ker(T-1_V).\]

With \(T^2(T-1_V)^2=0\) instead, one may take

\[P_1=3T^2-2T^3,\qquad P_0=1_V-3T^2+2T^3.\]

These integer-coefficient formulas work in every characteristic: their
defining congruences are polynomial identities over \(\mathbb Z\). There
is no division by \(2\) or \(3\). Only the distinctness of \(0\) and
\(1\) is used.

\textbf{What fails without coprimality.} Take \(R/(t^2)\). Its submodule
\(tR/(t^2)\) and quotient \(R/(t)\) are both one-dimensional over \(k\).
They do not give an invariant direct sum: the only vectors killed by
\(t\) are in \(tR/(t^2)\). Dividing one repeated primary factor into
smaller pieces does not create Bézout projectors.

\subsection{Jordan form and what is actually
canonical}\label{jordan-form-and-what-is-actually-canonical}

Factor the invariant factors into monic irreducibles. Chinese
remaindering gives blocks \(k[t]/(p^a)\). This elementary-divisor
decomposition exists over every field. If the minimal polynomial splits
over \(k\), each \(p=t-\lambda\). Writing \(N=T-\lambda1_V\) on a block
of size \(a\), choose a cyclic vector \(v\) with \(N^{a-1}v\ne0\) and
\(N^av=0\). The ordered basis

\[N^{a-1}v,N^{a-2}v,\ldots,Nv,v\]

gives the Jordan block \(J_a(\lambda)\) with \(1\) on the superdiagonal.
Thus Jordan form requires splitting of \(m_T\) over \(k\), not
necessarily that \(k\) itself be algebraically closed.

To recover block sizes without choosing chains, put
\(a_j=\dim\ker(T-\lambda1_V)^j\) and \(a_0=0\). Every \(\lambda\)-block
of size \(b\) contributes \(\min(j,b)\); other primary blocks contribute
zero because \(T-\lambda1_V\) is invertible there. Therefore

\[\begin{aligned}
a_j-a_{j-1}&=\#\{\lambda\text{-blocks of size at least }j\},\\
2a_j-a_{j-1}-a_{j+1}&=\#\{\lambda\text{-blocks of size exactly }j\}.
\end{aligned}\]

For example, nullities \(0,3,5,6,6\) give three \(\lambda\)-blocks, two
of size at least \(2\), one of size at least \(3\), and none of size at
least \(4\): the sizes are \(3,2,1\). These are all the blocks if
\(\dim_kV=6\); otherwise another primary summand can remain invisible to
these kernels. Even saying that \(\lambda\) is the only eigenvalue in
\(k\) does not exclude a nonsplitting summand. This gives uniqueness of
the \(\lambda\)-block sizes and a practical reconstruction procedure.
More generally a \(p^b\)-block contributes \(\deg(p)\min(j,b)\) to
\(\dim_k\ker p(T)^j\), so these filtrations recover the elementary
divisors over a nonsplitting field as well.

The primary subspaces and block-size multiset are intrinsic. Individual
cyclic summands, Jordan chains, their ordering, and a conjugating matrix
are choices. Conflating a canonical list of invariants with a canonical
basis leads to many false naturality claims.

\textbf{Diagonalizability and semisimplicity.} \(T\) is diagonalizable
over \(k\) if and only if \(m_T\) is a product of distinct linear
factors. The \(k[t]\)-module \(V_T\) is semisimple if and only if
\(m_T\) is a product of distinct monic irreducible factors. These differ
when the field does not split \(m_T\). A real plane rotation through a
nonzero angle other than \(\pi\) can be a simple \(\mathbb R[t]\)-module
without having a real eigenvector.

\subsection{Hom spaces and the
commutant}\label{hom-spaces-and-the-commutant}

A module map \(R/(f)\to R/(g)\) is determined by the image \(y\) of
\(1\), and its well-definedness condition is \(f y=0\) in \(R/(g)\). For
monic \(d=\gcd(f,g)\) write \(f=d f_0\), \(g=d g_0\) with
\(\gcd(f_0,g_0)=1\). The condition \(g\mid f y\) is equivalent to
\(g_0\mid y\). Thus

\[\operatorname{Hom}_R(R/(f),R/(g))
\cong R/(d),\qquad h\bmod d\longmapsto
\bigl(1\longmapsto g_0h\bmod g\bigr).\]

Cancellation here takes place in the domain \(k[t]\). In particular

\[\dim_k\operatorname{Hom}_R(R/(f),R/(g))=\deg\gcd(f,g).\]

The commutant of \(T\) is \(\operatorname{End}_{k[t]}(V_T)\). If the
elementary divisors are \((t-\lambda)^{b_{\lambda,i}}\), then

\[\dim_k\{S:ST=TS\}
=\sum_\lambda\sum_{i,j}\min(b_{\lambda,i},b_{\lambda,j}).\]

This computes the number of independent commuting maps without solving
all matrix-entry equations. For a single cyclic block its commutant is
\(k[T]\). For multiple identical blocks there are extra intertwiners, so
a commuting operator need not be a polynomial in \(T\).

\subsection{Fitting decomposition and the boundary of the
method}\label{fitting-decomposition-and-the-boundary-of-the-method}

For \(N\ge\dim V\), both the kernels and images of powers of \(T\) have
stabilized. Then

\[V=\ker T^N\oplus\operatorname{im}T^N.\]

If \(v=T^Nw\) lies in the intersection, then \(T^{2N}w=0\);
stabilization gives \(w\in\ker T^N\), so \(v=0\). Rank--nullity gives
the direct sum. The restriction to the first summand is nilpotent and to
the second is bijective. These subspaces are canonical, though a basis
for either is not. This is the primary split into the factor \(t\) and
the factors coprime to it, obtained without explicitly factoring the
minimal polynomial.

For two commuting operators the ring is \(k[s,t]\). It is not a PID: the
ideal \((s,t)\) is not principal. A simultaneous eigenbasis exists for a
finite family of commuting diagonalizable operators over a common
splitting field: decompose into eigenspaces of one operator and note
that the others preserve them. Their restrictions remain diagonalizable
because their minimal polynomials divide the original square-free split
minimal polynomials, so induction applies. Without diagonalizability,
the one-variable module classification does not provide simultaneous
Jordan form.

\textbf{Source route.} Xi II Theorems 2.51--2.57 and the
existence/uniqueness proofs (printed pp.~72--77; PDF pp.~78--83) already
contain Bézout decomposition, quotient induction, and intrinsic block
counts. Grillet VIII.6--7 supplies the module derivation (printed
pp.~336--346; PDF pp.~348--358); Hungerford IV.6 and VII.4--5 supplies
another module-first route. The presentation, centralizer, and Fitting
arguments above are included so the route is operational rather than a
list of references.

\section{Hom and tensor products}\label{hom-and-tensor-products}

\subsection{Hom records maps and
variance}\label{hom-records-maps-and-variance}

\(\operatorname{Hom}_k(V,W)\) is a vector space under pointwise
operations. Maps \(a:V'\to V\) and \(b:W\to W'\) induce

\[\operatorname{Hom}(V,W)\to\operatorname{Hom}(V',W'),
\qquad f\longmapsto b\circ f\circ a.\]

The source variable reverses arrows and the target variable preserves
them. For modules over a noncommutative ring,
\(\operatorname{Hom}_R(M,N)\) is always an abelian group but generally
has no natural left \(R\)-module structure. Commutativity or specified
bimodule actions must be supplied before treating it as one.

The kernel universal property gives left exactness. For a short exact
sequence \(0\to A\xrightarrow{i}B\xrightarrow{p}C\to0\) of
\(R\)-modules,

\[0\to\operatorname{Hom}_R(X,A)\xrightarrow{i\circ-}
\operatorname{Hom}_R(X,B)\xrightarrow{p\circ-}
\operatorname{Hom}_R(X,C)\]

is exact. A map \(X\to B\) killed by \(p\) factors uniquely through
\(i\). There need not be surjectivity at the right: lifting the identity
of \(C\) would be a section of \(p\). Dually,

\[0\to\operatorname{Hom}_R(C,X)\to\operatorname{Hom}_R(B,X)
\to\operatorname{Hom}_R(A,X)\]

is exact. A map on the submodule \(A\) need not extend to \(B\). Over a
field both failures disappear, by choosing a splitting. The resulting
exactness is canonical; its proof need not make a canonical choice.

\subsection{Tensor universality and a
construction}\label{tensor-universality-and-a-construction}

A tensor product of vector spaces \(V,W\) is a vector space
\(V\otimes_k W\) together with a bilinear map
\(\tau:V\times W\to V\otimes_k W\) such that every bilinear map
\(\beta:V\times W\to X\) factors uniquely as

\[\beta=\widetilde\beta\circ\tau,
\qquad \widetilde\beta:V\otimes_k W\to X\text{ linear}.\]

Write \(v\otimes w=\tau(v,w)\). The pair, not the bare vector space, is
the universal object.

For existence, take the free vector space on symbols \([v,w]\) and
quotient by the span of all bilinearity relations:

\[\begin{gathered}
\relax[v+v',w]-[v,w]-[v',w],\quad
[v,w+w']-[v,w]-[v,w'],\\
[av,w]-a[v,w],\quad[v,aw]-a[v,w].
\end{gathered}\]

Every bilinear \(\beta\) extends linearly from the symbols and
annihilates these relations, hence factors through the quotient. The
symbols span, giving uniqueness. Universal uniqueness is proved by
composing the two comparison maps, as for quotients.

If \((e_i)\) and \((f_j)\) are bases, then \((e_i\otimes f_j)\) is a
basis. Spanning follows by expansion. For independence, the coefficient
maps \(\beta_{ij}(v,w)=e^i(v)f^j(w)\) induce linear functionals on the
tensor product extracting each coefficient of a finite relation. This
proof works for infinite algebraic bases too, because each vector and
tensor uses only finite sums. It does not assert that dual bases span
the whole algebraic dual in that case.

\textbf{Recognition signal.} A bilinear expression should usually be
turned into a linear map on a tensor product before choosing
coordinates. A formula on pure tensors defines a linear map only after
its bilinearity or balancing has been checked. Pure tensors span; they
are not generally all the elements.

For instance, \(e_1\otimes f_1+e_2\otimes f_2\) is not pure when the two
displayed pairs are independent. Under coefficient matrices it has rank
\(2\), whereas a pure tensor has entries \(v_iw_j\) and rank at most
\(1\).

\subsection{Functorial maps, symmetry, and
associativity}\label{functorial-maps-symmetry-and-associativity}

For \(a:V\to V'\) and \(b:W\to W'\), bilinearity gives the unique linear
map

\[a\otimes b:V\otimes W\to V'\otimes W',\qquad
(a\otimes b)(v\otimes w)=a(v)\otimes b(w).\]

Checking on pure tensors proves
\((a'\otimes b')\circ(a\otimes b)=(a'\circ a)\otimes(b'\circ b)\). The
same method gives canonical isomorphisms

\[\begin{gathered}
V\otimes W\cong W\otimes V,\qquad
(U\otimes V)\otimes W\cong U\otimes(V\otimes W),\\
k\otimes V\cong V,\qquad
(U\oplus V)\otimes W\cong(U\otimes W)\oplus(V\otimes W).
\end{gathered}\]

These identify objects by specified maps; they are not literal
equalities of set-theoretic constructions. Associativity is especially
transparent by observing that both iterated tensor products represent
trilinear maps.

In chosen product bases, the matrix of \(a\otimes b\) is the Kronecker
product. If \(A\) is \(n\times n\) and \(B\) is \(m\times m\),

\[\operatorname{tr}(A\otimes B)=\operatorname{tr}(A)\operatorname{tr}(B),
\quad \det(A\otimes B)=\det(A)^m\det(B)^n.\]

For trace, sum the diagonal entries \(a_{ii}b_{jj}\). For determinant,
write \(A\otimes B=(A\otimes I_m)(I_n\otimes B)\); the second factor is
a direct sum of \(n\) copies of \(B\), and the first is conjugate by the
basis permutation \(e_i\otimes f_j\mapsto f_j\otimes e_i\) to a direct
sum of \(m\) copies of \(A\). This proves the identity over every field
and also over commutative rings. No diagonalization hypothesis is
required.

\subsection{Currying and the Hom--tensor
adjunction}\label{currying-and-the-homtensor-adjunction}

There is a canonical isomorphism, with no finite-dimensional hypothesis,

\[\operatorname{Hom}_k(U\otimes V,W)
\cong\operatorname{Hom}_k(U,\operatorname{Hom}_k(V,W)).\]

It sends \(F\) to the map \(u\mapsto(v\mapsto F(u\otimes v))\).
Conversely a linear \(g:U\to\operatorname{Hom}(V,W)\) yields the
bilinear map \((u,v)\mapsto g(u)(v)\), hence a unique \(F\). These
operations are inverse, and pre/postcomposition proves naturality in all
variables.

One can remember this as: ``a linear map out of a tensor is a bilinear
map, and a bilinear map is a linear family of linear maps.'' The
adjunction is the precise version of that sentence, not an extra theorem
to memorize separately.

For modules over a commutative ring \(R\) the same formulas work. Over a
noncommutative ring one uses a right module and a left module and tracks
the remaining bimodule actions. A correct underlying abelian-group
version is

\[\operatorname{Hom}_{\mathbb Z}(M\otimes_R N,A)
\cong\operatorname{Hom}_{R^{\mathrm{op}}}
\bigl(M,\operatorname{Hom}_{\mathbb Z}(N,A)\bigr),\]

where \(M\) is right \(R\), \(N\) is left \(R\), and
\((\phi\cdot r)(n)=\phi(rn)\) gives the Hom group its right action. The
tensor represents biadditive maps satisfying
\(\beta(mr,n)=\beta(m,rn)\), called balanced maps. Two arbitrary left
modules over a noncommutative ring cannot simply be tensored over \(R\)
using the commutative formulas.

\subsection{Finite-dimensional Hom as a tensor and trace as
evaluation}\label{finite-dimensional-hom-as-a-tensor-and-trace-as-evaluation}

Define

\[\Theta:W\otimes V^*\to\operatorname{Hom}(V,W),\qquad
\Theta(w\otimes\lambda)(v)=\lambda(v)w.\]

If \(V\) is finite-dimensional, its inverse sends \(f\) to
\(\sum_i f(e_i)\otimes e^i\). Indeed evaluation at
\(v=\sum_i e^i(v)e_i\) recovers \(f(v)\). Conversely applying the
inverse to a pure tensor recovers it. The inverse expression is
basis-independent because it is the inverse of a specified canonical
map. \(W\) need not be finite-dimensional for this result.

For \(W=V\), compose \(\Theta^{-1}\) with the evaluation contraction
\(v\otimes\lambda\mapsto\lambda(v)\). This defines trace invariantly,
and the basis formula is \(\operatorname{tr}T=\sum_i e^i(T(e_i))\). The
tensor corresponding to \(1_V\) is \(\sum_i e_i\otimes e^i\); it is
basis-independent even though each summand is not.

Composition of rank-one maps is

\[\Theta(v\otimes\lambda)\circ\Theta(w\otimes\mu)
=\lambda(w)\Theta(v\otimes\mu).\]

Taking traces gives \(\lambda(w)\mu(v)\) in either order. Rank-one maps
span all endomorphisms, so
\(\operatorname{tr}(ST)=\operatorname{tr}(TS)\). This is an evaluation
identity, which explains why trace is invariant under similarity and why
every commutator has trace zero. The converse is not needed here, and in
positive characteristic trace of the identity can be zero without the
identity being zero.

For infinite-dimensional \(V\), every tensor in \(W\otimes V^*\) still
maps to a finite-rank operator, and every finite-rank operator arises
this way. The canonical map is injective: express a tensor with linearly
independent \(W\)-components and use coordinate functionals to show its
coefficients vanish if the operator vanishes. Thus its image is exactly
the finite-rank maps. In particular, finite-dimensional \(W\) also makes
\(\Theta\) surjective, even for infinite-dimensional \(V\): for a basis
\(w_i\) of \(W\) with coordinate functionals \(w^i\), write
\(f(v)=\sum_i(w^i\circ f)(v)w_i\). The exact surjectivity criterion is
that every linear map \(V\to W\) have finite rank; finite dimension of
either space is sufficient. When \(V=W\) is infinite-dimensional the
identity is a concrete obstruction to surjectivity.

This also exposes the issue in Xi II's remark after Theorem 6.19,
printed p.~278 (PDF p.~284). The preceding definition uses \textbf{all}
multilinear forms on \(V^p\times(V^*)^q\). For \(p=q=1\), evaluation
\((v,\lambda)\mapsto\lambda(v)\) is such a form, but it cannot be a
finite sum \(\sum_i\alpha_i(v)\lambda(w_i)\): that would represent
\(1_V\) as a finite-rank map. The finite-dimensional theorem is sound;
its unqualified infinite-dimensional extension needs to be withdrawn or
the tensor space restricted. The printed remark and definition were both
checked as page images, so this is a mathematical correction, not an OCR
inference.

\subsection{Exactness, quotients, and change of
scalars}\label{exactness-quotients-and-change-of-scalars}

Over a field, tensoring is exact. One proof chooses a splitting of a
short exact sequence and uses distribution over direct sums.
Alternatively a basis of \(W\) identifies \(-\otimes W\) with a direct
sum of copies of the original map. These are proofs of a canonical
exactness statement, not canonical splittings.

Over a commutative ring, tensoring is always right exact:

\[A\xrightarrow{f}B\longrightarrow C\longrightarrow0
\quad\Longrightarrow\quad
A\otimes N\xrightarrow{f\otimes1}B\otimes N
\longrightarrow C\otimes N\longrightarrow0.\]

To see the kernel, take the quotient of \(B\otimes N\) by the image of
\(A\otimes N\). Bilinear maps from \(B\times N\) vanishing on that image
are precisely the ones factoring through \(C\times N\). The quotient and
\(C\otimes N\) therefore satisfy the same universal property. In
particular

\[(M/U)\otimes_R N\cong
(M\otimes_R N)/\operatorname{im}(U\otimes_R N).\]

The word ``image'' matters: the induced map
\(U\otimes_R N\to M\otimes_R N\) need not be injective. For
\(0\to\mathbb Z\xrightarrow{\times2}
\mathbb Z\to\mathbb Z/2\to0\), tensoring with \(\mathbb Z/2\) turns the
first displayed map into zero. Thus a tensor product does not always
preserve subobjects. A module for which tensoring preserves all short
exact sequences is called flat.

If \(L/k\) is a field extension, \(L\otimes_k V\) has an \(L\)-action
\(a(b\otimes v)=ab\otimes v\). A \(k\)-basis \((e_i)\) gives an
\(L\)-basis \((1\otimes e_i)\), so dimension over the \textbf{new} field
remains \(\dim_k V\). The map \(v\mapsto1\otimes v\) is injective, and
scalar extension is exact. A real operator may acquire eigenvectors
after extending to \(\mathbb C\); its complexified invariant
decomposition need not descend to real one-dimensional subspaces. The
coefficient field is part of every classification statement.

\textbf{Sources.} Xi II §6.2, especially Definition 6.11, Theorem 6.15,
and Examples 6.13--6.14 (printed pp.~273--278; PDF pp.~279--284),
contains actual tensor universality and its applications. Grillet
XI.5--8 (printed pp.~434--456; PDF pp.~446--468) supplies balanced
products, general exactness, dual modules, and flatness. The concrete
calculations here show why that extra generality matters.

\section{Free algebras, exterior powers, determinants, and
minors}\label{free-algebras-exterior-powers-determinants-and-minors}

\subsection{Free constructions remember their extension
property}\label{free-constructions-remember-their-extension-property}

The tensor algebra of a vector space is

\[T(V)=\bigoplus_{r\ge0}V^{\otimes r},\qquad V^{\otimes0}=k,\]

with multiplication by concatenation of tensors. Each element has
finitely many nonzero degree components. The inclusions of \(k\) and
\(V\) are part of the construction. For any unital associative
\(k\)-algebra \(A\), a linear map \(f:V\to A\) extends uniquely to an
algebra homomorphism

\[\widehat f:T(V)\to A,\qquad
\widehat f(v_1\otimes\cdots\otimes v_r)=f(v_1)\cdots f(v_r).\]

The product on the right is multilinear, giving a map on each tensor
power; concatenation proves multiplicativity, and generation by \(V\)
proves uniqueness. An algebra generated by a specified copy or image of
\(V\) is therefore a quotient of \(T(V)\) by its relations. The quotient
has to impose exactly the intended identities; it cannot be guessed from
a superficial resemblance of formulas.

The symmetric algebra is

\[\operatorname{Sym}(V)=T(V)/\langle v\otimes w-w\otimes v:v,w\in V\rangle.\]

It is the free \textbf{commutative} unital algebra on \(V\). The
extension above factors precisely when the images \(f(v)\) commute. With
a basis \(e_1,\ldots,e_n\), sending \(e_i\) to a variable \(x_i\)
identifies it with \(k[x_1,\ldots,x_n]\). The variables and monomial
basis depend on the chosen basis, while the free property does not.

The homogeneous degree-\(r\) part \(\operatorname{Sym}^r(V)\) represents
symmetric \(r\)-linear maps from \(V^r\). Notice the variance:
polynomial expressions on a space \(V\) belong to
\(\operatorname{Sym}(V^*)\), not \(\operatorname{Sym}(V)\), once its
vectors are regarded as points. Over a finite field, formal polynomials
are not determined by their values on all points: \(x^q-x\) vanishes on
\(\mathbb F_q\) but is a nonzero polynomial. Thus a coordinate ring
should not silently be replaced by the ring of set functions on its
rational points.

\subsection{Exterior algebra: alternation is the
relation}\label{exterior-algebra-alternation-is-the-relation}

Define

\[\Lambda(V)=T(V)/\langle v\otimes v:v\in V\rangle.\]

The product is denoted \(\wedge\). Expanding \((v+w)\wedge(v+w)=0\)
gives \(v\wedge w=-w\wedge v\), and exchanging blocks gives
\(\alpha\wedge\beta=(-1)^{rs}\beta\wedge\alpha\) for homogeneous degrees
\(r,s\). The defining equation \(v\wedge v=0\) remains necessary in
characteristic \(2\): a sign rule alone does not imply it there.

The degree-\(r\) piece \(\Lambda^r(V)\) represents alternating
\(r\)-linear maps: an alternating map vanishes whenever two arguments
are equal. Its unique factorization through \(V^{\otimes r}\) kills the
repeated-vector relations; conversely the wedge map is alternating. This
construction works over any commutative ring, not only in characteristic
zero.

With an ordered basis \(e_1,\ldots,e_n\), the wedges

\[e_{i_1}\wedge\cdots\wedge e_{i_r},\qquad
1\le i_1<\cdots<i_r\le n\]

form a basis. Reordering and removing repetitions proves spanning. For
independence, for each ordered subset \(I=(i_1,\ldots,i_r)\) define on
vectors with coordinates \(v_{ab}\) the alternating coefficient form

\[\beta_I(v_1,\ldots,v_r)
=\sum_{\sigma\in S_r}\operatorname{sgn}(\sigma)
\prod_{a=1}^r v_{i_a,\sigma(a)}.\]

When two columns coincide, pair the permutations exchanging those
columns; their terms sum to zero, in characteristic \(2\) as well. Thus
\(\beta_I\) factors through \(\Lambda^r(V)\) and takes value \(1\) on
the \(I\) basis wedge and \(0\) on the other ascending ones. This proves
independence without dividing by \(r!\). Therefore

\[\dim\Lambda^r(V)=\binom nr,\qquad\Lambda^r(V)=0\ (r>n).\]

For \(r=0\), \(\Lambda^0(V)=k\). The zero-dimensional top exterior power
is therefore a line too.

\subsection{Antisymmetrizing tensors is a model with extra
hypotheses}\label{antisymmetrizing-tensors-is-a-model-with-extra-hypotheses}

If \(r!\) is invertible in \(k\), the operator

\[\operatorname{Alt}_r
=\frac1{r!}\sum_{\sigma\in S_r}\operatorname{sgn}(\sigma)\sigma
:V^{\otimes r}\to V^{\otimes r}\]

is idempotent. Its image identifies with \(\Lambda^r(V)\): the quotient
map restricted to that image is inverse to the map induced by
\(\operatorname{Alt}_r\). With this normalization, exterior
multiplication in the image is
\(\alpha\wedge\beta=\operatorname{Alt}_{r+s}(\alpha\otimes\beta)\).
Other common conventions place factorial factors in the identification,
so one must check which model a formula uses.

The analogous averaging operator without signs models symmetric powers
when \(r!\) is invertible. In small characteristic the quotient
definitions remain valid but these normalized projectors do not.
Symmetric powers (coinvariants) and invariant symmetric tensors must not
be identified using a nonexistent \(1/r!\). Polarization formulas with
factorial denominators have the same limitation. Xi II §6.4 explicitly
assumes characteristic zero for its antisymmetric-tensor model (printed
p.~292, PDF p.~298); Grillet XIII.4 uses the quotient construction over
a commutative ring (printed pp.~523--527, PDF pp.~535--539). These are
related constructions with different stated scope.

\subsection{Functoriality, rank, and
determinants}\label{functoriality-rank-and-determinants}

For \(f:V\to W\), the alternating map
\((v_1,\ldots,v_r)\mapsto f(v_1)\wedge\cdots\wedge f(v_r)\) induces

\[\Lambda^r f:\Lambda^r V\to\Lambda^r W.\]

Checking decomposable wedges gives
\(\Lambda^r(g\circ f)=\Lambda^r g\circ\Lambda^r f\) and
\(\Lambda^r1_V=1_{\Lambda^rV}\). For \(r>0\),

\[\Lambda^r f\ne0\quad\Longleftrightarrow\quad\operatorname{rank}f\ge r.\]

If the image has dimension less than \(r\), every image wedge is zero.
If the rank is at least \(r\), choose \(r\) vectors whose images are
independent and extend those images to a basis of \(W\); their wedge is
nonzero by the basis theorem. This is the intrinsic meaning of testing
rank by minors.

For \(n=\dim V\) define the determinant line \(\det V=\Lambda^nV\). It
is one-dimensional. An endomorphism \(T:V\to V\) therefore acts on it by
a unique scalar:

\[\Lambda^n T=(\det T)1_{\det V}.\]

This scalar is the determinant. Functoriality immediately gives
\(\det(ST)=\det S\det T\), and \(\det T\ne0\) precisely when \(T\) is
invertible. For \(V=0\) the determinant is \(1\) by the degree-zero
convention. Expanding \(T(e_1)\wedge\cdots\wedge T(e_n)\) gives the
usual signed permutation formula. Thus elimination formulas can be
reconstructed from alternation and the chosen volume basis.

If \(f:V\to W\) maps two different \(n\)-dimensional spaces, it
determines a map \(\det V\to\det W\). There is no preferred scalar
determinant until trivializations of both lines are chosen. This
distinction is useful in coordinate changes, orientation, and families
of vector spaces.

\subsection{Minors and Cauchy--Binet from
composition}\label{minors-and-cauchybinet-from-composition}

Let \(A\) be the matrix of \(f:k^n\to k^m\). In ascending wedge bases,
the entry of \(\Lambda^r f\) in row subset \(I\) and column subset \(J\)
is the minor \(\det A_{I,J}\). This follows by expanding the image of
\(e_{j_1}\wedge\cdots\wedge e_{j_r}\) and collecting the \(I\)
coefficient.

For a second matrix \(B:k^m\to k^\ell\), functoriality gives

\[\det(BA)_{I,J}
=\sum_{\substack{K\subseteq\{1,\ldots,m\}\\|K|=r}}
\det B_{I,K}\det A_{K,J}.\]

This is the full minors version of Cauchy--Binet. The familiar
determinant formula for a rectangular product is the case \(r=n=\ell\).
The same proof works over commutative rings using finite free modules
and the exterior basis theorem. A rank interpretation over a field
should not be carried unchanged to rings with zero divisors.

Expanding \(\det(xI-T)\) by the columns in which \(-T\) is selected
gives

\[\chi_T(x)=\sum_{r=0}^n(-1)^r
\operatorname{tr}(\Lambda^rT)x^{n-r}.\]

The trace of \(\Lambda^rT\) is the sum of principal \(r\)-minors. In
particular \(\operatorname{tr}T\) and \(\det T\) are the first and last
nontrivial elementary symmetric coefficients. This works without
eigenvalues in \(k\). In positive characteristic, recovering these
coefficients from power traces via Newton identities can require
division by a nonunit integer; the exterior-power formula itself does
not.

\subsection{Determinant lines of exact
sequences}\label{determinant-lines-of-exact-sequences}

For a short exact sequence of finite-dimensional vector spaces
\(0\to A\xrightarrow{i}B\xrightarrow{p}C\to0\), put \(a=\dim A\) and
\(c=\dim C\). There is a canonical isomorphism

\[\det A\otimes\det C\xrightarrow{\ \sim\ }\det B.\]

Send a decomposable tensor to

\[(a_1\wedge\cdots\wedge a_a)\otimes(c_1\wedge\cdots\wedge c_c)
\longmapsto
i(a_1)\wedge\cdots\wedge i(a_a)
\wedge\widetilde c_1\wedge\cdots\wedge\widetilde c_c,\]

where \(p(\widetilde c_j)=c_j\). Changing a lift adds a vector of
\(i(A)\), whose extra occurrence makes that term zero. Multilinearity
and alternation give a map on the determinant lines, and lifting a basis
of \(C\) gives a basis of \(B\), so it is an isomorphism. It is
independent of the chosen lifts, although the splitting of the original
sequence is not. The order ``\(A\) then \(C\)'' is part of the
convention; reversing it contributes \((-1)^{ac}\).

An operator preserving \(A\) therefore satisfies
\(\det T_B=\det T_A\det T_C\). Applied to \(xI-T\) over \(k(x)\), or
directly to its block triangular matrix, this gives
\(\chi_{T_B}=\chi_{T_A}\chi_{T_C}\) even without an invariant
complement. Multiplication of invariants across an exact sequence does
not imply splitting.

\subsection{Contraction and
reconstruction}\label{contraction-and-reconstruction}

For \(\lambda\in V^*\) and \(r\ge1\) define contraction
\(\iota_\lambda:\Lambda^rV\to\Lambda^{r-1}V\) by

\[\iota_\lambda(v_1\wedge\cdots\wedge v_r)
=\sum_{j=1}^r(-1)^{j-1}\lambda(v_j)
v_1\wedge\cdots\widehat v_j\cdots\wedge v_r.\]

The right side is alternating and multilinear in the \(v_j\), so the
universal property makes it well defined. Set contraction on
\(\Lambda^0V=k\) to zero; together these maps define a degree \(-1\)
operator on the exterior algebra. Splitting the sum into the positions
of \(\alpha\) and \(\beta\) gives

\[\iota_\lambda(\alpha\wedge\beta)
=(\iota_\lambda\alpha)\wedge\beta
+(-1)^r\alpha\wedge\iota_\lambda\beta
\quad(\deg\alpha=r).\]

This supplies a controlled introduction to signs in graded
constructions. Exterior powers, contractions, and determinant lines are
useful now; differential forms additionally require a manifold and
smooth variation. They should not be identified merely because the same
wedge symbol appears.

The core memory is short: free algebra extends linear maps, quotient
imposes relations, exterior power transports alternating data, and the
top exterior power is a line. Minors, rank tests, determinant
multiplication, and Cauchy--Binet can be rebuilt from those principles.

\section{Forms, adjoints, and the correct equivalence
relation}\label{forms-adjoints-and-the-correct-equivalence-relation}

\subsection{Similarity, equivalence, and congruence classify different
data}\label{similarity-equivalence-and-congruence-classify-different-data}

For an endomorphism \(T:V\to V\), changing the basis by columns \(S\)
changes its matrix from \(A\) to \(S^{-1}AS\). The source and target are
the same space and use the same changed basis: this is similarity.

For a map \(f:V\to W\) with independently changed source and target
bases, the matrix becomes \(Q^{-1}AP\). This is matrix equivalence; over
a field its only invariant is rank.

For a bilinear form \(b:V\times V\to k\) with matrix \(B\),
\(b(x,y)=x^{\mathsf T}By\). The same change of vector basis in both
arguments gives \(B\mapsto S^{\mathsf T}BS\), called congruence. It is
not similarity. The associated map

\[b^\flat:V\to V^*,\qquad b^\flat(v)(w)=b(v,w)\]

has a dual space as target; the two basis changes are linked by duality.
A nondegenerate form is exactly one for which this map is an
isomorphism. This identification depends on \(b\).

For a Hermitian form over \(\mathbb C\), this study uses the convention

\[h(x,y)=\overline{x}^{\mathsf T}Hy,\]

conjugate-linear in the first argument and linear in the second, with
\(H=\overline H^{\mathsf T}\). The matrix changes by \(H\mapsto S^*HS\)
where \(S^*=\overline S^{\mathsf T}\). The map \(v\mapsto h(v,-)\) into
the algebraic dual is conjugate-linear. Xi's inner-product convention is
linear in the first argument instead; translate the convention rather
than copying an adjoint formula unchanged.

\textbf{A decisive example.} Over \(\mathbb R\), \(I_2\) and \(2I_2\)
are congruent via \(S=\sqrt2 I_2\), but are not similar because they
have different eigenvalues. Similarly the real symmetric matrices
\(\operatorname{diag}(1,-1)\) and \(\operatorname{diag}(2,-2)\) have the
same inertia but different similarity classes. A theorem about Jordan
form cannot replace a theorem about forms.

\subsection{Symmetric forms and quadratic
forms}\label{symmetric-forms-and-quadratic-forms}

Assume \(\operatorname{char}k\ne2\). A symmetric bilinear form
determines \(q(v)=b(v,v)\), and

\[b(v,w)=\frac{q(v+w)-q(v)-q(w)}2.\]

Thus symmetric forms and quadratic forms encode the same data in this
characteristic. In characteristic \(2\), \(b(v,v)\) can be identically
zero for a nonzero alternating form, and a quadratic form needs its own
definition and polarization theory. Dividing by \(2\) conceals an
essential hypothesis.

Every symmetric bilinear form over a field of characteristic different
from \(2\) has a diagonal matrix in some basis. If \(b\ne0\),
polarization shows there is \(v\) with \(b(v,v)\ne0\). Then for any
\(w\),

\[w=\frac{b(v,w)}{b(v,v)}v+
\left(w-\frac{b(v,w)}{b(v,v)}v\right)\]

is a decomposition into \(kv\) and its \(b\)-orthogonal complement. The
sum is direct since \(b(v,v)\ne0\). Induct on dimension. If the
remaining restriction is zero, it contributes a radical summand. This
proof explains both diagonalization and its field dependence: nonzero
diagonal entries can be rescaled only by squares available in \(k\).

Over \(\mathbb R\), rescale the nonzero coefficients to \(+1\) or
\(-1\). The numbers \((p,q,r)\) of positive, negative, and zero
coefficients are unique (Sylvester's law of inertia). To see uniqueness
of \(p\), the diagonal form has a positive-definite subspace of
dimension \(p\). Any subspace of dimension greater than \(p\) meets the
span of the negative and zero coordinate vectors nontrivially by the
dimension inequality, and hence cannot be positive definite. Thus \(p\)
is the maximal dimension of a positive-definite subspace, an intrinsic
number. The same argument gives \(q\), and \(r\) is the dimension of the
radical.

Over an algebraically closed field of characteristic different from
\(2\), every nonzero coefficient is a square, so rank classifies
symmetric forms up to congruence. Over \(\mathbb Q\) or a number field,
rank alone does not: square classes and further arithmetic invariants
matter. ``Diagonal'' is not the same as ``classified.''

\subsection{Alternating forms and symplectic
pairs}\label{alternating-forms-and-symplectic-pairs}

A bilinear form is alternating if \(b(v,v)=0\) for all \(v\). This
implies \(b(v,w)=-b(w,v)\) in every characteristic. Conversely
skew-symmetry implies alternation only when \(2\ne0\).

If \(b\ne0\), choose \(u,v\) with \(b(u,v)\ne0\) and rescale to
\(b(u,v)=1\). The span \(P=\langle u,v\rangle\) has matrix
\(\begin{pmatrix}0&1\\-1&0\end{pmatrix}\) and is nondegenerate. For any
\(w\) define

\[w'=w+b(v,w)u-b(u,w)v.\]

Then \(b(u,w')=b(v,w')=0\), so \(V=P\oplus P^\perp\). Induction
decomposes an alternating form into symplectic pairs and a zero radical.
Its rank is even, and a nondegenerate alternating space has even
dimension. This argument works in characteristic \(2\) with the
definition of alternation retained.

The essential memory is the nondegenerate two-dimensional pair and the
orthogonal-complement induction, not a special sequence of row
operations.

\subsection{Adjoint operators require a chosen
form}\label{adjoint-operators-require-a-chosen-form}

Let \(h\) be a nondegenerate Hermitian form on a finite-dimensional
complex space. There is a unique \(T^\dagger\) such that

\[h(Tv,w)=h(v,T^\dagger w)\qquad(v,w\in V).\]

In a basis with Gram matrix \(H\), its matrix is

\[[T^\dagger]=H^{-1}[T]^*H.\]

This proves existence and uniqueness and shows explicitly why the
adjoint is ordinary conjugate transpose only in an orthonormal basis. It
also gives

\[\begin{gathered}
(ST)^\dagger=T^\dagger S^\dagger,\qquad
(aS+bT)^\dagger=\bar aS^\dagger+\bar bT^\dagger,\\
(T^\dagger)^\dagger=T.
\end{gathered}\]

For a nondegenerate real symmetric form the analogous formula is
\(H^{-1}A^{\mathsf T}H\). The algebraic dual map \(T^*\) is defined
without a form and has a different type; the adjoint uses the form to
return to \(V\).

If \(h\) is positive definite, every subspace has a canonical orthogonal
complement relative to \(h\). Moreover

\[U\text{ is }T\text{-invariant}
\quad\Longleftrightarrow\quad
U^\perp\text{ is }T^\dagger\text{-invariant}.\]

For example, if \(u\in U\) and \(w\in U^\perp\),
\(h(u,T^\dagger w)=h(Tu,w)=0\). This is the mechanism behind invariant
orthogonal decompositions. For an indefinite nondegenerate form a
subspace can intersect its own orthogonal complement: the isotropic line
\(\mathbb R(1,1)\) for \(\operatorname{diag}(1,-1)\) does. Nondegeneracy
of the ambient form alone does not give a direct sum with every
orthogonal complement.

\subsection{The finite-dimensional spectral theorem, with the hypotheses
visible}\label{the-finite-dimensional-spectral-theorem-with-the-hypotheses-visible}

\textbf{Theorem.} On a finite-dimensional positive-definite complex
inner-product space, \(T\) has an orthonormal eigenbasis if and only if
\(TT^\dagger=T^\dagger T\).

\textbf{Proof.} One direction follows from a diagonal matrix and its
conjugate transpose. For the other, call \(T\) normal. Any
\(N=T-\lambda1_V\) is normal and

\[\|Nv\|^2=h(v,N^\dagger Nv)
=h(v,NN^\dagger v)=\|N^\dagger v\|^2.\]

Hence \(Tv=\lambda v\) implies \(T^\dagger v=\bar\lambda v\). The
complex characteristic polynomial has a root, so choose a unit
eigenvector \(v\). The line \(kv\) is invariant under both \(T\) and
\(T^\dagger\); its orthogonal complement is therefore also invariant
under both. The restrictions commute and are adjoints there, so
induction gives the orthonormal eigenbasis. \(\square\)

A self-adjoint operator has real eigenvalues: if \(Tv=\lambda v\), then
\(h(v,Tv)=\lambda\|v\|^2\) is real. A unitary operator has eigenvalues
of modulus \(1\) by norm preservation. Real self-adjoint operators admit
real orthonormal eigenbases, for instance by complexifying and observing
that real eigenvalues yield real eigenspaces of the required dimensions.
A real normal operator need not have real eigenvectors: a planar
rotation gives a \(2\times2\) real block instead. The field and the
positive metric are doing actual work.

The nonzero nilpotent matrix \(\begin{pmatrix}0&1\\0&0\end{pmatrix}\)
has a Jordan basis, but no orthonormal eigenbasis. It is not normal
under the standard positive inner product. A purely module-theoretic
classification does not replace the spectral theorem's metric
conclusion.

\subsection{Gram matrices, least squares, and determinant
bounds}\label{gram-matrices-least-squares-and-determinant-bounds}

For vectors \(v_1,\ldots,v_r\) in a positive inner-product space, their
Gram matrix is \(G=(h(v_i,v_j))\). If \(X\) is the matrix with these
vectors as columns in an orthonormal basis, then \(G=X^*X\). Thus

\[c^*Gc=\left\|\sum_jc_jv_j\right\|^2\ge0,\]

and \(G\) is positive definite exactly when the vectors are independent.
This is the correct reason the normal-equations matrix \(X^*X\) is
invertible for a full-column-rank least-squares problem. If \(X\) is
rank deficient, the minimizer in coefficient space need not be unique,
though the orthogonal projection of the target onto
\(\operatorname{im}X\) is unique.

Gram--Schmidt writes independent vectors as \(X=QR\) with orthonormal
columns of \(Q\) and upper triangular \(R\). Its diagonal entries are
the lengths of successive orthogonal residuals, so

\[\det G=\prod_j|r_{jj}|^2\le\prod_j\|v_j\|^2.\]

Dependent vectors give determinant zero, so the inequality holds in
general. Equality for nonzero vectors occurs precisely when all
projections onto earlier vectors vanish, equivalently the vectors are
pairwise orthogonal. For a square real matrix this yields Hadamard's
determinant inequality \(|\det X|\le\prod_j\|v_j\|\). A sign matrix with
orthogonal rows attains the bound, linking metric geometry directly to
the intended Hadamard problem without using an unnecessary abstraction.

\subsection{Affine and projective information that must be
retained}\label{affine-and-projective-information-that-must-be-retained}

For an affine quadratic polynomial \(q(x)=b(x,x)+\ell(x)+c\) in
characteristic different from \(2\), a translation \(x=y+a\) changes its
linear part to \(\ell+2b(a,-)\). If \(b\) is nondegenerate one can
choose a unique \(a\) to cancel that part. If \(b\) is degenerate, the
component of \(\ell\) on the radical cannot be removed this way and
controls further geometric types. The constant and the permitted
transformations also matter. Similarity invariants of a coefficient
matrix do not classify affine quadrics.

Projectivizing identifies nonzero scalar multiples of a vector, and
homogeneous quadratic equations are themselves defined up to nonzero
scale. Real points, signature, isotropic directions, and the distinction
between affine and projective transformations are information a general
algebra text may not teach in the same way as Xi II Chapters 4--5. Those
chapters are selectively deferable for the operator route; their
geometric content is not thereby proved replaceable.

\textbf{Source anchors.} Xi II §1.7 (printed pp.~25--43), Chapter 3
(pp.~85--160), and Chapters 4--5 (pp.~161--262) supply forms, adjoints,
metric examples, orthogonal projection, and geometry. The normal
spectral theorem is Theorem 3.62 (printed pp.~126--128; PDF
pp.~132--134). The book comparison distinguishes these substantial
results from the graduate books' module and category-theoretic
strengths.

\section{Exact sequences, extensions, and calibrated homological
algebra}\label{exact-sequences-extensions-and-calibrated-homological-algebra}

\subsection{Exactness records how pieces
fit}\label{exactness-records-how-pieces-fit}

A sequence \(A\xrightarrow{f}B\xrightarrow{g}C\) is exact at \(B\) if
\(\operatorname{im}f=\ker g\). The weaker condition \(g\circ f=0\) only
gives containment. The quotient \(\ker g/\operatorname{im}f\) measures
the difference. For a short exact sequence of finite-dimensional vector
spaces, \(\dim B=\dim A+\dim C\); the dimension identity alone does not
specify the maps or prove exactness.

An exact sequence of modules splits when its quotient map has a
module-linear section. A module \(P\) is projective if every map
\(P\to C\) lifts across every surjection \(B\to C\). Equivalently every
short exact sequence ending in \(P\) splits: apply the lifting property
to \(1_P\), or form the pullback of a proposed map along the surjection.
A module is projective if and only if it is a direct summand of a free
module. For the nontrivial direction, choose a free surjection
\(F\to P\) and split it; for the converse, lift basis vectors of a free
module and restrict the lift to the summand.

All vector spaces over \(k\) are free, hence projective. They are also
injective: a map defined on a subspace extends by a basis extension.
Consequently

\[\operatorname{Ext}^j_k(V,W)=0,\qquad
\operatorname{Tor}^k_j(V,W)=0\qquad(j>0).\]

One may use the projective resolution concentrated in degree \(0\) to
see both vanishings. They are statements about the coefficient ring
\(k\) and its module category. They do not imply that every complex of
vector spaces is exact, nor that \(\operatorname{Ext}_{k[t]}\) or
\(\operatorname{Ext}_{k[G]}\) vanishes.

\subsection{The smallest equivariant nonsplitting
example}\label{the-smallest-equivariant-nonsplitting-example}

Let \(R=k[t]\), \(M=R/(t^2)\), and \(K=R/(t)\). Define

\[0\longrightarrow K\xrightarrow{i}M\xrightarrow{p}K\longrightarrow0,
\qquad i(a\bmod t)=at\bmod t^2,\quad
p(f\bmod t^2)=f\bmod t.\]

The maps are well defined and \(R\)-linear; the image of \(i\) is
exactly the kernel of \(p\). As vector spaces, the sequence splits via
\(s(a)=a\bmod t^2\). As modules, it does not split. Any lift of
\(1\in K\) has the form \(1+ct\) in \(M\), and

\[t(1+ct)=t\ne0,\]

whereas \(t\cdot1=0\) in \(K\). Thus no module-linear section exists. In
the vector basis \((t,1)\), multiplication by \(t\) has matrix
\(\begin{pmatrix}0&1\\0&0\end{pmatrix}\). Its invariant line \(kt\) has
no invariant complement. This is the exact same obstruction expressed as
a module extension, a failure of equivariant splitting, or a Jordan
block.

The underlying vector-space splitting is useful, but forgetting the
action has discarded the information responsible for the obstruction.

\subsection{Extensions and the first Ext group, with an explicit
computation}\label{extensions-and-the-first-ext-group-with-an-explicit-computation}

An extension of \(M\) by \(N\) is a short exact sequence
\(0\to N\to E\to M\to0\). Two extensions are equivalent when an
isomorphism of their middle terms commutes with the specified inclusion
and quotient maps and is the identity on \(N\) and \(M\). Merely having
isomorphic middle modules is a weaker relation. The equivalence classes
form \(\operatorname{Ext}^1_R(M,N)\) under Baer sum, with the split
extension as zero. The general equivalence with the derived-functor
definition is cited from Hilton--Stammbach III.1--2; we can construct
the important cyclic case directly.

Let \(R\) be a commutative domain, \(0\ne f\in R\), and \(N\) an
\(R\)-module. The free resolution

\[0\longrightarrow R\xrightarrow{\times f}R
\longrightarrow R/(f)\longrightarrow0\]

is exact. Apply \(\operatorname{Hom}_R(-,N)\). Evaluation at \(1\)
identifies the two Hom modules with \(N\), and the induced map is
multiplication by \(f\). Hence

\[\operatorname{Hom}_R(R/(f),N)=\ker(f:N\to N),
\qquad\operatorname{Ext}^1_R(R/(f),N)\cong N/fN,\]

and the positive Ext groups of degree greater than \(1\) vanish for this
first argument. The injectivity of \(\times f\) uses both the domain and
\(f\ne0\) hypotheses; over a ring with zero divisors this two-term
complex may not be a resolution.

The extension meaning of \(N/fN\) can be seen without the
derived-functor terminology. For an extension, choose \(e\in E\) over
\(1\bmod f\). Then \(fe=i(n)\) for a unique \(n\in N\). Replacing \(e\)
by \(e+i(n_0)\) changes \(n\) by \(fn_0\). Conversely, for \(n\in N\)
define

\[E_n=(N\oplus Re)/R(-n,fe),\]

meaning the relation \(fe=n\). The copy of \(N\) injects because
\(rf=0\) forces \(r=0\). The quotient is \(R/(f)\); every extension has
this form by expressing its elements using \(N\) and a lift of \(1\).
Changes of the lift give precisely the equivalences \(n\sim n+fn_0\).
The zero class permits a lift killed by \(f\), and therefore a section.
Pullback followed by addition of the two copies of \(N\) sends classes
\([n],[n']\) to \([n+n']\), explaining the addition law.

For \(R=k[t]\) and \(N=K=R/(t)\) this gives

\[\operatorname{Ext}^1_{k[t]}(K,K)\cong k.\]

The class \(a\in k\) has generators \(u,e\) and relations \(tu=0\),
\(te=au\). If \(a=0\) it is split. For every \(a\ne0\), the middle
module is isomorphic to \(R/(t^2)\), but the extension classes for
different \(a\) remain different with the endpoints fixed. Thus an
extension invariant contains more information than an isomorphism class
of a single operator.

\subsection{Tor detects lost
injections}\label{tor-detects-lost-injections}

Tensor the same free resolution with \(N\). Its positive homology is

\[\operatorname{Tor}^R_1(R/(f),N)=\ker(f:N\to N),\qquad
R/(f)\otimes_R N\cong N/fN.\]

The groups of degree greater than \(1\) vanish for this first argument.
For \(R=\mathbb Z\), \(f=m\ne0\), and \(N=\mathbb Z/n\),

\[\mathbb Z/m\otimes\mathbb Z/n\cong\mathbb Z/\gcd(m,n),
\qquad
\operatorname{Tor}^{\mathbb Z}_1(\mathbb Z/m,\mathbb Z/n)
\cong\mathbb Z/\gcd(m,n).\]

For the kernel, multiplication by \(m\) on \(\mathbb Z/n\) kills the
subgroup generated by \(n/\gcd(m,n)\), which has that order. This
supplies a simple nonzero Tor example and explains the injection failure
after tensoring the multiplication-by-\(m\) sequence.

An obstruction in Ext asks whether maps can lift or extend; an
obstruction in Tor asks what tensoring has collapsed. They are useful
descriptions only when the coefficient ring and the relevant sequence
are retained.

\subsection{Complexes and homology over a
field}\label{complexes-and-homology-over-a-field}

A chain complex is a sequence of maps \(d_n:C_n\to C_{n-1}\) satisfying
\(d_{n-1}\circ d_n=0\). Define

\[Z_n=\ker d_n,\qquad B_n=\operatorname{im}d_{n+1},\qquad
H_n(C)=Z_n/B_n.\]

A chain map \(f:C\to D\) satisfies
\(d_n^D\circ f_n=f_{n-1}\circ d_n^C\). It carries cycles to cycles and
boundaries to boundaries, hence induces \(H_n(f)\). The quotient
property constructs that map, and composition makes homology a functor.

\textbf{Why the field does not force zero homology.} The complex with a
single nonzero term \(C_0=k\) has \(H_0=k\). In a two-term complex
\(0\to k^2\xrightarrow{d}k^2\to0\) with \(d(x,y)=(x,0)\), both \(H_1\)
and \(H_0\) have dimension \(1\). All vector-space short exact sequences
involved split; the map \(d\) still need not have maximal rank or have
image equal to a prescribed cycle space.

In fact splitting gives a normal description of a complex. Choose
complements

\[Z_n=B_n\oplus\mathcal H_n,\qquad
C_n=Z_n\oplus L_n.\]

Then \(\mathcal H_n\cong H_n(C)\) and \(d_n:L_n\to B_{n-1}\) is an
isomorphism: it is injective because \(L_n\cap\ker d_n=0\), and its
image is all \(B_{n-1}\). Thus

\[C_n=B_n\oplus\mathcal H_n\oplus L_n,\]

with differential zero on \(B_n\oplus\mathcal H_n\) and an isomorphism
from \(L_n\) onto \(B_{n-1}\). A complex is therefore a direct sum of
its homology with zero differential and two-term contractible pieces,
after choices. The decomposition does not say the homology summands are
zero.

More explicitly let \(j:\mathcal H\to C\) and \(q:C\to\mathcal H\) be
inclusion and projection. Define \(s_n:C_n\to C_{n+1}\) as the inverse
of \(d_{n+1}:L_{n+1}\to B_n\) on \(B_n\), and zero on
\(\mathcal H_n\oplus L_n\). Then

\[d_{n+1}s_n+s_{n-1}d_n=1_{C_n}-j_nq_n.\]

This is a chain homotopy from the identity to the projection onto chosen
homology representatives. In particular an acyclic complex of vector
spaces is contractible. Neither the contraction nor representatives are
generally natural. Over \(k[t]\), the required invariant complements can
fail, so the same decomposition need not exist in that module category.

\subsection{A connecting homomorphism is a controlled lifting
construction}\label{a-connecting-homomorphism-is-a-controlled-lifting-construction}

Let \(0\to A\xrightarrow{i}B\xrightarrow{p}C\to0\) be a short exact
sequence of chain complexes, exact in every degree. If \(c\in C_n\) is a
cycle, lift it to \(b\in B_n\). Then \(p(d b)=d c=0\), so uniquely
\(d b=i(a)\) for \(a\in A_{n-1}\). Since \(i(d a)=d^2b=0\), injectivity
gives \(d a=0\). Define

\[\delta:H_n(C)\to H_{n-1}(A),\qquad\delta([c])=[a].\]

A different lift changes \(b\) by \(i(a_0)\) and therefore changes \(a\)
by \(d a_0\), a boundary. Changing \(c\) by a boundary and lifting its
primitive changes \(b\) by a boundary in \(B\), which does not alter
\(a\) except for the preceding lift ambiguity. Thus \(\delta\) is well
defined and linear.

The resulting sequence

\[\cdots\to H_n(A)\to H_n(B)\to H_n(C)
\xrightarrow{\delta}H_{n-1}(A)\to H_{n-1}(B)\to\cdots\]

is exact. For example, \(\delta([c])=0\) means \(a=d a_0\); replacing
\(b\) by \(b-i(a_0)\) gives a cycle lifting \(c\), so \([c]\) comes from
\(H_n(B)\). Conversely a cycle lift plainly gives \(\delta([c])=0\). At
\(H_{n-1}(A)\), a cycle \(a\) dies in \(B\) precisely when \(i(a)=d b\);
then \(p(b)\) is a cycle and maps under \(\delta\) to \([a]\). At
\(H_n(B)\), if \(p(b)=d c'\) lift \(c'\) to \(b'\) and observe that
\(b-d b'\) lies in \(i(A)\) and is a cycle. These checks prove exactness
at all three kinds of position.

This explains why exact sequences are not decorative notation: the
connecting map measures the obstruction to lifting a cycle as a cycle.

\subsection{A sensible stopping point for the homological
toolkit}\label{a-sensible-stopping-point-for-the-homological-toolkit}

For ordinary finite-dimensional linear algebra, retain exactness,
splitting, the explicit \(k[t]/(t^2)\) example, complexes, and the
one-step resolution calculations. Derived functors become worthwhile
when a concrete functor fails to be exact or an equivariant extension
remains nonsplit. General spectral sequences, derived categories, and
long resolutions can be deferred until a specific problem needs them.

For the Hadamard/Williamson direction, a common-lift obstruction is not
automatically an Ext class: sign constraints define a nonlinear finite
set, not a module extension with an identified coefficient category.
Calling it ``homological'' without constructing that category and
obstruction would add a label without mathematical content.

\textbf{Sources.} Grillet X.1, X.3, XI.2, XI.6, and XII.1--5 give the
general framework. Hilton--Stammbach III.1--4 (printed pp.~84--99; PDF
pp.~95--110) is especially useful for extension equivalence and
computations; IV.1--4 (printed pp.~117--130; PDF pp.~128--141) supplies
complexes, connecting maps, homotopy, and resolutions. The cyclic
extension construction and field-complex decomposition above are fully
derived so the reader can use the ideas before reading a whole
homological algebra text.

\section{Representations and
semisimplicity}\label{representations-and-semisimplicity}

\subsection{An action becomes a
module}\label{an-action-becomes-a-module}

A representation of a group \(G\) on a \(k\)-vector space is a
homomorphism \(\rho:G\to\operatorname{GL}(V)\). Extend linearly to the
group algebra \(k[G]\), whose elements are finite sums \(\sum_g a_g g\)
and whose product extends the group law bilinearly. Then

\[\left(\sum_g a_g g\right)\cdot v=\sum_g a_g\rho(g)(v)\]

makes \(V\) a \(k[G]\)-module. Conversely such a module restricts to a
group action because \(g^{-1}\) supplies the inverse operator.
Intertwining maps are exactly module maps, invariant subspaces are
submodules, and irreducible representations are nonzero simple modules.

For the infinite cyclic group, \(k[G]\cong k[t,t^{-1}]\); this describes
an invertible operator. For a cyclic group of order \(m\),
\(k[G]\cong k[t]/(t^m-1)\); its representations are operators satisfying
\(T^m=1_V\). If \(\operatorname{char}k\) does not divide \(m\), the
polynomial \(t^m-1\) is square-free, since its derivative is
\(mt^{m-1}\). Over a field containing its roots, such operators are
diagonalizable. Both separability and the splitting field enter this
elementary representation example.

\subsection{Maschke's averaging
argument}\label{maschkes-averaging-argument}

\textbf{Theorem.} If \(G\) is finite and \(|G|\) is invertible in \(k\),
every invariant subspace of a finite-dimensional representation has an
invariant complement. Hence every finite-dimensional representation is a
direct sum of irreducibles.

\textbf{Proof.} Let \(U\subseteq V\) be invariant. Choose an ordinary
linear projection \(P:V\to V\) with image \(U\) and \(P(u)=u\) for every
\(u\in U\). Define the endomorphism

\[\overline P=\frac1{|G|}\sum_{g\in G}\rho(g)P\rho(g)^{-1}.\]

Each summand has image in \(U\) and is the identity on \(U\), so the
average does too; in particular \(\overline P^2=\overline P\).
Reindexing the sum proves that \(\overline P\) commutes with every
\(\rho(h)\). Therefore \(\ker\overline P\) is an invariant complement.
If \(V\) is not irreducible, split it and repeat, decreasing dimensions.
\(\square\)

The averaged projection depends on the original projection; the theorem
asserts existence, not a preferred decomposition. Over \(\mathbb R\) or
\(\mathbb C\), one can instead average a positive inner product. It
remains positive because a nonzero vector has a positive contribution in
each summand. Its invariant orthogonal complements prove complete
reducibility with a geometric meaning.

\textbf{Characteristic failure.} In characteristic \(p\), let a
generator of \(C_p\) act on \(k^2\) by \(1+N\), where
\(N=\begin{pmatrix}0&1\\0&0\end{pmatrix}\). Since \(N^2=0\),
\((1+N)^p=1\) in characteristic \(p\) (also when \(p=2\)). This is a
representation with invariant line \(ke_1\) and no invariant complement,
exactly as in the \(k[t]/(t^2)\) example. Averaging would require
\(1/p\).

Finite-dimensional representations of infinite groups need not be
semisimple even over \(\mathbb C\): send \(n\in\mathbb Z\) to \(1+nN\).
The same invariant line fails to split. Finite-dimensionality alone does
not remove extensions.

\subsection{Schur's lemma and the role of the
field}\label{schurs-lemma-and-the-role-of-the-field}

If \(S,T\) are irreducible representations, any intertwiner \(f:S\to T\)
is either zero or an isomorphism. Its kernel and image are invariant,
and irreducibility forces the result. Thus \(\operatorname{End}_G(S)\)
is a division algebra over \(k\).

If \(k\) is algebraically closed and \(S\) is finite-dimensional, every
endomorphism \(f\) has an eigenvalue \(\lambda\). The intertwiner
\(f-\lambda1_S\) has nonzero kernel, so it is zero. Therefore
\(\operatorname{End}_G(S)=k1_S\). Neither the algebraic-closure nor
finite-dimensional hypothesis should be suppressed.

Over \(\mathbb R\), let \(C_4\) act on the plane by a quarter-turn
\(J=\begin{pmatrix}0&-1\\1&0\end{pmatrix}\). There is no real invariant
line, so the representation is irreducible, but its commuting
endomorphisms include all \(aI+bJ\), a copy of \(\mathbb C\). Thus
``irreducible implies scalar commutant'' is false over \(\mathbb R\).

For a finite-dimensional semisimple representation over an algebraically
closed field, let \(S_i\) represent the distinct simple types occurring
in \(V\). The evaluation map gives the canonical isotypic description

\[\bigoplus_i S_i\otimes_k\operatorname{Hom}_G(S_i,V)
\xrightarrow{\ \sim\ }V,\qquad s\otimes f\longmapsto f(s).\]

The group acts only on the first factor. To prove the map is an
isomorphism, choose a decomposition \(V\cong\bigoplus_i S_i^{m_i}\) and
use Schur's lemma: the Hom space for \(S_i\) is \(k^{m_i}\) and
contributes exactly that summand. The map was specified without this
decomposition, so invertibility does not make the chosen copies
canonical. The isotypic subspaces are intrinsic; a basis of each
multiplicity space is not. Their commutant is a product of matrix
algebras \(\prod_i M_{m_i}(k)\) after choosing those bases.

\subsection{Duals and tensors of
representations}\label{duals-and-tensors-of-representations}

Define the dual action by

\[\rho^\vee(g)(\lambda)=\lambda\circ\rho(g^{-1}).\]

The inverse compensates for reversal under duality, ensuring
\(\rho^\vee(gh)=\rho^\vee(g)\rho^\vee(h)\). Plain transpose without the
inverse would reverse the group multiplication. In dual bases its matrix
is \(\rho(g^{-1})^{\mathsf T}\), without complex conjugation.

For two representations, the diagonal tensor action is

\[g\cdot(v\otimes w)=\rho(g)(v)\otimes\sigma(g)(w).\]

Functoriality proves the group law. The canonical identification
\(W\otimes V^*\cong\operatorname{Hom}(V,W)\) carries the action to
\(f\mapsto\sigma(g)f\rho(g)^{-1}\). Its invariant vectors are the
intertwiners. This is a practical use of tensor variance, not merely an
alternative notation.

Characters are traces \(\chi_V(g)=\operatorname{tr}\rho(g)\). Direct
sums and tensor products satisfy

\[\chi_{V\oplus W}=\chi_V+\chi_W,\qquad
\chi_{V\otimes W}=\chi_V\chi_W.\]

For finite groups over \(\mathbb C\), characters determine isomorphism
classes; the full orthogonality theorem proving that statement is a
cited next step, not required to prove the two formulas. In modular
characteristic the ordinary complex character package cannot be copied
without changes.

\subsection{From groups to Lie algebras: a short concrete
bridge}\label{from-groups-to-lie-algebras-a-short-concrete-bridge}

For a differentiable representation of a matrix Lie group,
differentiation at the identity yields a Lie algebra representation
\(d\rho:\mathfrak g\to\operatorname{End}(V)\) preserving brackets.
Tensor actions differentiate according to the product rule:

\[X\cdot(v\otimes w)=(X\cdot v)\otimes w+v\otimes(X\cdot w).\]

Thus the Lie algebra operator is \(A\otimes I+I\otimes B\), whereas the
group operator is \(\rho(g)\otimes\sigma(g)\). Confusing these two
formulas changes the representation. Dualization differentiates to
\(X\cdot\lambda=-\lambda\circ d\rho(X)\).

Here is an explicit algebraic example over a characteristic-zero field.
On homogeneous polynomials of degree \(m\) in \(x,y\), set

\[E=x\frac{\partial}{\partial y},\qquad
F=y\frac{\partial}{\partial x},\qquad
H=x\frac{\partial}{\partial x}-y\frac{\partial}{\partial y}.\]

On the basis \(v_j=x^{m-j}y^j\), \(0\le j\le m\),

\[Ev_j=jv_{j-1},\qquad Fv_j=(m-j)v_{j+1},\qquad Hv_j=(m-2j)v_j,\]

with missing endpoint vectors interpreted as zero. Direct calculation
gives \([H,E]=2E\), \([H,F]=-2F\), and \([E,F]=H\). This defines a
representation of \(\mathfrak{sl}_2\).

It is irreducible. The distinct eigenvalues \(m-2j\) of \(H\) allow
polynomial projections onto its individual weight spaces by
interpolation. A nonzero invariant subspace therefore contains some
\(v_j\). Repeated application of \(E\) reaches \(v_0\), since the
intervening integers are nonzero in characteristic zero; repeated
application of \(F\) then reaches every \(v_j\). The argument explicitly
explains where small positive characteristic can fail.

Over \(\mathbb C\), these are all finite-dimensional irreducible
\(\mathfrak{sl}_2\) representations. Hall §4.6 proves this
classification using a highest-weight vector and the commutator identity
\(E F^jv=j(m-j+1)F^{j-1}v\) (printed pp.~96--101, PDF pp.~109--114). The
notes above prove the model and its irreducibility, while the exhaustion
statement is cited.

The associated \(SU(2)\) representation is
\(\operatorname{Sym}^m\mathbb C^2\). The element \(-I\) acts by
\((-1)^m\), so it descends through \(SU(2)\to SO(3)\) exactly when \(m\)
is even. Integrating a Lie algebra representation and descending through
a covering are extra global questions; Lie algebra isomorphisms alone do
not identify the corresponding groups or their representations. Hall's
group-versus-Lie-algebra sections and Xi III §3.6 are useful precisely
at this boundary.

\textbf{Selective sources.} Xi III §3.2 gives an actual averaging proof
(Theorem 3.20, printed pp.~109--111, PDF pp.~115--117), §3.7 constructs
dual and tensor representations (pp.~152--162, PDF pp.~158--168), and
§4.6 introduces \(\mathfrak{sl}_2\) modules (starting p.~215, PDF
p.~221). Grillet IX.2 and IX.7 place these in semisimple module theory.
Hall §§4.3--4.7 gives the finite-dimensional Lie-theoretic version.
Compact-group averaging requires integration and continuity, whereas the
finite-group proof above requires only an invertible group order.

\textbf{Terminology warning.} Grillet IX.7 initially uses
``irreducible'' to mean indecomposable, that is, admitting no nontrivial
invariant direct-sum decomposition (printed p.~381, PDF p.~393). He
explicitly relates that definition to indecomposable modules on p.~383
(PDF p.~395). This study and Xi/Hall use the standard invariant-subspace
definition, corresponding to a simple module. They coincide in a
semisimple setting, but not for a nontrivial Jordan block. Thus Grillet
IX.7.1, asserting a decomposition of every finite-dimensional
representation into his ``irreducibles,'' is not the
complete-reducibility theorem with the terminology of these notes. This
is a verified convention difference, not a mathematical counterexample
to his stated definition.

\section{Bridges to arithmetic, algebraic geometry, and differential
topology}\label{bridges-to-arithmetic-algebraic-geometry-and-differential-topology}

These sections explain how the same constructions lead outward. They are
selective derivations and prerequisite diagnoses, not compressed
substitutes for the destination books. The Hadamard/Williamson bridge
follows separately and uses only the directions that have a concrete
role there.

\subsection{Finite extensions as linear algebra: trace and
norm}\label{finite-extensions-as-linear-algebra-trace-and-norm}

Let \(L/k\) be a finite field extension of degree \(n\). Every
\(a\in L\) defines the \(k\)-linear multiplication map

\[m_a:L\to L,\qquad m_a(x)=ax.\]

Define \(\operatorname{Tr}_{L/k}(a)=\operatorname{tr}(m_a)\) and
\(N_{L/k}(a)=\det(m_a)\). Since \(m_{ab}=m_a\circ m_b\) and
\(m_{a+b}=m_a+m_b\), norm is multiplicative and trace is \(k\)-linear.
For \(c\in k\), \(m_c=c1_L\), so trace is \(nc\) and norm is \(c^n\).
The map \(a\mapsto m_a\) is faithful: evaluate at \(1\) to recover
\(a\).

For \(L=k(\alpha)\) with monic minimal polynomial
\(f(t)=t^n+a_{n-1}t^{n-1}+\cdots+a_0\), multiplication by \(\alpha\) has
the companion matrix of \(f\). Hence

\[\operatorname{Tr}_{L/k}(\alpha)=-a_{n-1},\qquad
N_{L/k}(\alpha)=(-1)^n a_0.\]

If \(L/k\) is separable, let \(\Omega\) be an algebraic closure of
\(k\). The primitive-element theorem writes \(L=k(\alpha)\); it is cited
here, not assumed for inseparable extensions. The polynomial \(f\) has
distinct roots \(\alpha_1,\ldots,\alpha_n\), and

\[\Omega\otimes_k L\cong\Omega[t]/(f)
\cong\prod_{i=1}^n\Omega\]

by scalar extension and Chinese remaindering. Multiplication by \(a\)
becomes diagonal with entries \(\sigma_i(a)\) for the \(n\) embeddings
\(\sigma_i:L\to\Omega\). Trace and determinant are unchanged by scalar
extension of the matrix, so

\[\operatorname{Tr}_{L/k}(a)=\sum_i\sigma_i(a),\qquad
N_{L/k}(a)=\prod_i\sigma_i(a).\]

This is the same polynomial-splitting mechanism used for operators. The
new content is that the regular representation comes from a field
multiplication.

The trace pairing \(b(a,b)=\operatorname{Tr}_{L/k}(ab)\) is
nondegenerate for finite separable extensions. In the power basis its
Gram matrix is \(V^{\mathsf T}V\), where \(V_{ij}=\alpha_i^{j-1}\). The
Vandermonde determinant is nonzero because the roots are distinct;
therefore its square, the determinant of the trace Gram matrix, is
nonzero. This proves nondegeneracy even in positive characteristic when
the extension is separable; positivity has no role.

For a finite purely inseparable counterexample, take
\(k=\mathbb F_p(s)\) and \(L=k(\alpha)\) with \(\alpha^p=s\). In the
basis \(1,\alpha,\ldots,\alpha^{p-1}\), multiplication by
\(\sum_{j=0}^{p-1}a_j\alpha^j\) has every diagonal entry \(a_0\); the
other powers produce cyclic shifts. Thus its trace is \(pa_0=0\). The
trace pairing is identically zero although \([L:k]=p\). A field
extension being finite does not suffice for trace duality.

Xi III §5.8 (printed pp.~272--278, PDF pp.~278--284) is the immediate
Galois-theoretic source. Lang III.1 (printed pp.~57--62, PDF pp.~71--76)
uses the separable trace pairing to define complementary integral
modules. The latter retains an integral lattice, which is extra
information beyond the \(k\)-vector space and its nondegenerate pairing.

\subsection{Why arithmetic needs more than field
extensions}\label{why-arithmetic-needs-more-than-field-extensions}

In a number field, the ring of integers is a \(\mathbb Z\)-module inside
a \(\mathbb Q\)-vector space. Tensoring with \(\mathbb Q\) recovers the
ambient vector space but loses torsion and integral index. For example
\(2\mathbb Z\subset\mathbb Z\) becomes an isomorphism after tensoring
with \(\mathbb Q\), while its cokernel is \(\mathbb Z/2\) before
tensoring. Integral bases and discriminants cannot be read from the
rational vector space alone.

Localizing a commutative ring \(R\) at a multiplicative set \(S\)
produces \(S^{-1}R\) with a map from \(R\) making all elements of \(S\)
invertible. It is universal for that property. For modules,
\(S^{-1}M\cong S^{-1}R\otimes_R M\) and localization is exact. For
instance if a localized element maps to zero, some denominator in \(S\)
annihilates its original image; multiplying the numerator by that
denominator reduces the kernel assertion to the original exact sequence.
This proves the middle exactness directly, with surjectivity obtained by
lifting numerators.

A finitely generated torsion-free module over a Dedekind domain is
projective (Lang I.9, Proposition 26, printed p.~29, PDF p.~43). It need
not be free: a nonprincipal invertible ideal is a rank-one example. Its
scalar extension to the fraction field is one-dimensional, so
vector-space dimension misses its ideal class. Over a PID, this
distinction vanishes for finitely generated torsion-free modules. That
is the arithmetic boundary of the elementary structure theorem.

To read Lang efficiently, the missing bridges are integral dependence
and integral closure, prime ideal factorization in Dedekind domains,
discrete valuations, completions, and local ramification. Some algebraic
ingredients appear in Hungerford VIII and Grillet's commutative algebra
chapters; their existence in a table of contents does not establish
complete preparation for Lang's local and analytic arguments. Lang I.9
explicitly says it is not used later in that book, so projectivity
should be read for this conceptual bridge, not falsely advertised as a
prerequisite for all his subsequent chapters.

\subsection{Affine geometry: symmetric algebra, ideals, and tangent
vectors}\label{affine-geometry-symmetric-algebra-ideals-and-tangent-vectors}

For a finite-dimensional vector space \(V\), the affine coordinate
algebra is \(\operatorname{Sym}(V^*)\). Imposing polynomial equations
amounts to taking a quotient by an ideal. Over an algebraically closed
field, passage between radical ideals and algebraic sets requires the
Nullstellensatz; it is a substantial theorem, not an immediate
consequence of the quotient property. Nonradical ideals retain nilpotent
information invisible to a set of points.

For a finitely generated \(k\)-algebra \(A\) and a \(k\)-rational point
\(x:A\to k\), put \(\mathfrak m=\ker x\). A tangent vector at \(x\) can
be defined as a \(k\)-linear derivation \(D:A\to k\) satisfying

\[D(ab)=x(a)D(b)+x(b)D(a).\]

It vanishes on constants and on \(\mathfrak m^2\). Conversely every
functional \(\ell:\mathfrak m/\mathfrak m^2\to k\) defines
\(D(a)=\ell((a-x(a))\bmod\mathfrak m^2)\). Expanding a product proves
the derivation law, and these assignments are inverse. Thus

\[T_x\operatorname{Spec}A\cong
\operatorname{Hom}_k(\mathfrak m/\mathfrak m^2,k).\]

This statement uses a rational point, so its residue field is \(k\). At
a general scheme point the residue field and the relative differential
module must be tracked.

For \(A=k[x_1,\ldots,x_n]/I\) at a rational point \(a\), expand every
\(f\in I\) modulo the square of the maximal ideal:

\[f(a+u)\equiv f(a)+\sum_i
\frac{\partial f}{\partial x_i}(a)u_i\pmod{(u)^2}.\]

The tangent space is the common kernel of these linear equations. For
the cusp \(y^2=x^3\) at the origin, there is no nonzero linear equation
from its defining polynomial, so the tangent space is two-dimensional
even though the curve has dimension one. For
\(A=k[\varepsilon]/(\varepsilon^2)\), the point has a one-dimensional
tangent space despite its underlying point set having one point.
Nilpotents store infinitesimal directions.

Tensor products now do geometric work: affine fiber products are
represented by \(A\otimes_R B\). Fibers of a module use
\(M\otimes_R\kappa(\mathfrak p)\). Because this tensor operation may
fail to preserve injections, flatness is needed to control families. The
finite-dimensional vector-space version of exactness is insufficient
once the coefficient ring varies.

\subsection{What the scheme and sheaf bridge
adds}\label{what-the-scheme-and-sheaf-bridge-adds}

Hartshorne II.1--2 introduces sheaves and schemes; II.5 develops sheaves
of modules, and II.8 introduces differentials. These topics add local
data and gluing, not just new names for global vector spaces. A locally
free sheaf is locally a finite direct sum of the structure sheaf; a
global free basis is an extra condition. On an affine scheme, finite
locally free sheaves correspond to finitely generated projective
modules, a theorem requiring the module--sheaf correspondence and its
local criteria.

Global sections form a left-exact functor; their higher derived functors
are sheaf cohomology. The obstruction here concerns sheaves and gluing
even when the ground field is a field. It is not contradicted by
\(\operatorname{Ext}^j_k(V,W)=0\): these are different categories and
different functors. Hartshorne III.1--3 develops that transition. His
introduction marks externally cited commutative-algebra and homological
results with an ``A'' instead of proving them all in place (printed
pp.~xv--xvi, PDF pp.~14--15). The collection provides many algebraic
ingredients, but a reader still needs localization, local rings, sheaves
and stalks, and substantial commutative-algebra fluency before this
becomes a comfortable route.

\subsection{Differential topology: linear algebra varies
smoothly}\label{differential-topology-linear-algebra-varies-smoothly}

For a smooth map \(F:M\to N\), differentiation gives maps
\(dF_x:T_xM\to T_{F(x)}N\). The chain rule makes differentiation
functorial. At a regular point of a level set \(F^{-1}(y)\), the
implicit-function theorem identifies

\[T_xF^{-1}(y)=\ker dF_x.\]

For a smooth embedded submanifold \(M\subset N\), there is an exact
sequence of vector bundles over \(M\),

\[0\to TM\to TN|_M\to\nu\to0,\]

where \(\nu\) is the quotient normal bundle. A Riemannian metric on
\(N\) identifies it with the orthogonal normal bundle. On a paracompact
smooth base, a smooth positive bundle metric exists by partitions of
unity, so every short exact sequence of smooth vector bundles splits.
The splitting need not be natural without the chosen metric or other
extra data.

This is stronger than choosing complements separately at each point:
arbitrary pointwise choices need not vary continuously. Conversely
splitting does not make either summand trivial. A nontrivial line bundle
can be a direct summand of a trivial higher-rank bundle. Smooth
variation and a global frame are different requirements.

Hirsch proves the bundle splitting theorem as IV.2.2 (printed p.~95, PDF
p.~105), using partitions of unity, and defines the orthogonal structure
immediately afterward. This is actual proof coverage, not an analogy
inferred from the chapter title. The necessary new bridges for his
development are multivariable calculus, the inverse/implicit-function
theorems, general topology, manifolds, partitions of unity, and then
approximation and transversality.

The Hessian of a smooth real function at a critical point is a symmetric
bilinear form on the tangent space. Its inertia survives coordinate
change, which explains why Sylvester's law underlies the Morse index.
The Morse lemma and Morse theory additionally require analytic and
topological arguments; diagonalization alone does not prove their local
normal form or global conclusions. Hirsch VI is a destination for those
arguments.

Exterior algebra can build differential forms once smooth bundles exist,
but Hirsch's preface expressly chooses a route without tensors or
differential forms. It would be misleading to force homological algebra,
algebraic geometry, or tensor calculus into his prerequisite chain. The
common foundation is careful linear geometry; the subsequent methods are
genuinely additional.

\section{A concrete destination: Hadamard and Williamson
foundations}\label{a-concrete-destination-hadamard-and-williamson-foundations}

\subsection{Scope and the exact
question}\label{scope-and-the-exact-question}

The confirmed existing direction is the paused Hadamard/Williamson
project in this repository. Its general existence question asks whether,
for every positive integer \(k\), there is a matrix

\[H\in\{-1,1\}^{4k\times4k},\qquad HH^{\mathsf T}=4kI_{4k}.\]

The prior work at commit
\texttt{2834fea31fb4ca7203531ca220e55a2fb7640119} records three rounds
of bounded research in a restricted family. Its proofs and limitations
are read here as the starting point. No prior benchmark is presented as
a fresh run, and no historical claim about the smallest missing order or
the worldwide current status is imported without a separate literature
audit. This chapter does not restart the computational search.

The structural chain actually needed is

\[\text{Gram matrices}\longrightarrow\text{tensor constructions},
\qquad
\text{finite fields}\longrightarrow\text{quadratic characters},
\]

\[\text{cyclic group algebras}\longrightarrow\text{Fourier decomposition}
\longrightarrow\text{correlation equations and lifting constraints}.
\]

It makes direct use of Xi II's metric and tensor material and Xi III's
groups, representations, and finite fields. General Ext theory, schemes,
and differential topology have no established role in this chain.

\subsection{Gram matrices and the order
restriction}\label{gram-matrices-and-the-order-restriction}

The identity \(HH^{\mathsf T}=nI\) says precisely that the rows are
mutually orthogonal in the usual positive real inner product and have
squared length \(n\). It also gives \(H^{-1}=n^{-1}H^{\mathsf T}\) and
\(|\det H|=n^{n/2}\), attaining the determinant bound proved in the
forms chapter. Orthogonality alone does not enforce sign entries; both
conditions are essential.

Multiplying rows or columns by signs and permuting them preserves the
Hadamard condition. Normalize the first row to all \(+1\). A second row
must have equally many signs, so \(n\) is even. When \(n>2\), permute
its positive entries into the first half and let a third row have \(a\)
positives in the first half and \(b\) in the second. Its products with
the first two rows give

\[2(a+b)=n,\qquad2(a-b)=0.\]

Hence \(n=4a\). This proves the necessary orders \(1,2\), or a multiple
of \(4\). It supplies no construction at each admissible order. The
proof is over the integers or reals; reducing the equation modulo \(2\)
would collapse all signs and lose the argument.

\subsection{Tensor closure and what it leaves
unsolved}\label{tensor-closure-and-what-it-leaves-unsolved}

If \(H\) and \(K\) have orders \(a,b\), functoriality of the tensor
product and transpose gives

\[(H\otimes K)(H\otimes K)^{\mathsf T}
=(HH^{\mathsf T})\otimes(KK^{\mathsf T})=abI_{ab}.\]

Its entries remain signs. Tensoring with
\(S=\begin{pmatrix}1&1\\1&-1\end{pmatrix}\) gives
\(\begin{pmatrix}H&H\\H&-H\end{pmatrix}\) and yields all powers-of-two
orders from \([1]\). This is a genuine infinite construction theorem. It
does not produce all multiples of four without additional seeds or
constructions. The universal tensor identities explain why the
construction transports orthogonality, while the sign-entry requirement
explains why arbitrary orthogonal matrices are not sufficient seeds.

\subsection{Finite fields and a character identity proved by
counting}\label{finite-fields-and-a-character-identity-proved-by-counting}

Let \(\mathbb F_q\) be a finite field of odd order. Its nonzero squares
form an index-two subgroup: the squaring map has exactly two preimages
\(x,-x\) for each nonzero square. Define the quadratic character
\(\chi:\mathbb F_q\to\{0,\pm1\}\) with \(\chi(0)=0\), \(+1\) on nonzero
squares, and \(-1\) on nonsquares. The quotient-group description proves
multiplicativity, and equal coset sizes prove \(\sum_x\chi(x)=0\).

The multiplicative group of a finite field is cyclic. One short proof
lets \(e\) be the least common multiple of its element orders. In a
finite abelian group, combine elements of maximal prime-power orders to
obtain an element of order \(e\). Every nonzero field element is a root
of \(X^e-1\), so there are at most \(e\) of them, whereas Lagrange gives
\(e\le q-1\). Thus \(e=q-1\). It follows that
\(\chi(-1)=(-1)^{(q-1)/2}\).

For \(a\ne0\), the crucial correlation identity is

\[\sum_x\chi(x)\chi(x+a)=-1.\]

Count pairs \((x,y)\) with \(y^2=x(x+a)\). For fixed \(x\), their number
is \(1+\chi(x(x+a))\), so the total is \(q\) plus the displayed sum.
Scale \(x=az\), \(y=aw\) and rewrite as

\[(2z+1-2w)(2z+1+2w)=1.\]

Since \(2\) is invertible, each nonzero first factor determines a unique
second factor and a unique \((z,w)\). There are \(q-1\) pairs, proving
the sum. For \(a=0\) the corresponding sum is \(q-1\). This argument
uses a finite field of odd order, not the special fact that its elements
are residue classes modulo a prime.

\subsection{Paley constructions with explicit
signs}\label{paley-constructions-with-explicit-signs}

Index \(Q\) by \(x,y\in\mathbb F_q\) and put \(Q_{xy}=\chi(y-x)\). The
character identities give

\[Q\mathbf1=0,\qquad QQ^{\mathsf T}=qI-J,
\qquad Q^{\mathsf T}=\chi(-1)Q.\]

If \(q\equiv3\pmod4\), set

\[H=\begin{pmatrix}1&\mathbf1^{\mathsf T}\\
\mathbf1&Q-I\end{pmatrix}.\]

Its entries are signs, its top row is orthogonal to every lower row
because \(1+(Q-I)\mathbf1=0\), and its lower Gram block is

\[J+(Q-I)(Q^{\mathsf T}-I)
=J+qI-J-Q-Q^{\mathsf T}+I=(q+1)I.\]

Therefore \(HH^{\mathsf T}=(q+1)I\). The prior study used the equivalent
sign convention with lower first column \(-\mathbf1\) and core \(Q+I\).
The conventions have been derived separately here so a copied sign does
not silently spoil orthogonality.

For \(q\equiv1\pmod4\), define the symmetric conference matrix

\[C=\begin{pmatrix}0&\mathbf1^{\mathsf T}\\\mathbf1&Q\end{pmatrix}.\]

It has zero diagonal, off-diagonal signs, and \(C^2=qI\). The block
matrix

\[\mathcal H=\begin{pmatrix}C+I&C-I\\C-I&-C-I\end{pmatrix}\]

has only sign entries. The diagonal Gram block is
\((C+I)^2+(C-I)^2=2(q+1)I\) and the off-diagonal block vanishes by
commutation. It is Hadamard of order \(2(q+1)\).

These proofs add a short prime-power and conference-matrix development
to the prior study's prime, \(q\equiv3\pmod4\) proof. They are classical
constructions, not new results or claims that the earlier scripts
implemented them. Finite fields of every prime-power order exist (Xi III
Theorems 5.40--5.45, printed pp.~260--261, PDF pp.~266--267); these
pages and their proofs were inspected visually. Small exact instances
are checked separately in this deliverable, without resuming an
enumeration campaign.

\subsection{Circulants as an integral group
algebra}\label{circulants-as-an-integral-group-algebra}

Fix \(m\ge1\) and let \(P\) be the cyclic shift with \(P_{ij}=1\) when
\(j=i+1\pmod m\). The map

\[\mathbb Z[z]/(z^m-1)\longrightarrow M_m(\mathbb Z),
\qquad f\longmapsto f(P)\]

is injective: the first row of \(\sum_{j=0}^{m-1}a_jP^j\) is its
coefficient vector. Its image is the circulant algebra, so circulants
commute. Define \(f^\iota(z)=f(z^{-1})\). Transpose corresponds to this
involution. For a real sequence \(x\) with polynomial
\(X(z)=\sum_jx_jz^j\), the coefficient of \(z^s\) in \(XX^\iota\) is

\[\operatorname{PAF}_x(s)=\sum_jx_jx_{j+s}.\]

This follows by collecting terms with exponent difference \(s\);
autocorrelation is even in \(s\). It is convolution with the reversed
sequence, so the group algebra retains the integral coefficients needed
for exact verification.

If \(x_j=x_{-j}\), then \(X=X^\iota\) and the circulant is symmetric.
Four such sign sequences form a Williamson quadruple when

\[AA^{\mathsf T}+BB^{\mathsf T}+CC^{\mathsf T}+DD^{\mathsf T}
=4mI_m,\]

equivalently their four autocorrelations sum to zero at every nonzero
shift. Symmetry permits replacing the four Gram matrices by
\(A^2,B^2,C^2,D^2\). Commutation and symmetry then make

\[W=\begin{pmatrix}
A&B&C&D\\-B&A&-D&C\\-C&D&A&-B\\-D&-C&B&A
\end{pmatrix}\]

Hadamard of order \(4m\). Multiplying blocks, every diagonal Gram block
is the sum of the four squares, while an off-diagonal block has terms
such as \(-AB+BA-CD+DC=0\). Transpose and commutation hypotheses are
indispensable in that cancellation. Existence in the Williamson family
is sufficient for a Hadamard matrix of that order; it is not a necessary
characterization of all Hadamard matrices.

\subsection{Fourier characters are Chinese remainder
coordinates}\label{fourier-characters-are-chinese-remainder-coordinates}

Let \(k\) have characteristic not dividing \(m\) and contain a primitive
\(m\)th root \(\zeta\). The roots of \(z^m-1\) are distinct, so

\[k[z]/(z^m-1)\xrightarrow{\ \sim\ }k^m,
\qquad f\longmapsto(f(1),f(\zeta),\ldots,f(\zeta^{m-1}))\]

is a Chinese remainder isomorphism. Its inverse is the finite Fourier
formula

\[a_j=\frac1m\sum_{r=0}^{m-1}f(\zeta^r)\zeta^{-rj}.\]

To verify it, the geometric sum \(\sum_r\zeta^{r\ell}\) equals \(m\)
when \(m\mid\ell\) and \(0\) otherwise. In the regular representation,
the vector \((1,\zeta^r,\ldots,\zeta^{(m-1)r})\) is an eigenvector of
\(P\) with eigenvalue \(\zeta^r\). Thus all circulants diagonalize in
the same basis; this is the group representation decomposition into
characters.

Over \(\mathbb C\), divide these vectors by \(\sqrt m\) to get a unitary
basis. A real symmetric circulant then has real eigenvalues, since
\(X(\zeta^{-r})=X(\zeta^r)=\overline{X(\zeta^r)}\). The correlation
identity becomes

\[\sum_{X\in\{A,B,C,D\}}|X(\zeta^r)|^2=4m\qquad(0\le r<m).\]

The frequency-zero equation is the row-sum identity
\(r_a^2+r_b^2+r_c^2+r_d^2=4m\). A single row or partial tuple must
satisfy the corresponding upper bounds, but separate feasible spectral
values need not come from any common sign coefficients.

The field assumptions explain failures: over \(\mathbb Q\) the algebra
need not split into one-dimensional factors, although it splits after
adjoining \(\zeta\); if the characteristic divides \(m\), \(z^m-1\) has
repeated roots and this semisimple Fourier argument breaks. Positivity
and absolute-square inequalities require the real/complex setting; they
have no automatic meaning over a finite field.

\subsection{A quotient test is not a lift
theorem}\label{a-quotient-test-is-not-a-lift-theorem}

For \(d\mid m\), the group homomorphism \(C_m\to C_d\) gives an algebra
quotient sending \(z\) to the generator of order \(d\). The compressed
coefficient is \(y_r=\sum_{j\equiv r\ (d)}x_j\). Applying the quotient
to the four correlation identity gives

\[\sum_x C_{y_x}C_{y_x}^{\mathsf T}=4mI_d.\]

The right side remains \(4m\), not \(4d\), and compressed coefficients
are integer sums rather than signs. This is a necessary consequence of a
full solution. The quotient discards coefficients, so its converse needs
a separate common-lift theorem.

The prior proof gives an exact illustrative failure at \(m=9\). The
normalized symmetric masks \((2,10,10,13)\), where bit \(j-1\) negates
positions \(\pm j\), compress at \(d=3\) to

\[(3,1,1),\quad(3,-1,-1),\quad(3,-1,-1),\quad(-1,-1,-1).\]

Their total compressed autocorrelation is \((36,0,0)\) and their
row-square sum is \(36\), so both proper cyclic quotients of \(C_9\)
pass. But their uncompressed autocorrelation at shift \(1\) sums to
\(-8\). This shows that these particular lifts fail; it alone would not
prove that every other lift of the compressed multiset fails. The prior
independent review separately records a stronger no-common-lift fixture:
the compressed multiset

\[(-1,-1,-1),\quad(-1,1,1),\quad(-1,3,3),\quad(3,-1,-1).\]

It too has total compressed autocorrelation \((36,0,0)\). There are 16
normalized symmetric sign rows at core nine; among their four-row
multisets, exactly 27 map to this specified compressed multiset. The new
exact checker independently reconstructs and rejects all 27 by their
full autocorrelations. Thus no common lift exists \textbf{in this
normalized symmetric Williamson family}. That is stronger evidence than
one failed lift, and still says nothing about nonexistence of arbitrary
Hadamard matrices of order 36.

Even positive spectral residuals need not have a sign lift. At \(m=7\),
let three sequences have negative positions \(\{\pm1\}\), \(\{\pm2\}\),
and \(\{\pm3\}\). Their squared circulants sum to \(20I+J\), leaving
\(8I-J\) toward the target \(28I\). This residual is positive definite:
its eigenvalues are \(1\) on constants and \(8\) on their orthogonal
complement. Nevertheless the odd Williamson product constraint below
forces the fourth sequence to be all positive; its row-sum square is
\(49\), while the remaining row-square budget is \(1\). No fourth
Williamson row exists. A positive real Gram factor does not necessarily
satisfy the required discrete sign pattern.

\subsection{Why four colors appear in the prior
reduction}\label{why-four-colors-appear-in-the-prior-reduction}

Now let \(m=2h+1\) be odd and normalize every sequence by \(x_0=1\).
Write \(T(z)=\sum_{j=0}^{m-1}z^j\) and \(N_x(z)=\sum_{x_j=-1}z^j\), so
\(X=T-2N_x\). The identities \(T^2=mT\) and \(TN_x=|N_x|T\), substituted
into the Williamson equation and divided coefficientwise by \(4\) in the
integral group algebra, give

\[\sum_xN_x^2+(m-R)T=m,\qquad R=\sum_x|N_x|.\]

Each \(|N_x|\) is even by symmetry and normalization. Modulo \(2\),
squaring removes cross terms and sends \(z^j\) to \(z^{2j}\). Since
doubling permutes residues modulo odd \(m\), each nonzero position
contains an odd number of negative entries among the four rows.
Equivalently

\[a_jb_jc_jd_j=-1\qquad(j\ne0).\]

This is the odd Williamson product theorem in the prior proof. Its
hypotheses are structural restrictions; it is not a constraint on
arbitrary Hadamard matrices.

Let \(v_c=\mathbf1-2e_c\in\mathbb Z^4\), \(c=1,2,3,4\). The eight
four-entry sign columns with product \(-1\) are uniquely
\(\varepsilon v_c\) with \(\varepsilon=\pm1\), and

\[v_c^{\mathsf T}v_d=4[c=d],\qquad
\mathbf1^{\mathsf T}v_c=2.\]

Thus each nonzero column has a color \(\kappa_j\) and sign
\(\varepsilon_j\), reflected at opposite indices. Write \([j]=\{j,-j\}\)
for a nonzero opposite-index class, represented in \(\{1,\ldots,h\}\).
For a nonzero shift \(t\), define \(E_t\) as follows, with indices
modulo \(m\). Partition \(j\in\mathbb Z/m\mathbb Z\) into the orbits of
the involution \(j\mapsto-t-j\). Remove the two-element orbit
\(\{0,-t\}\) and the unique fixed point solving \(2j=-t\). For each
remaining two-element orbit, choose either member \(j\) and include one
unordered edge \(\{[j],[j+t]\}\) in \(E_t\), retaining one edge per
orbit. Replacing \(j\) by \(-t-j\) swaps the endpoints, so the
definition is independent of the representative. Both endpoints are
nonzero and distinct; equality of the classes would force either \(t=0\)
or the removed fixed point. There are \(h-1\) edges, counted with their
orbit multiplicities.

In the correlation sum, the removed pair involving zero contributes
\(4\varepsilon_t\); the fixed term contributes \(4\). Each remaining
pair with endpoint classes \(u,v\) contributes
\(8\varepsilon_u\varepsilon_v[\kappa_u=\kappa_v]\). Hence exact
correlation is equivalent to

\[\varepsilon_t+1+
2\sum_{\{u,v\}\in E_t}\varepsilon_u\varepsilon_v
[\kappa_u=\kappa_v]=0.\]

If \(c_t\) counts its monochromatic edges, reducing this equation modulo
\(4\) gives \(\varepsilon_t=-(-1)^{c_t}\). All signs are therefore
forced by a complete coloring. The full integer equations still have to
hold; the parity consequence alone is insufficient.

Permuting the four block rows merely permutes color names. Naming colors
in order of first appearance, with at most four colors, gives a unique
restricted-growth representative of each orbit. This is a quotient by
block permutations, including repeated blocks. It is not classification
up to all Hadamard matrix equivalences, nor a quotient by arbitrary
coordinate changes.

\subsection{Boolean polynomials and the precise relaxation
boundary}\label{boolean-polynomials-and-the-precise-relaxation-boundary}

Encode each color minus one by two bits \(a_i,b_i\). In the Boolean
algebra

\[B=\mathbb F_2[a_i,b_i\ (1\le i\le h)]/
(a_i^2-a_i,b_i^2-b_i),\]

every function on the bit cube has a unique squarefree polynomial.
Induct on the last bit: any function is \(f_0+x(f_1-f_0)\), which proves
existence and uniqueness recursively. Equality of colors has polynomial

\[\delta_{uv}=(1+a_u+a_v)(1+b_u+b_v).\]

It is \(1\) exactly when both bits agree. Therefore

\[n_t=1+\sum_{\{u,v\}\in E_t}\delta_{uv},\qquad
\varepsilon_t=(-1)^{n_t}\]

has degree at most two in the bits. Substituting a partial assignment
preserves this degree bound. Write a residual \(n_i=\alpha_i+p_i\) with
constant \(\alpha_i\) and zero-constant polynomial \(p_i\), and put
\(\eta_i=(-1)^{\alpha_i}\), \(\xi_p=(-1)^p\). Then

\[\varepsilon_i=\eta_i\xi_{p_i},\qquad
\varepsilon_u\varepsilon_v=\eta_u\eta_v\xi_{p_u+p_v}.\]

These are exact identities of functions. Equal residual polynomials
therefore represent the same unknown phase even before its sign is
known. Summing their integer coefficients first can reveal cancellation
that separate sign ranges miss. Constants and linear terms must be
included when deciding equality; equal quadratic parts alone do not
suffice.

The exponents are computed in \(B\), of characteristic two. The
correlation coefficients and their sums are computed in \(\mathbb Z\).
The map \(p\mapsto(-1)^p\) converts addition of Boolean functions into
multiplication of signs; it is not a ring map allowing integer
coefficients to be reduced modulo two.

After grouping, a lag expression can have the exact form
\(K+\sum_pc_p\xi_p\). Treating distinct \(\xi_p\) as independent signs
gives a safe superset of its attainable values, hence a sound rejection
test when zero is absent. Presence of zero is only a necessary
condition. Relations such as \(\xi_{p+q}=\xi_p\xi_q\) remain; even
preserving all such linear phase relations need not enforce the
quadratic origin. For example, the values of \(a,b,ab\) cannot be
\((0,0,1)\), though these three polynomials are linearly independent
over \(\mathbb F_2\). A prospective stronger method must specify which
extra dependencies it keeps and prove its rejection rule.

The prior work's core-11 false-positive prefix is a concrete warning
that separately feasible row and lag relaxations can lack a common
completion. Its existing audit, not a new run here, supports that
bounded observation. Neither the algebraic reduction nor the word
``exact'' makes every downstream relaxation complete.

\subsection{What to learn next for this
direction}\label{what-to-learn-next-for-this-direction}

\begin{longtable}[]{@{}
  >{\raggedright\arraybackslash}p{(\columnwidth - 4\tabcolsep) * \real{0.3333}}
  >{\raggedright\arraybackslash}p{(\columnwidth - 4\tabcolsep) * \real{0.3333}}
  >{\raggedright\arraybackslash}p{(\columnwidth - 4\tabcolsep) * \real{0.3333}}@{}}
\toprule\noalign{}
\begin{minipage}[b]{\linewidth}\raggedright
Tool needed for the study
\end{minipage} & \begin{minipage}[b]{\linewidth}\raggedright
Source in the collection
\end{minipage} & \begin{minipage}[b]{\linewidth}\raggedright
Supplement supplied or still needed
\end{minipage} \\
\midrule\noalign{}
\endhead
\bottomrule\noalign{}
\endlastfoot
Orthogonality, Gram matrices, positive residuals & Xi II §§3.1--3.5, 3.8
& Gram-factor versus sign-factor distinction; examples above \\
Kronecker products and tensor identities & Xi II §6.2; Hungerford IV.5;
Grillet XI.5--6 & Sylvester closure proof above \\
Cyclic groups, quotient rings, polynomial actions & Xi III §§1.2,
4.2--4.4; Hungerford III--IV; Grillet III, VIII & Integral
cyclic-algebra involution and correlation dictionary above \\
Finite fields and multiplicative characters & Xi III §5.6; Hungerford
V.5; Grillet IV.1.6 and V.1.3 & Quadratic-character sum and both Paley
block proofs above \\
Character decomposition and invariant splitting & Xi III §§3.1--3.5;
Grillet IX.2, IX.7--9 & Fourier inversion as CRT, with
splitting/characteristic hypotheses above \\
Compression and common lifting & Quotients and maps throughout the core
& A real additional theorem would have to couple sign lifts; none is
assumed \\
Quadratic Boolean constraints & Polynomial rings, quotients, functions
on finite fields & Boolean normal form and phase dependence above;
algorithmic feasibility remains a separate subject \\
\end{longtable}

The strongest useful next reading task is to reconstruct each displayed
dictionary and counterexample, then read the pinned prior proofs with
the coefficient rings visible. A later research proposal should target a
specific lost dependency, a genuinely wider construction family, or a
proved lift criterion. Such proposals are prospective; no improvement in
search complexity or progress toward universal existence is claimed by
this foundational chapter.

\section{Invariant complements as a computable extension
obstruction}\label{invariant-complements-as-a-computable-extension-obstruction}

\textbf{Status and scope.} The extension classification, Sylvester
criterion, trace duality, Bézout argument, and Smith-form consequences
below are classical mathematics. The organization, connecting
derivations, and worked examples are original exposition, with no claim
of research novelty. This chapter builds on the operator-module
translation in Chapter 2 and the cyclic extension construction in
Chapter 6. It supplies the missing machinery for an actual
invariant-complement problem: compute the obstruction, exhibit a
certificate when it is nonzero, construct a complement when it vanishes,
and determine how a nonzero obstruction joins cyclic blocks. Every main
claim is proved below; no new textbook proof-comparison claim is made
here.

Throughout, \(k\) is a field and \(R=k[t]\). Let \(U,W\) be
finite-dimensional \(k\)-spaces, with operators \(A:U\to U\) and
\(B:W\to W\). Write \(U_A,W_B\) for the corresponding \(R\)-modules.
Several arguments work without finite dimension; those extensions are
identified explicitly. Matrices act on column vectors, and \(J_n(0)\)
has ones on the \textbf{superdiagonal}.

\subsection{The datum an invariant-subspace problem actually
contains}\label{the-datum-an-invariant-subspace-problem-actually-contains}

Suppose an operator \(T\) on \(E\) preserves a specified subspace \(U\).
Its restriction is \(A\); its induced operator on \(W=E/U\) is \(B\).
This gives an exact sequence with specified endpoint identifications:

\[0\longrightarrow U_A\xrightarrow{i}E_T\xrightarrow{q}W_B
\longrightarrow0.\tag{1}\]

A \(k\)-linear section \(s:W\to E\) exists because \(k\) is a field. The
map \((u,w)\mapsto i(u)+s(w)\) identifies the underlying vector space
with \(U\oplus W\). The defect of \(s\) from being an \(R\)-linear map
is

\[Ts(w)-s(Bw)=i(Cw),\]

for a unique \(k\)-linear map \(C:W\to U\). Thus

\[T_C=\begin{pmatrix}A&C\\0&B\end{pmatrix}.\tag{2}\]

Every \(C\) gives an extension, but \(C\) depends on the chosen section.
Put

\[H=\operatorname{Hom}_k(W,U),\qquad
\delta_{A,B}:H\longrightarrow H,\quad \delta(h)=Ah-hB.\tag{3}\]

Replace \(s\) by \(s+i h\). Direct substitution gives

\[C\longmapsto C+Ah-hB=C+\delta(h).\tag{4}\]

Equivalently, with \(P_h=\begin{pmatrix}1&h\\0&1\end{pmatrix}\),

\[P_h^{-1}T_CP_h=T_{C+\delta(h)}.\tag{5}\]

This conjugacy fixes the embedded \(U\) and induces the identity on the
quotient \(W\). It is much more restricted than arbitrary similarity.

For clarity about signs, an extension isomorphism written in the
\emph{direction} \(F_h:E_C\to E_D\), \(F_h(u,w)=(u+hw,w)\), satisfies

\[F_hT_C=T_DF_h\quad\Longleftrightarrow\quad D=C-\delta(h).\tag{6}\]

There is no discrepancy between (4) and (6): \(P_h\) maps the new
coordinate model to the old one. Both say that equivalent extensions
have the same class in \(H/\delta H\).

\subsubsection{Complements, sections, and the solution
set}\label{complements-sections-and-the-solution-set}

Any vector-space complement to \(U\oplus0\) is the graph

\[\Gamma_h=\{(hw,w):w\in W\}\]

of a unique \(h\in H\): the quotient projection restricts to an
isomorphism from the complement to \(W\). Applying \(T_C\) to a graph
vector shows

\[T_C(hw,w)=(Ahw+Cw,Bw)\in\Gamma_h
\quad\Longleftrightarrow\quad Ah-hB=-C.\tag{7}\]

Consequently the following are equivalent:

\begin{enumerate}
\def\labelenumi{\arabic{enumi}.}
\tightlist
\item
  \(U\) has a \(T_C\)-invariant complement;
\item
  \begin{enumerate}
  \def\labelenumii{(\arabic{enumii})}
  \tightlist
  \item
    admits an \(R\)-linear section;
  \end{enumerate}
\item
  the Sylvester equation \(\delta(h)=-C\) has a solution;
\item
  \([C]=0\) in \(H/\delta H\).
\end{enumerate}

If \(h_0\) is one solution, every solution is

\[h_0+f,\qquad f\in\ker\delta
=\operatorname{Hom}_R(W_B,U_A).\tag{8}\]

The set of invariant complements is therefore either empty or an affine
space modeled on the intertwiners \(W_B\to U_A\). It has no
distinguished origin until a complement is chosen. An invariant
complement is unique exactly when a solution exists and this intertwiner
space is zero.

\subsection{Why the cokernel is precisely the first Ext
group}\label{why-the-cokernel-is-precisely-the-first-ext-group}

The preceding calculation is already a classification by
endpoint-preserving extension equivalence: any extension admits a linear
section; any equivalence fixing both endpoints must have the form
\(F_h\); (6) describes exactly when it is \(R\)-linear. Two further
derivations explain the group law and connect this classification to a
free resolution.

\subsubsection{A canonical one-step free
resolution}\label{a-canonical-one-step-free-resolution}

Let \(F=R\otimes_k W\). Define

\[\begin{aligned}
d_B:F&\longrightarrow F,&
d_B(t^j\otimes w)&=t^{j+1}\otimes w-t^j\otimes Bw,\\
\varepsilon:F&\longrightarrow W_B,&
\varepsilon(t^j\otimes w)&=B^jw.
\end{aligned}\]

Then

\[0\longrightarrow F\xrightarrow{d_B}F
\xrightarrow{\varepsilon}W_B\longrightarrow0\tag{9}\]

is exact. Here is a proof without a determinant calculation. If a
nonzero vector polynomial \(z=\sum_{j=0}^m t^j\otimes w_j\) has
\(w_m\ne0\), then the coefficient of \(t^{m+1}\) in \(d_Bz\) is \(w_m\).
Hence \(d_B\) is injective. Surjectivity of \(\varepsilon\) follows from
\(\varepsilon(1\otimes w)=w\). Finally,

\[t^j\otimes w-1\otimes B^jw
=d_B\left(\sum_{\ell=0}^{j-1}t^{j-1-\ell}\otimes B^\ell w\right)
\qquad(j\ge1).\tag{10}\]

Thus every vector polynomial is congruent modulo
\(\operatorname{im}d_B\) to its constant evaluation; exactly the
elements with evaluation zero lie in the image. Choosing a basis of
\(W\) shows that \(F\) is free over \(R\). In matrix notation
\(d_B=tI-B\).

An \(R\)-map \(F\to U_A\) is uniquely specified by a \(k\)-map
\(h:W\to U\), its value on \(1\otimes W\). Precomposition with \(d_B\)
sends \(h\) to

\[w\longmapsto Ah(w)-h(Bw)=\delta(h)(w).\]

Applying \(\operatorname{Hom}_R(-,U_A)\) to the resolution gives

\[0\longrightarrow\operatorname{Hom}_R(W_B,U_A)
\longrightarrow H\xrightarrow{\delta}H
\longrightarrow\operatorname{Ext}^1_R(W_B,U_A)
\longrightarrow0,\tag{11}\]

and \(\operatorname{Ext}^j_R(W_B,U_A)=0\) for \(j\ge2\). In particular,

\[\boxed{\operatorname{Ext}^1_{k[t]}(W_B,U_A)
\cong\operatorname{coker}(h\mapsto Ah-hB).}\tag{12}\]

The sign in this identification can be checked on the extension itself.
Extend its chosen section \(s\) to the \(R\)-map
\(\widetilde s:F\to E_T\), \(t^j\otimes w\mapsto T^js(w)\). Then
\(\widetilde s\,d_B\) takes values in \(i(U)\) and its value on
\(1\otimes w\) is \(i(Cw)\). Therefore the resolution cocycle is \(C\),
not an unspecified sign of \(C\).

This resolution and the cokernel computation remain valid for arbitrary
\(k\)-vector spaces, assuming the usual existence of bases. Finite
dimension will matter for the trace pairing and the dimension counts
below.

\subsubsection{Baer addition becomes addition of upper-right
blocks}\label{baer-addition-becomes-addition-of-upper-right-blocks}

Take \(E_C\) and \(E_D\). Their fiber product over \(W\) consists of
triples \((u_1,u_2,w)\), and its operator is

\[(u_1,u_2,w)\longmapsto(Au_1+Cw,Au_2+Dw,Bw).\]

The Baer sum pushes out the kernel \(U\oplus U\) along addition
\((u_1,u_2)\mapsto u_1+u_2\). Concretely quotient the fiber product by
\(\{(u,-u,0):u\in U\}\). The resulting coordinate \((u_1+u_2,w)\) has
operator \(T_{C+D}\). Hence

\[[E_C]+[E_D]=[E_{C+D}].\tag{13}\]

Pushout along \(\lambda1_U\) similarly gives \([E_{\lambda C}]\). Thus
(12) respects the \(k\)-vector-space structure as well as the set of
extension classes. The zero class is the split sequence.

\subsection{A complete certificate for failure: trace-dual
intertwiners}\label{a-complete-certificate-for-failure-trace-dual-intertwiners}

A failed attempt to solve a matrix equation does not by itself explain
why it fails. Here the dual equation supplies a finite certificate.

Put \(K=\operatorname{Hom}_k(U,W)\). There is a perfect bilinear pairing

\[K\times H\longrightarrow k,\qquad
\langle g,C\rangle=\operatorname{tr}_W(gC)
=\operatorname{tr}_U(Cg).\tag{14}\]

The two traces agree because both are \(\sum_{i,j}g_{ij}C_{ji}\) in
bases. The pairing is nondegenerate: a matrix unit in one rectangular
matrix space reads a single entry of the other. This proof works in
every characteristic. In particular it does \textbf{not} require
\(\operatorname{char}k\nmid\dim U\) or
\(\operatorname{char}k\nmid\dim W\).

For \(h\in H\),

\[\begin{aligned}
\langle g,\delta(h)\rangle
&=\operatorname{tr}_W(gAh-ghB)\\
&=\operatorname{tr}_W((gA-Bg)h).
\end{aligned}\tag{15}\]

Thus the dual linear map is

\[\delta^\vee:K\to K,\qquad g\longmapsto gA-Bg.\tag{16}\]

Its kernel consists of the intertwiners in the \textbf{reverse
direction}:

\[\ker\delta^\vee=\operatorname{Hom}_R(U_A,W_B).\tag{17}\]

Equation (15) proves that every \(\delta(h)\) is annihilated by those
intertwiners. Conversely, a vector \(C\) outside
\(\operatorname{im}\delta\) can be separated from that subspace by a
linear functional: extend a basis of the image by \(C\) and prescribe
the functional to be zero on the image and one on \(C\). The perfect
pairing represents that functional by some \(g\). Equation (15) then
forces \(gA=Bg\). Hence

\[\boxed{\ C\in\operatorname{im}\delta
\iff \operatorname{tr}_W(gC)=0
\text{ for every }g:U\to W\text{ with }gA=Bg.\ }\tag{18}\]

In particular, a single \(g\) satisfying

\[gA=Bg,\qquad \operatorname{tr}_W(gC)\ne0\tag{19}\]

certifies that \(U\) has no invariant complement in \(E_C\). A basis of
the reverse-intertwiner space suffices to test all witnesses. A complete
vanishing test, followed by solving \(Ah-hB=-C\), gives a constructive
complement.

The quotient pairing is perfect and identifies

\[\operatorname{Ext}^1_R(W_B,U_A)^*
\cong\operatorname{Hom}_R(U_A,W_B),\tag{20}\]

or equivalently, in finite dimension,
\(\operatorname{Ext}^1_R(W_B,U_A)\cong
\operatorname{Hom}_R(U_A,W_B)^*\). This identification uses the
specified polynomial variable \(t\) through (9). It is independent of
bases and of the section. Its endpoint orientation is not optional: maps
\(W\to U\) parametrize changes of a splitting, while maps \(U\to W\)
test the obstruction.

For completeness, the expected compatibility with endpoint maps follows
from traces. Pushing an extension out along \(a:U\to U'\) replaces \(C\)
by \(aC\) and pairs with \(g':U'\to W\) as \(\operatorname{tr}(g'aC)\).
Pulling back along \(b:W'\to W\) replaces \(C\) by \(Cb\) and
\(\operatorname{tr}_{W'}(gCb)=\operatorname{tr}_W(bgC)\). These
identities explain the natural dual pairing; they do not identify a
vector space with its own dual without additional data.

\subsection{Coprime annihilators give an explicit
inverse}\label{coprime-annihilators-give-an-explicit-inverse}

Suppose \(a,b\in k[t]\) satisfy

\[a(A)=0,\qquad b(B)=0,\qquad u(t)a(t)+v(t)b(t)=1.\tag{21}\]

One often stops after using Bézout to prove that a complement exists. A
divided-difference identity constructs the actual solution of every
Sylvester equation.

For \(b(t)=\sum_{j=0}^d b_jt^j\), define the polynomial

\[Q_b(x,y)=\sum_{j=1}^d b_j\sum_{\ell=0}^{j-1}x^{j-1-\ell}y^\ell.
\tag{22}\]

It satisfies

\[(x-y)Q_b(x,y)=b(x)-b(y).\tag{23}\]

The notation \((b(x)-b(y))/(x-y)\) abbreviates this polynomial identity;
no inversion of \(x-y\) is being assumed. On \(H\), let \(L(h)=Ah\) and
\(R_B(h)=hB\). These operators commute, even though there is no
meaningful product \(AB\) when \(U\ne W\). Therefore (23) gives

\[(L-R_B)Q_b(L,R_B)=L_{b(A)}-R_{b(B)}=L_{b(A)}.\tag{24}\]

Evaluating (21) at \(A\) yields \(v(A)b(A)=1_U\). Every operator
involved in (24) commutes with left multiplication by \(v(A)\). It
follows that

\[\boxed{\ \delta^{-1}(D)=v(A)
\sum_{j=1}^d b_j\sum_{\ell=0}^{j-1}A^{j-1-\ell}DB^\ell.\ }\tag{25}\]

Multiplying on either side by \(\delta\) gives the identity; this is a
two-sided inverse, not merely a recipe for selected right-hand sides.
The unique invariant complement is the graph of \(h=-\delta^{-1}(C)\).

For example, if \(A^2=0\) and \(B=1_W\), use \(a=t^2\), \(b=t-1\), and
\(t^2-(t+1)(t-1)=1\). Equation (25) becomes

\[\delta^{-1}(D)=-(1_U+A)D,\qquad
h=(1_U+A)C.\]

Directly, \(A(1_U+A)C-(1_U+A)C=-C\). No Jordan basis is needed, and the
formula also works in characteristic two.

A less trivial exact check is \(A=N\), \(B=1+M\), where \(N^2=M^2=0\).
Take \(a=t^2\), \(b=(t-1)^2\) and

\[(3-2t)t^2+(1+2t)(t-1)^2=1.\]

Then (25) reduces to

\[\delta^{-1}(D)=-D-ND+DM+2NDM.\tag{26}\]

The integers in this identity are interpreted in \(k\); no division by
two or three occurs.

\subsubsection{A finite geometric inverse for distinct Jordan
eigenvalues}\label{a-finite-geometric-inverse-for-distinct-jordan-eigenvalues}

Suppose \(A=\lambda1_U+N\) and \(B=\mu1_W+M\), with \(N^a=0\),
\(M^b=0\), and \(d=\lambda-\mu\ne0\). On \(H\), put \(K(D)=ND-DM\). Its
two summands commute, so the binomial expansion shows \(K^{a+b-1}=0\):
each term has either at least \(a\) left factors \(N\) or at least \(b\)
right factors \(M\). Since \(\delta=d1_H+K\),

\[\delta^{-1}=\sum_{r=0}^{a+b-2}(-1)^r d^{-r-1}K^r,\]

where

\[K^r(D)=\sum_{\ell=0}^r(-1)^\ell\binom r\ell
N^{r-\ell}DM^\ell.\]

Multiplication by \(d1_H+K\) telescopes to the identity, with the last
term zero by nilpotence. This is a finite algebraic identity, with no
convergence assumption. It works in every characteristic when \(d\ne0\);
cancellation of binomial coefficients in small characteristic can
shorten the nilpotence bound without invalidating it. Formula (26) is
its case \(d=-1\), \(a=b=2\).

\subsubsection{A coordinate-free description of the same
complement}\label{a-coordinate-free-description-of-the-same-complement}

Let \(E_T\) be any extension in (1) under (21). Since \(b(B)=0\),
\(b(T)E\subseteq U\); since \(a(A)=0\), it follows that \(a(T)b(T)=0\).
The candidate complement is

\[K=\ker b(T).\tag{27}\]

Indeed \(b(A)\) is invertible with inverse \(v(A)\), so \(K\cap U=0\).
For any lift \(e\) of \(w\in W\), write \(b(T)e=i(u)\) and replace \(e\)
by \(e-i(v(A)u)\). This new lift lies in \(K\). Thus \(q|_K:K\to W\) is
an isomorphism and \(E=U\oplus K\).

Every invariant complement is isomorphic to \(W_B\), hence is killed by
\(b(T)\) and lies in \(K\); it must equal \(K\). This proves uniqueness
without finite-dimensional dimension counting. Both this proof and (25)
work for arbitrary \(k\)-spaces with the stated coprime annihilators.
Over a commutative base ring the algebraic identities also hold, but
arbitrary extensions need not admit a section over that base ring; the
block classification then needs that additional hypothesis.

\subsection{What a shared spectral factor does and does not
imply}\label{what-a-shared-spectral-factor-does-and-does-not-imply}

In finite dimension, take \(a=m_A\) and \(b=m_B\). The following are
equivalent:

\begin{itemize}
\tightlist
\item
  \(m_A\) and \(m_B\) are coprime;
\item
  \(\delta_{A,B}\) is invertible;
\item
  every extension of \(W_B\) by \(U_A\) splits, with a unique invariant
  complement;
\item
  the spectra of \(A\) and \(B\) are disjoint over an algebraic closure
  of \(k\).
\end{itemize}

Coprimality implies invertibility by (25). Invertibility is equivalent
to \(\ker\delta=\operatorname{coker}\delta=0\) by finite-dimensional
linear algebra, and (8), (12) give the extension assertion. For the
converse to coprimality, suppose an irreducible \(p\) divides both
minimal polynomials. The primary cyclic decompositions contain blocks
\(R/(p^a)\) and \(R/(p^b)\) for some \(a,b\ge1\). A nonzero map from the
latter to the former sends \(1\) to \(p^{\max(a-b,0)}\) modulo \(p^a\).
Extending it by zero on the other blocks gives a nonzero \(W_B\to U_A\)
intertwiner, so \(\delta\) is not injective. Finally, two polynomials
over a field share an irreducible factor exactly when they have a common
root in an algebraic closure, proving the spectral equivalence.

The conclusion concerns \textbf{all possible \(C\)}. If a factor is
shared, some right-hand sides are obstructed, but a particular \(C\) may
still lie in \(\operatorname{im}\delta\). In particular \(C=0\) always
splits, and nonzero coboundaries split as well. A common eigenvalue is
not a certificate for a particular nonsplitting extension; (19) is such
a certificate.

Checking only eigenvalues \emph{in the ground field} is insufficient.
Over \(\mathbb R\) let

\[A=B=\begin{pmatrix}0&-1\\1&0\end{pmatrix},\qquad C=1_2.\]

Neither operator has a real eigenvalue, so their real-eigenvalue sets
are both empty. Nonetheless \(g=1_2\) is a reverse intertwiner and
\(\operatorname{tr}(gC)=2\ne0\): this extension does not split. The
shared factor is \(t^2+1\).

Nor does separate semisimplicity of the endpoints suffice. For
\(A=B=(\lambda)\), \(\delta=0\); every \(C\ne0\) gives a nonsplit Jordan
block of size two. Semisimplicity of the middle module would force every
submodule to split, but semisimplicity is not preserved by arbitrary
extensions.

\subsection{Cyclic extensions expose exactly how blocks are
joined}\label{cyclic-extensions-expose-exactly-how-blocks-are-joined}

Let \(f,g\in R\) be nonzero monic polynomials, with \(U=R/(f)\) and
\(W=R/(g)\). The resolution \(0\to R\xrightarrow{g}R\to R/(g)\to0\)
gives

\[\operatorname{Ext}^1_R(R/(g),R/(f))
\cong(R/(f))/g(R/(f))
\cong R/(f,g)=R/(\gcd(f,g)).\tag{28}\]

The following presentation retains the endpoint maps, which (28) alone
could hide. Let \(u\) be the specified generator of \(U\), and choose a
lift \(e\) of the specified generator of \(W\). There is a polynomial
\(c\) such that

\[fu=0,\qquad ge=cu.\tag{29}\]

Conversely these relations define

\[E_c=\operatorname{coker}
\begin{pmatrix}f&-c\\0&g\end{pmatrix}.\tag{30}\]

Columns are relation vectors in the ordered generators \((u,e)\). The
map \(R/(f)\to E_c\), \(1\mapsto u\), is injective: if \((r,0)\) is a
relation combination \(x(f,0)+y(-c,g)\), then \(yg=0\); since \(R\) is a
domain and \(g\ne0\), \(y=0\), so \(f\mid r\). The quotient by this copy
of \(U\) is \(R/(g)\) generated by the image of \(e\).

Replacing \(e\) by \(e+hu\) changes \(c\) to \(c+gh\) modulo \(f\).
Therefore the extension class is precisely \(c\bmod(f,g)\), with the
specified cyclic generators understood. The sequence splits exactly when
that residue is zero. This is a classifying statement with fixed
endpoints; it is not merely an assertion about the isomorphism type of
\(E_c\).

\subsubsection{Middle modules: a smaller
invariant}\label{middle-modules-a-smaller-invariant}

The first Smith factor of the matrix in (30) is the monic polynomial

\[d_1=\gcd(f,g,c),\]

because its entries generate that ideal. Its determinant is \(fg\), so
the second factor is \(d_2=fg/d_1\). Since \(d_1\) divides both \(f\)
and \(g\), \(d_1^2\mid fg\), and hence \(d_1\mid d_2\), as required.
Thus

\[E_c\cong R/(d_1)\oplus R/(d_2).\tag{31}\]

A unit factor contributes the zero module. This is an application of the
Smith theorem already proved in Chapter 2. Row and column operations
classify the middle module; they need not preserve its specified
submodule, quotient, or their generators.

Consequently the residue \(c\bmod\gcd(f,g)\) classifies the extension,
while \(\gcd(f,g,c)\) determines the middle-module isomorphism type in
this cyclic situation. Different residues can have the same gcd and
hence the same middle module.

\subsubsection{Primary powers and an explicit Smith
reduction}\label{primary-powers-and-an-explicit-smith-reduction}

Fix a monic irreducible \(p\) and positive integers \(a,b\). Take
\(f=p^a\), \(g=p^b\). Define \(v_p(0)=\infty\) and

\[m=\min(a,b),\qquad \nu=\min(a,b,v_p(c)).\tag{32}\]

Changing \(c\) by a multiple of \(p^m\) leaves this truncated valuation
unchanged. Equations (28) and (31) become

\[\operatorname{Ext}^1_R(R/(p^b),R/(p^a))\cong R/(p^m),\tag{33}\]

\[\boxed{\ E_c\cong R/(p^\nu)\oplus
R/(p^{a+b-\nu}).\ }\tag{34}\]

The following row and column operations exhibit the mechanism in (34).
When \(\nu<m\), write \(c=p^\nu q\), \(P=p^{a-\nu}\), \(Q=p^{b-\nu}\).
Then \((P,q)=1\), so choose \(s,r\in R\) with \(sP+rq=1\). Both matrices
below have determinant one, and

\[\begin{pmatrix}1&0\\rQ&1\end{pmatrix}
\begin{pmatrix}P&-q\\0&Q\end{pmatrix}
\begin{pmatrix}s&q\\-r&P\end{pmatrix}
=\begin{pmatrix}1&0\\0&PQ\end{pmatrix}.\tag{35}\]

Multiplying the middle matrix by \(p^\nu\) produces exactly the Smith
factors \(p^\nu\) and \(p^{a+b-\nu}\). For \(\nu=m\) the class is zero:
\(c\) is divisible by \(p^m\). If \(a\le b\), add \((c/p^a)\) times
column one to column two; if \(b\le a\), add \((c/p^b)\) times row two
to row one. This removes \(-c\) and gives the split diagonal
presentation, reordered if necessary. These operations also cover
\(c=0\).

The three boundary behaviors deserve to be kept separate:

\begin{itemize}
\tightlist
\item
  \textbf{Zero class:} \(\nu=m\), including \(c=0\) and every
  \(c\in p^mR\). The middle module is \(R/(p^a)\oplus R/(p^b)\), and the
  specified sequence splits.
\item
  \textbf{Unit residue:} \(p\nmid c\), so \(\nu=0\). The first Smith
  factor is \(1\) and \(E_c\cong R/(p^{a+b})\). Both cyclic lengths join
  into one. Here ``unit'' means a unit modulo \(p^m\); \(c\) need not be
  a constant unit of \(k[t]\).
\item
  \textbf{Intermediate valuation:} \(0<\nu<m\). A cyclic piece of length
  \(\nu\) remains, while the other has length \(a+b-\nu\). For
  \(p=t-\lambda\), these are the two Jordan block sizes. For irreducible
  \(p\) of degree greater than one, these are elementary-divisor
  exponents; their \(k\)-dimensions are multiplied by \(\deg p\).
\end{itemize}

The formulas keep total \(k\)-dimension \((a+b)\deg p\) fixed. They do
not say that arbitrary simultaneous matrix conjugation preserves the
chosen extension class.

\subsection{Two size-two Jordan blocks: the full quotient
calculation}\label{two-size-two-jordan-blocks-the-full-quotient-calculation}

Take \(U=W=k^2\) and

\[A=B=N=\begin{pmatrix}0&1\\0&0\end{pmatrix},\qquad
C=\begin{pmatrix}c_{11}&c_{12}\\c_{21}&c_{22}\end{pmatrix}.\]

For \(h=\begin{pmatrix}x&y\\z&w\end{pmatrix}\),

\[Nh-hN=\begin{pmatrix}z&w-x\\0&-z\end{pmatrix}.\tag{36}\]

Thus \(\operatorname{im}\delta\) is spanned by \(E_{12}\) and
\(E_{11}-E_{22}\). The two quotient coordinates are

\[\alpha=c_{21},\qquad \beta=c_{11}+c_{22}=\operatorname{tr}C.\tag{37}\]

Every class has the unique representative

\[C_{\alpha,\beta}=\begin{pmatrix}\beta&0\\\alpha&0\end{pmatrix}.
\tag{38}\]

For example, choose \(z=c_{22}\) and \(w-x=-c_{12}\) in (36); then
\(C+\delta(h)=C_{\alpha,\beta}\). This works in every characteristic,
including two.

The reverse intertwiners are exactly \(g=r1+sN\). Their pairings are

\[\operatorname{tr}(gC)=r\beta+s\alpha.\tag{39}\]

So the abstract trace criterion recovers precisely the two explicit
quotient coordinates. There is a complement if and only if
\(\alpha=\beta=0\). In that case the solution space has dimension two,
since \(\ker\delta=k1\oplus kN\).

To connect this computation with cyclic presentations, let \(u_1,u_2\)
be the basis of \(U\), with \(Nu_2=u_1\); the cyclic generator is
\(u=u_2\). Similarly choose \(w_1,w_2\) in \(W\) and let \(e\) be the
coordinate lift \((0,w_2)\). As \(B^2=0\),

\[T_C^2e=ACw_2+CBw_2
=(c_{11}+c_{22})u_1+c_{21}u_2
=(\alpha+\beta t)u.\tag{40}\]

Thus the identification with the cyclic Ext group is

\[[C]\longmapsto\alpha+\beta t\pmod{t^2},\qquad
\operatorname{Ext}^1_R(R/(t^2),R/(t^2))\cong R/(t^2).\tag{41}\]

The three module types now follow from (34):

\begin{longtable}[]{@{}
  >{\raggedright\arraybackslash}p{(\columnwidth - 6\tabcolsep) * \real{0.2500}}
  >{\raggedright\arraybackslash}p{(\columnwidth - 6\tabcolsep) * \real{0.2500}}
  >{\raggedright\arraybackslash}p{(\columnwidth - 6\tabcolsep) * \real{0.2500}}
  >{\raggedright\arraybackslash}p{(\columnwidth - 6\tabcolsep) * \real{0.2500}}@{}}
\toprule\noalign{}
\begin{minipage}[b]{\linewidth}\raggedright
Extension coordinates
\end{minipage} & \begin{minipage}[b]{\linewidth}\raggedright
Representative \(C\)
\end{minipage} & \begin{minipage}[b]{\linewidth}\raggedright
Middle-module Jordan type
\end{minipage} & \begin{minipage}[b]{\linewidth}\raggedright
\((\operatorname{rank}T,\operatorname{rank}T^2,\operatorname{rank}T^3)\)
\end{minipage} \\
\midrule\noalign{}
\endhead
\bottomrule\noalign{}
\endlastfoot
\(\alpha=\beta=0\) & \(0\) & \(J_2(0)\oplus J_2(0)\) & \((2,0,0)\) \\
\(\alpha=0\), \(\beta\ne0\) & \(\beta E_{11}\) & \(J_3(0)\oplus J_1(0)\)
& \((2,1,0)\) \\
\(\alpha\ne0\) & \(\beta E_{11}+\alpha E_{21}\) & \(J_4(0)\) &
\((3,2,1)\) \\
\end{longtable}

These are all possibilities, over every field. Several comparisons
prevent misreadings of the table:

\begin{itemize}
\tightlist
\item
  \(C=E_{12}\ne0\) nevertheless represents the split class, since
  \(E_{12}=\delta(\operatorname{diag}(0,1))\). A complement is the graph
  of \(h=-\operatorname{diag}(0,1)\).
\item
  \(E_{11}\) and \(E_{22}\) are the same extension class: their
  difference is \(\delta(E_{21})\).
\item
  \(E_{21}\) and \(E_{21}+E_{11}\) are distinct extension classes, with
  cyclic residues \(1\) and \(1+t\), but both have middle module
  \(R/(t^4)\). This distinction persists over \(\mathbb F_2\).
\item
  \(C=1_2\) gives the residue \(2t\). The middle module is
  \(J_3\oplus J_1\) if \(\operatorname{char}k\ne2\), but the extension
  is split in characteristic two. In that characteristic
  \(1_2=E_{11}-E_{22}=\delta(E_{21})\).
\end{itemize}

The last example does not invalidate the trace criterion. The pairing on
\emph{all} rectangular matrices is nondegenerate in every
characteristic; a particular trace, or the pairing restricted to a
smaller subspace, can vanish.

\subsection{Self-extensions and the commutant: dimensions without false
identities}\label{self-extensions-and-the-commutant-dimensions-without-false-identities}

Since \(\delta:H\to H\) is an endomorphism of a finite-dimensional
space, rank--nullity gives

\[\dim_k\operatorname{Ext}^1_R(W_B,U_A)
=\dim_k\operatorname{Hom}_R(W_B,U_A).\tag{42}\]

This dimension equality alone supplies no preferred isomorphism between
its two sides. Trace duality supplies a different, precise statement:
the dual of Ext is the \textbf{reverse} Hom space in (20).

If \(U_A=\bigoplus_iR/(f_i)\) and \(W_B=\bigoplus_jR/(g_j)\), the block
decomposition of \(\delta\), or finite additivity of the resolution,
gives

\[\operatorname{Ext}^1_R(W_B,U_A)
\cong\bigoplus_{i,j}R/(\gcd(f_i,g_j)).\tag{43}\]

The displayed coordinates depend on the chosen cyclic decompositions and
cyclic generators. The resulting dimension is intrinsic:

\[\dim_k\operatorname{Ext}^1_R(W_B,U_A)
=\sum_{i,j}\deg\gcd(f_i,g_j).\tag{44}\]

For a single module with elementary divisors \(p^{a_{p,i}}\), this
yields

\[\dim_k\operatorname{Ext}^1_R(V_T,V_T)
=\dim_k\operatorname{End}_R(V_T)
=\sum_p\deg(p)\sum_{i,j}\min(a_{p,i},a_{p,j}).\tag{45}\]

There is also the canonical duality, relative to the fixed \(t\),

\[\operatorname{Ext}^1_R(V_T,V_T)
\cong\operatorname{End}_R(V_T)^*.\tag{46}\]

One must not replace the dual in (46) by the commutant itself merely by
writing a trace pairing \emph{inside} the commutant. For \(T=N=J_2(0)\),
that commutant is \(k1\oplus kN\), and

\[\operatorname{tr}\bigl(N(r1+sN)\bigr)=0\quad\text{for all }r,s.\]

Its restricted trace form is degenerate in every characteristic. Yet
\(N\) detects the nonzero extension represented by \(E_{21}\) because
\(\operatorname{tr}(NE_{21})=1\). The duality pairs a commutant element
with a \emph{quotient class of arbitrary matrices}, not necessarily with
another commutant element. This is why choosing a direct complement to
\(\operatorname{im}\delta\) in the full matrix space is an extra choice.

A useful related calculation gives the same quotient another meaning.
Over the dual numbers \(k[\epsilon]/(\epsilon^2)\), a first-order change
of basis \(1+\epsilon h\) transforms \(T+\epsilon C\) into

\[(1-\epsilon h)(T+\epsilon C)(1+\epsilon h)
=T+\epsilon(C+Th-hT).\tag{47}\]

Thus first-order perturbations modulo changes of basis that reduce to
the identity have the same quotient as self-extensions. This explains
the appearance of a commutator cokernel, but does not classify finite
perturbations or give a multiplication on Ext matching composition in
the commutant.

\subsection{The determinant bridge: resultants and integral
obstructions}\label{the-determinant-bridge-resultants-and-integral-obstructions}

For \(m=\dim U\) and \(n=\dim W\), column-by-column vectorization gives

\[[\delta_{A,B}]=1_n\otimes A-B^{\mathsf T}\otimes1_m.\tag{48}\]

With the convention
\(\operatorname{Res}(f,g)=\prod_{f(\alpha)=0}g(\alpha)\) for monic
polynomials, roots counted with multiplicity, one has

\[\boxed{\det\delta_{A,B}
=\operatorname{Res}(\chi_A,\chi_B)
=\prod_{i=1}^m\prod_{j=1}^n(\alpha_i-\beta_j).}\tag{49}\]

Here \(\alpha_i,\beta_j\) are the characteristic roots in an algebraic
closure. To prove this without diagonalizability or separability, extend
the field to an algebraic closure and triangularize \(A,B\). Such a
triangular form follows by repeatedly choosing an eigenvector and
passing to the quotient; it does not require an eigenbasis. In (48),
\(B^{\mathsf T}\) is lower triangular, so the matrix is lower block
triangular with diagonal blocks \(A-\beta_j1_m\). Its determinant is
\(\prod_j\det(A-\beta_j1_m)=\prod_{i,j}(\alpha_i-\beta_j)\).
Determinants are unchanged by scalar extension, proving (49) over \(k\).
Repeated roots simply repeat factors in the product.

Now let \(U=\mathbb Z^m\) and \(W=\mathbb Z^n\), with integral operators
\(A,B\). There are two different questions: splitting after tensoring
with \(\mathbb Q\), and splitting as \(\mathbb Z[t]\)-modules. Because
\(W\) is free over \(\mathbb Z\), every underlying exact sequence of
abelian groups ending in \(W\) splits. Thus the block and section
calculations still apply, with

\[\delta_{\mathbb Z}:\operatorname{Hom}_{\mathbb Z}(W,U)
\longrightarrow\operatorname{Hom}_{\mathbb Z}(W,U),\quad
h\longmapsto Ah-hB.\tag{50}\]

Moreover (9) remains a free \(\mathbb Z[t]\)-resolution: its
leading-term and telescoping proofs work verbatim, and
\(\mathbb Z[t]\otimes_{\mathbb Z}W\cong\mathbb Z[t]^n\) is free.
Consequently, for these endpoints,

\[\operatorname{Ext}^1_{\mathbb Z[t]}(W_B,U_A)
\cong\operatorname{coker}\delta_{\mathbb Z}.\tag{51}\]

The first argument being free over \(\mathbb Z\) is essential to this
free resolution argument. We are not asserting that \(\mathbb Z[t]\) is
a PID or that every \(\mathbb Z[t]\)-module has a resolution of length
one.

If the resultant in (49) is nonzero, \(\delta_{\mathbb Q}\) is
invertible. The integer matrix \(\delta_{\mathbb Z}\) is then injective
with finite cokernel, and integer Smith form gives

\[\left|\operatorname{coker}\delta_{\mathbb Z}\right|
=|\det\delta_{\mathbb Z}|
=|\operatorname{Res}(\chi_A,\chi_B)|.\tag{52}\]

Indeed a Smith diagonal \(d_1,\ldots,d_{mn}\) identifies the cokernel
with \(\bigoplus_i\mathbb Z/d_i\mathbb Z\); the product of its finite
orders is \(|\prod_i d_i|=|\det\delta|\). Thus a nonzero resultant gives
unique rational splitting, whereas resultant \(\pm1\) gives unique
integral splitting for every \(C\). For a particular integral \(C\) with
nonzero resultant, the unique rational solution
\(h=-\delta_{\mathbb Q}^{-1}C\) is integral exactly when that extension
splits integrally. More generally, Smith coordinates turn the integral
solvability test into divisibility by the \(d_i\).

\subsubsection{The smallest lattice
defect}\label{the-smallest-lattice-defect}

Take \(A=(0)\), \(B=(2)\), \(C=(1)\), so

\[T=\begin{pmatrix}0&1\\0&2\end{pmatrix},\qquad
\delta_{\mathbb Z}(h)=-2h.\tag{53}\]

The rational complement is uniquely the graph of \(h=1/2\). But its
class in \(\operatorname{coker}\delta_{\mathbb Z}=\mathbb Z/2\mathbb Z\)
is nonzero, so the specified copy \(U=\mathbb Z(1,0)\) has no invariant
integral complement. The integral points on the other eigenline are
generated by \((1,2)\), and

\[[\mathbb Z^2:\mathbb Z(1,0)+\mathbb Z(1,2)]=2.\tag{54}\]

The quotient projection sends this eigenline lattice to \(2\mathbb Z\),
not onto \(W=\mathbb Z\). Independent rational eigenvectors have
produced a sublattice of finite index, not a direct-sum decomposition of
the specified lattice. The resultant is
\(\operatorname{Res}(t,t-2)=-2\), whose absolute value \(2\) is the
order of the obstruction.

This also pinpoints the boundary of trace witnesses. Here the reverse
equation \(gA=Bg\) forces \(g=0\) over both \(\mathbb Q\) and
\(\mathbb Z\). There is no nonzero trace witness over either ring, yet
the integral extension is nonsplit. Linear functionals into
\(\mathbb Z\) cannot detect the torsion cokernel \(\mathbb Z/2\).
Reducing (53) modulo two does expose it, but complete integral tests use
the Smith divisibility conditions, including higher prime powers when
present.

This is the same arithmetic distinction that appears in integral CRT:
polynomials can be coprime over \(\mathbb Q[t]\) while the inverse or
Bézout coefficients require denominators. The resultant measures a
finite failure of invertibility on a lattice. Its nonvanishing is the
field-level splitting test; its absolute value need not be one.

\subsection{Exercises with worked
answers}\label{exercises-with-worked-answers}

\textbf{Exercise 1 --- a split extension with many complements.} Let
\(A=B=J_2(0)\) and \(C=E_{12}\). Find every invariant complement, and
verify invariance without using the Ext terminology.

\emph{Answer.} Equation (36) must equal \(-E_{12}\), so \(z=0\) and
\(w-x=-1\), while \(x,y\) are arbitrary. The complements are the graphs
of

\[h=\begin{pmatrix}x&y\\0&x-1\end{pmatrix}.\]

For these matrices \(Ah+C=hB\), hence \(T_C(hw,w)=(hBw,Bw)\) belongs to
the same graph. Distinct \(h\) give distinct complements because
projection to \(W\) fixes the graph parameter.

\textbf{Exercise 2 --- produce an obstruction witness.} Let \(U=k^2\),
\(W=k\), \(A=J_2(0)\), \(B=0\), and \(C(1)=(c_1,c_2)^{\mathsf T}\).
Determine exactly when a complement exists. When it does not, provide a
one-line certificate.

\emph{Answer.} Write \(h(1)=(x,y)^{\mathsf T}\). Then
\(\delta(h)(1)=(y,0)^{\mathsf T}\), so a complement exists exactly when
\(c_2=0\), with \(y=-c_1\) and \(x\) arbitrary. The reverse intertwiner
\(g=(0,1)\) satisfies \(gA=0=Bg\), and \(\operatorname{tr}(gC)=c_2\). It
certifies nonsplitting when \(c_2\ne0\). This example also makes the
direction \(g:U\to W\) visible from matrix sizes.

\textbf{Exercise 3 --- compute a complement without solving entries.}
Suppose \(A^3=0\) and \(B=1_W\). Find the unique invariant complement
for arbitrary \(C:W\to U\).

\emph{Answer.} \(\delta(h)=(A-1_U)h\), and \((1_U-A)(1_U+A+A^2)=1_U\).
Thus \(h=(1_U+A+A^2)C\) solves \(\delta(h)=-C\). The same identity
proves uniqueness. This is the specialization of the divided-difference
formula to \(b=t-1\).

\textbf{Exercise 4 --- valuation controls the middle module.} Let
\(U=R/(p^3)\) and \(W=R/(p^5)\) for a monic irreducible \(p\). Classify
the possible middle-module types, and count extension classes of each
type when \(k=\mathbb F_q\) and \(\deg p=d\).

\emph{Answer.} The extension group is \(R/(p^3)\); put \(Q=q^d\). There
is one zero class, giving \(R/(p^3)\oplus R/(p^5)\). For \(\nu=0,1,2\),
the type is \(R/(p^\nu)\oplus R/(p^{8-\nu})\), and the number of classes
of exact valuation \(\nu\) is

\[|p^\nu R/(p^3)|-|p^{\nu+1}R/(p^3)|
=Q^{3-\nu}-Q^{2-\nu}.\]

These counts refer to endpoint-preserving extension classes. Their sum
with the split class is \(Q^3\). Each of the four rows of this
classification specifies one middle-module isomorphism type, so counting
middle modules would give four, not \(Q^3\).

\textbf{Exercise 5 --- the characteristic-two trap.} Over
\(\mathbb F_2\), let \(A=B=J_2(0)\) and \(C=1_2\). Exhibit an invariant
complement and reconcile it with the fact that \(C\) is invertible.

\emph{Answer.} \(\delta(E_{21})=1_2=-1_2\) in characteristic two. The
graph of \(E_{21}\) is therefore invariant. More explicitly its basis
consists of \((u_2,w_1)\) and \((0,w_2)\); \(T\) sends the first to zero
and the second to the first. Invertibility of the off-diagonal map is
unrelated to whether its class in \(H/\delta H\) is zero.

\textbf{Exercise 6 --- a larger self-extension count.} If
\(T=J_3(0)\oplus J_1(0)\), determine the dimension of its self-extension
space. Explain why the answer is not yet a classification of middle
modules.

\emph{Answer.} Formula (45) gives \(3+1+1+1=6\). These are six
independent parameters for extensions with the two copies of \(V_T\)
identified in advance. Changing endpoint automorphisms or forgetting the
inclusion and quotient can identify different extension classes.
Determining the isomorphism types of the resulting eight-dimensional
middle modules is a further problem; the number six is not a number of
Jordan types.

\subsection{Hypotheses, failed shortcuts, and reconstructive
output}\label{hypotheses-failed-shortcuts-and-reconstructive-output}

\begin{longtable}[]{@{}
  >{\raggedright\arraybackslash}p{(\columnwidth - 4\tabcolsep) * \real{0.3333}}
  >{\raggedright\arraybackslash}p{(\columnwidth - 4\tabcolsep) * \real{0.3333}}
  >{\raggedright\arraybackslash}p{(\columnwidth - 4\tabcolsep) * \real{0.3333}}@{}}
\toprule\noalign{}
\begin{minipage}[b]{\linewidth}\raggedright
Claim or construction
\end{minipage} & \begin{minipage}[b]{\linewidth}\raggedright
Hypotheses actually used
\end{minipage} & \begin{minipage}[b]{\linewidth}\raggedright
Failed shortcut or boundary
\end{minipage} \\
\midrule\noalign{}
\endhead
\bottomrule\noalign{}
\endlastfoot
Block form and section-change classification & A field, or an explicitly
split underlying sequence over the base ring & An arbitrary exact
sequence over \(\mathbb Z\) need not admit such a section;
\(0\to\mathbb Z/2\to\mathbb Z/4\to\mathbb Z/2\to0\) is an example \\
One-step resolution (9) & \(R=k[t]\) and a \(k\)-space \(W\) & It is a
one-variable resolution; it does not automatically compute extensions
with two commuting operators \\
Complete trace witnesses and perfect Ext--Hom duality &
Finite-dimensional spaces over a field & Infinite-dimensional
endomorphisms have no general trace; a trace on the commutant alone can
be degenerate \\
Explicit inverse (25) & Stated annihilators and a Bézout identity
\(ua+vb=1\) & Shared factors prevent a universal inverse but do not
obstruct every individual \(C\) \\
Spectral disjointness criterion & Finite dimension, eigenvalues taken
over an algebraic closure & Empty real-eigenvalue sets can hide a shared
irreducible quadratic \\
Cyclic Ext and Smith formulas & Nonzero \(f,g\in k[t]\); fixed cyclic
endpoint generators for the residue coordinate & Smith operations
classify the middle module and may change endpoint data \\
Jordan interpretation of (34) & \(p=t-\lambda\) & Higher-degree
irreducibles give rational elementary divisors rather than Jordan blocks
over \(k\) \\
Equality of commutant and self-Ext dimensions & Finite dimension & Equal
dimensions do not identify the spaces or their algebraic structures; the
natural statement here is duality \\
Integral cokernel has order equal to absolute resultant & Finite free
integral endpoints and nonzero resultant & Rational splitting can leave
a finite lattice-index obstruction; rational trace tests do not detect
it \\
\end{longtable}

The reconstruction route is short enough to use repeatedly. Record the
specified submodule and quotient; choose any linear section and extract
\(C\). Compute its class modulo \(Ah-hB\), or test it against reverse
intertwiners. If the class vanishes, solve once and add every forward
intertwiner to obtain all complements. If annihilators are coprime, (25)
solves immediately and the complement is uniquely \(\ker b(T)\). If an
obstruction remains, cyclic coordinates and Smith reduction explain what
it does to the middle module. This keeps four different outputs
distinct: a yes-or-no splitting test, a witness or complement, an
extension class with fixed endpoints, and the coarser isomorphism type
obtained after forgetting those endpoints.

\textbf{Local validation.} The companion script
\texttt{check\_sylvester.py} checks the section/equivalence signs, both
sides of the Bézout inverse, the trace-adjoint identity, all size-two
Jordan extensions over \(\mathbb F_2\) and \(\mathbb F_3\),
representative polynomial Smith reductions over \(\mathbb Q[t]\), and
the integral resultant/index bridge. These finite checks supplement the
field-independent proofs; they are not substitutes for the hypotheses or
the general arguments.

\section{Repairing linear maps: spectral projection, separability, and
averaging}\label{repairing-linear-maps-spectral-projection-separability-and-averaging}

An arbitrary vector-space complement rarely respects an operator or a
group action. There are two ways to repair it. One can solve a separate
equation for the complement, as in the Sylvester calculation, or
construct a single operation that repairs every linear map between
modules for the algebra in question. This chapter constructs the second
operation explicitly.

The distinction matters. A particular extension may split even when no
repair defined by a separability tensor exists. Conversely, a repair
operation for modules over a quotient algebra says nothing about
extensions whose middle term is not a module over that quotient.

\subsection{Two identities that repair every linear
map}\label{two-identities-that-repair-every-linear-map}

Let \(A\) be an associative unital algebra over a field \(k\). For the
moment, suppose that finitely many elements \(a_i,b_i\in A\) satisfy

\[
\sum_i a_i b_i=1,\qquad
\sum_i c a_i\otimes b_i=\sum_i a_i\otimes b_i c
\quad(c\in A).                                      \tag{A1}
\]

The tensor products in this displayed identity are tensor products of
\(k\)-vector spaces. The order of the two factors in \(b_i c\) matters.
For left \(A\)-modules \(M,N\) and a \(k\)-linear map \(f:M\to N\),
define

\[
\mathcal A_{M,N}(f)(v)=\sum_i a_i f(b_i v).            \tag{A2}
\]

No dimension assumption on \(M\) or \(N\) is needed: the sum defining
the repair is finite.

\textbf{Repair lemma.} The operation \(\mathcal A_{M,N}\) is a
\(k\)-linear projection

\[
\operatorname{Hom}_k(M,N)\longrightarrow
\operatorname{Hom}_A(M,N).
\]

If \(u:M'\to M\) and \(v:N\to N'\) are \(A\)-linear, then

\[
\mathcal A_{M',N'}(vfu)=v\,\mathcal A_{M,N}(f)\,u.    \tag{A3}
\]

\textbf{Proof.} Apply the \(k\)-linear map \(a\otimes b\mapsto a f(bw)\)
to the second identity in (A1). It gives

\[
c\,\mathcal A(f)(w)=\sum_i a_i f(b_i c w)
                  =\mathcal A(f)(cw),
\]

so the repaired map is \(A\)-linear. If \(f\) was already \(A\)-linear,

\[
\mathcal A(f)(w)=\sum_i a_i b_i f(w)=f(w).
\]

These two assertions give \(\mathcal A^2=\mathcal A\) and its stated
image. For (A3), move the \(A\)-linear maps \(u,v\) through the scalar
actions in (A2). \(\square\)

Consider an exact sequence of \(A\)-modules

\[
0\longrightarrow U\xrightarrow{\iota}V
 \xrightarrow{\pi}W\longrightarrow0.
\]

Choose a \(k\)-linear section \(s:W\to V\) of \(\pi\). Then

\[
\pi\mathcal A(s)=\mathcal A(\pi s)
               =\mathcal A(1_W)=1_W.
\]

Thus \(\mathcal A(s)\) is an \(A\)-linear section. A \(k\)-linear
retraction \(r:V\to U\) of \(\iota\) is repaired in the same way:
\(\mathcal A(r)\iota=1_U\). Consequently every such exact sequence
splits as a sequence of \(A\)-modules. Ordinary vector spaces admit the
initial section by choosing a basis; finite-dimensionality is not the
source of this splitting theorem.

The identities also explain why averaging a projection works in
Maschke's proof. If \(P:V\to V\) has image in an invariant submodule
\(U\) and is the identity on \(U\), each term in (A2) has image in
\(U\), and \(\mathcal A(P)|_U=1_U\). These facts imply
\(\mathcal A(P)^2=\mathcal A(P)\).

\textbf{Boundary.} Averaging is not an algebra homomorphism on
\(\operatorname{End}_k(V)\), and it does not preserve arbitrary
idempotents. Over \(\mathbb Q\), let \(T=\operatorname{diag}(1,-1)\) and
average by \(\mathcal A(f)=(f+TfT)/2\). It removes the off-diagonal
entries. Hence

\[
\mathcal A(E_{12})=\mathcal A(E_{21})=0,\qquad
\mathcal A(E_{12}E_{21})=E_{11}.
\]

Moreover, the projection
\(Q=\frac12\begin{pmatrix}1&1\\1&1\end{pmatrix}\) averages to \(I/2\),
which is not a projection. Its image was not \(T\)-invariant. The valid
repair argument uses the invariant target and the identity on that
target; it is not an assertion about arbitrary averages of idempotent
matrices.

\subsection{What the tensor means}\label{what-the-tensor-means}

Write

\[
e=\sum_i a_i\otimes b_i^{\mathrm{op}}
\quad\text{in}\quad A^e=A\otimes_k A^{\mathrm{op}}.
\]

The superscript records the opposite multiplication. In \(A^e\),

\[
(a\otimes b^{\mathrm{op}})(c\otimes d^{\mathrm{op}})
   =ac\otimes(db)^{\mathrm{op}}.
\]

It acts on \(A\) by \((a\otimes b^{\mathrm{op}})x=axb\). The
multiplication map \(\mu:A^e\to A\),
\(\mu(a\otimes b^{\mathrm{op}})=ab\), is a homomorphism of
\(A\)-bimodules, with the outer bimodule action on its source. In
general it is not a homomorphism of algebras.

The identities (A1) become

\[
\mu(e)=1,\qquad
(c\otimes1)e=(1\otimes c^{\mathrm{op}})e.             \tag{A4}
\]

They imply that \(\sigma(c)=(c\otimes1)e\) is an \(A\)-bimodule section
of \(\mu\). Conversely, such a section gives (A4) on putting
\(e=\sigma(1)\). An algebra with this property is called
\textbf{separable over \(k\)}. This is a condition on the algebra
together with its ground field.

The name \textbf{separability idempotent} is literal:

\[
\begin{aligned}
e^2
 &=\sum_i(a_i\otimes1)(1\otimes b_i^{\mathrm{op}})e\\
 &=\sum_i(a_i b_i\otimes1)e=e.
\end{aligned}
\]

It need not be central in \(A^e\). The outer-action identity (A4),
rather than centrality, is what the repair proof needs.

For commutative \(A\), the opposite algebra is just \(A\), and these
notations simplify. The next construction gives the tensor without
choosing an eigenbasis or passing to a splitting field.

\subsection{A divided difference constructs the polynomial
repair}\label{a-divided-difference-constructs-the-polynomial-repair}

Let \(m\in k[t]\) be monic of positive degree and suppose

\[
\gcd(m,m')=1.                                      \tag{A5}
\]

Set \(S=k[t]/(m)\) and identify

\[
S\otimes_kS=k[x,y]/(m(x),m(y)).
\]

The expression

\[
q(x,y)=\frac{m(x)-m(y)}{x-y}
\]

is a polynomial: division here abbreviates the polynomial identity
\(m(x)-m(y)=(x-y)q(x,y)\). The Euclidean algorithm supplies a polynomial
\(h\) with \(h(t)m'(t)\equiv1\pmod m\). Define

\[
e=q(x,y)h(y)
 =\frac{m(x)-m(y)}{x-y}\,m'(y)^{-1}
\quad\text{in }S\otimes_kS.                         \tag{A6}
\]

The inverse in this formula exists by (A5); it is an inverse in \(S\),
not a rational function to be evaluated at a zero of its denominator.

In \(S\otimes S\) we have \((x-y)e=(m(x)-m(y))h(y)=0\). For any
polynomial \(f\), \(f(x)-f(y)\) is divisible by \(x-y\), so

\[
(f(x)-f(y))e=0.
\]

This is precisely the second identity in (A4). The multiplication map
\(\mu\) substitutes \(x=y=\bar t\), and \(q(t,t)=m'(t)\); hence
\(\mu(e)=m'(\bar t)h(\bar t)=1\). Formula (A6) therefore proves all the
repair conclusions for \(S\)-modules.

More concretely, if \(T_M\) and \(T_N\) satisfy \(m(T_M)=m(T_N)=0\) and
\(e=\sum_{i,j}c_{ij}x^i y^j\), then

\[
\mathcal A(f)=\sum_{i,j}c_{ij}T_N^i f T_M^j.         \tag{A7}
\]

The result satisfies \(T_N\mathcal A(f)=\mathcal A(f)T_M\). It is a
polynomial construction in the two actions, defined over \(k\).

Over a splitting field \(K\), write \(m(t)=\prod_{j=1}^d(t-\lambda_j)\)
with distinct \(\lambda_j\), and let \(e_j(t)\) be the usual Lagrange
idempotents. At the pair of roots \((\lambda_r,\lambda_s)\), formula
(A6) is \(0\) when \(r\ne s\) and \(1\) when \(r=s\). The evaluation
description of \(K[t]/(m)\otimes_KK[t]/(m)\) therefore gives

\[
e=\sum_j e_j\otimes e_j,\qquad
\mathcal A(f)=\sum_j P_j^N fP_j^M,                 \tag{A8}
\]

where \(P_j^M=e_j(T_M)\) and \(P_j^N=e_j(T_N)\). The operation keeps the
maps between matching eigenspaces and deletes the maps between different
eigenspaces. This is the spectral meaning of the tensor, but the
divided-difference proof did not require roots in \(k\).

In the commutative case the tensor is unique. Indeed, if \(e,e'\) both
satisfy (A4), moving each factor of a term of \(e\) across \(e'\) gives
\(ee'=\mu(e)e'=e'\). Interchanging them gives \(e'e=e\), so
commutativity implies \(e=e'\). Interchanging the two tensor factors
produces another such tensor, which is consequently the same tensor.

\subsection{Three calculations that can be
reconstructed}\label{three-calculations-that-can-be-reconstructed}

For \(m=t^2-1\) and \(\operatorname{char}k\ne2\), \(q=x+y\) and
\((2y)^{-1}=y/2\) in \(S\). Thus

\[
e=\frac{1+xy}{2},\qquad
\mathcal A(f)=\frac{f+T_N fT_M}{2}.                 \tag{A9}
\]

The projections are \(P_\pm=(I\pm T)/2\), and (A8) becomes
\(P_+^NfP_+^M+P_-^NfP_-^M\).

For a non-diagonal example over \(\mathbb Q\), put

\[
T=\begin{pmatrix}1&2\\0&-1\end{pmatrix},\qquad
U=\mathbb Q e_1,\qquad
P=\begin{pmatrix}1&3\\0&0\end{pmatrix}.
\]

Then \(T^2=I\), \(P\) projects onto \(U\), and

\[
\mathcal A(P)=\frac{P+TPT}{2}
             =\begin{pmatrix}1&1\\0&0\end{pmatrix}.
\]

Its kernel is \(\mathbb Q(-1,1)\), the desired invariant complement.
Equivalently, the quotient has operator \(-1\); the ordinary section
\(s(1)=(0,1)\) repairs to \((-1,1)\).

For \(m=t^2+1\) over \(\mathbb R\), write \(J_M^2=J_N^2=-I\). Now
\((2y)^{-1}=-y/2\), so

\[
e=\frac{1-xy}{2},\qquad
\mathcal A(f)=\frac{f-J_N fJ_M}{2}.                \tag{A10}
\]

This extracts the complex-linear part of a real-linear map between
complex vector spaces. Its complementary part \((f+J_N fJ_M)/2\) is
complex-antilinear. On \(\mathbb C\), a real-linear map has the form
\(f(z)=az+b\bar z\); (A10) retains \(az\). A real eigenbasis is neither
available nor needed.

For a cubic example over \(\mathbb Q\), take \(m=t^3-2\). Since
\((3y^2)^{-1}=y/6\),

\[
e=\frac{2+x y^2+x^2y}{6},\qquad
\mathcal A(f)=\frac{2f+T_N fT_M^2+T_N^2fT_M}{6}.   \tag{A11}
\]

Multiplying (A11) on the appropriate side by \(T_N\) and \(T_M\) gives

\[
6\bigl(T_N\mathcal A(f)-\mathcal A(f)T_M\bigr)
=2T_Nf-2fT_M+T_N^3fT_M-T_NfT_M^3=0,
\]

where the final cancellation uses both cubic relations. The repair can
therefore be checked before knowing the eigenvalues of a cubic operator.

Characteristic \(2\) itself does not preclude a polynomial repair. For
\(m=t^2+t+1\) over \(\mathbb F_2\), the derivative is \(1\) and
\(q=x+y+1\). Thus \(e=1+x+y\) and \(\mathcal A(f)=f+T_N f+fT_M\). This
is a separable quadratic field example in the same characteristic where
\(t^2-1\) fails.

\subsection{The derivative is the exact obstruction to this
tensor}\label{the-derivative-is-the-exact-obstruction-to-this-tensor}

Condition (A5) is necessary as well as sufficient. A useful proof
exposes the role of derivations.

For a commutative \(k\)-algebra \(S\) and an \(S\)-module \(L\), a
\(k\)-derivation \(D:S\to L\) is a \(k\)-linear map satisfying

\[
D(ab)=aD(b)+bD(a).
\]

Suppose \(e=\sum_i a_i\otimes b_i\) satisfies (A4). Apply
\(a\otimes b\mapsto aD(b)\) to \((c\otimes1)e=(1\otimes c)e\). The two
sides are

\[
c\sum_i a_iD(b_i)
\quad\text{and}\quad
c\sum_i a_iD(b_i)+\left(\sum_i a_i b_i\right)D(c).
\]

Therefore \(D(c)=0\) for every \(c\). A separability tensor kills all
such derivations.

For \(S=k[t]/(m)\), however, formal differentiation defines

\[
D:S\longrightarrow k[t]/(m,m'),\qquad
D(\bar f)=\overline{f'}.                           \tag{A12}
\]

It is well-defined: replacing \(f\) by \(f+mh\) changes the derivative
by \(m'h+mh'\), which is zero in the target. It obeys the product rule.
If \(\gcd(m,m')\ne1\), its target is nonzero and \(D(\bar t)=1\ne0\).
Consequently no tensor (A4) exists. We have proved

\[
\boxed{\ k[t]/(m)\text{ is separable over }k
       \iff \gcd(m,m')=1.\ }                       \tag{A13}
\]

This is the elementary content of the differential formula
\(\Omega_{S/k}\cong S/(m')\,dt\): differentiating the relation
\(m(t)=0\) imposes exactly \(m'(t)\,dt=0\). No differential-module
formalism is needed for the proof (A12).

For the smallest failure, take the dual numbers
\(D_0=k[\epsilon]/(\epsilon^2)\). Write a putative tensor as

\[
e=a+b(\epsilon\otimes1)+c(1\otimes\epsilon)
  +d(\epsilon\otimes\epsilon).
\]

The relation \((\epsilon\otimes1)e=(1\otimes\epsilon)e\) forces \(a=0\)
and \(b=c\). Its product is then \(2b\epsilon\), never \(1\). This
argument works in every characteristic. In characteristic \(2\),
\(t^2-1=(t-1)^2\), so the involution algebra itself becomes this
counterexample after putting \(\epsilon=t-1\).

\subsection{Splitting over one field is weaker than
separability}\label{splitting-over-one-field-is-weaker-than-separability}

For a monic polynomial \(m\), factor it in \(k[t]\) as
\(m=\prod_i p_i^{r_i}\) with distinct monic irreducibles \(p_i\). The
algebra \(S=k[t]/(m)\) is semisimple precisely when every \(r_i=1\).
Here is a proof that keeps the field issue visible.

If all \(r_i=1\), the polynomial CRT identifies \(S\) with the product
of the fields \(k[t]/(p_i)\). Every \(S\)-module decomposes by these
central idempotents, and every submodule has a complement by taking a
vector-space complement over each component field.

If some \(r_i\ge2\), consider the component \(B=k[t]/(p_i^{r_i})\). The
quotient \(B\to B/(p_i)\) has no \(B\)-linear section. The image of
\(1\) under such a section would be an element \(v\equiv1\pmod{p_i}\)
killed by \(p_i\). But \(p_i v=0\) in \(B\) forces \(v\) to be divisible
by \(p_i^{r_i-1}\), contradicting \(v\equiv1\pmod{p_i}\). Extending this
module structure through the CRT projection gives a nonsplit sequence of
\(S\)-modules.

This semisimplicity criterion uses irreducible factorization over \(k\).
It does not assert that the roots in an algebraic closure are distinct.
The derivative criterion (A13) imposes that stronger requirement.

For example, let

\[
k=\mathbb F_p(s),\qquad m(t)=t^p-s.
\]

Eisenstein at the prime \(s\) in \(\mathbb F_p[s]\), followed by Gauss's
lemma, proves that \(m\) is irreducible over \(k\). Thus \(S=k[t]/(m)\)
is a field: every sequence of \(S\)-vector spaces splits. Nevertheless
\(m'=0\), so (A12) is a nonzero derivation \(S\to S\) and no
separability tensor exists.

For any fixed pair of modules, ordinary vector-space projections onto
the subspace of intertwiners may still be chosen. What fails here is the
separability-tensor construction, which repairs all pairs coherently as
in (A3).

Adjoin an element \(\alpha\) with \(\alpha^p=s\), and write
\(K=k(\alpha)\). Then

\[
K\otimes_k S\cong K[t]/((t-\alpha)^p)
              \cong K[u]/(u^p).
\]

This algebra has a nonsplit sequence with middle term \(K[u]/(u^2)\) and
trivial one-dimensional end terms. The base change has exposed nilpotent
structure that was invisible to factorization over \(k\).

In contrast, a tensor satisfying (A4) survives every extension of the
ground field. Its coefficients can be extended, and the same repair
proof still works. For the polynomial algebras here, separability is
equivalent to semisimplicity after every field extension. One direction
follows from base change of the tensor. For the other, a repeated root
over a splitting field gives a quotient by \((t-\lambda)^2\) and the
nonsplitting example above.

Over a perfect field, every irreducible polynomial is separable, so the
two criteria coincide. Over an imperfect field, ``no repeated
irreducible factors in \(k[t]\)'' and ``coprime to the derivative'' must
be kept separate.

\subsection{Which extensions is the repair allowed to
split?}\label{which-extensions-is-the-repair-allowed-to-split}

The previous proof applies to exact sequences \textbf{of \(S\)-modules}.
It requires \(m\) to annihilate the middle term, not merely the two end
terms.

Take \(R=k[t]\) and \(S=R/(t)\cong k\). The exact sequence

\[
0\longrightarrow k
 \xrightarrow{\ 1\mapsto\bar t\ }R/(t^2)
 \xrightarrow{\ \bar f\mapsto f(0)\ }k
 \longrightarrow0                                \tag{A14}
\]

has both end terms annihilated by \(t\). It does not split as a sequence
of \(R\)-modules: any lift of \(1\) in the quotient is \(1+ct\), and
\(t(1+ct)=t\ne0\). Its middle term is not an \(S\)-module, because \(t\)
does not act as zero there.

Thus the correct comparison is

\[
\operatorname{Ext}^1_S(k,k)=0,\qquad
\operatorname{Ext}^1_R(k,k)\cong k.
\]

The latter group is computed in the Sylvester/extension chapter. There
is no contradiction. Changing the scalar ring changes which middle terms
are eligible extensions.

This also locates the difference between two splitting mechanisms. For
operators with coprime minimal polynomials on the two end terms, the
Sylvester equation gives a unique invariant section of every extension
with those end actions. The repair in this chapter may allow the end
terms to share eigenvalues, but restricts the entire middle operator to
a separable annihilating polynomial. In that category an invariant
section exists, and its choices need not be unique.

\subsection{Maschke's sum is the same tensor
construction}\label{maschkes-sum-is-the-same-tensor-construction}

Let \(G\) be a finite group and suppose \(|G|\) is invertible in \(k\).
In \(A=k[G]\), put

\[
e_G=\frac1{|G|}\sum_{g\in G}
           g\otimes(g^{-1})^{\mathrm{op}}.          \tag{A15}
\]

Its product is \(1\). For \(h\in G\), reindexing by \(r=hg\) gives

\[
\sum_g hg\otimes g^{-1}
 =\sum_r r\otimes r^{-1}h.
\]

Linear extension establishes (A1) for every element of \(k[G]\). The
general repair lemma now reads

\[
\mathcal A(f)=\frac1{|G|}
  \sum_{g\in G}\rho_N(g)f\rho_M(g^{-1}).            \tag{A16}
\]

It repairs intertwiners, sections, and retractions simultaneously.
Maschke's invariant-complement theorem is the retraction case. No
irreducible-character theory, radical theory, or classification of
semisimple rings is needed for this proof.

There is also a short converse for the assertion that \emph{every}
\(k[G]\)-module sequence splits. Give \(k\) the trivial action and
consider the augmentation

\[
\varepsilon:k[G]\longrightarrow k,\qquad
\sum_g c_g g\longmapsto\sum_g c_g.
\]

An invariant vector in the left regular representation has all its
coefficients equal, since left multiplication acts transitively on the
group basis. Hence it is \(c\sum_g g\). If \(|G|=0\) in \(k\), every
such vector has augmentation zero. An equivariant section would have to
send \(1\) to an invariant vector of augmentation \(1\), which is
impossible. Thus the division by \(|G|\) marks a genuine obstruction for
group algebras.

\subsection{A repair without division by a
dimension}\label{a-repair-without-division-by-a-dimension}

The group formula should not suggest that every separable algebra needs
the inverse of its dimension. Let \(A=M_r(k)\) and denote its matrix
units by \(E_{ij}\). For any characteristic,

\[
e=\sum_{i=1}^r E_{i1}\otimes E_{1i}^{\mathrm{op}}    \tag{A17}
\]

is a separability idempotent. Its product is \(\sum_i E_{ii}=I\). For
\(c=E_{ab}\), both sides of the centralizing identity are
\(E_{a1}\otimes E_{1b}\); linearity proves the identity for every
matrix. The repair formula uses no denominator.

The associated module decomposition is equally concrete. Put
\(L=E_{11}M\). Then

\[
k^r\otimes_k L\longrightarrow M,\qquad
e_i\otimes v\longmapsto E_{i1}v
\]

is an \(A\)-module isomorphism, with inverse

\[
m\longmapsto\sum_i e_i\otimes E_{1i}m.
\]

The matrix algebra acts on the \(k^r\) factor. These formulas prove that
its modules are built from copies of the standard module, and show
directly why the conclusion still holds when the characteristic divides
\(r\). This is the elementary matrix-unit case of the relation between
matrix algebras and their coefficient fields.

\subsection{The denominator also remembers integral
gluing}\label{the-denominator-also-remembers-integral-gluing}

The rational involution formula \(e=(1+xy)/2\) has an arithmetic
boundary. Work over \(\mathbb Z\) with
\(A_{\mathbb Z}=\mathbb Z[t]/(t^2-1)\). Rationally its two eigenspaces
are separated by \((1\pm t)/2\). Integrally the evaluation map is

\[
A_{\mathbb Z}\longrightarrow\mathbb Z\oplus\mathbb Z,
\qquad a+bt\longmapsto(a+b,a-b).
\]

Its image consists exactly of pairs congruent modulo \(2\); it has index
\(2\). This is the resultant obstruction for \(t-1\) and \(t+1\) in the
integral CRT chapter.

There is no separability tensor over \(\mathbb Z\) for this algebra.
Such a tensor would reduce modulo \(2\) to one for
\(\mathbb F_2[t]/((t-1)^2)\), contradicting the dual-number calculation.
Thus the rational denominator, the integral index, and the repeated root
modulo \(2\) describe three aspects of the same failure in this example.
This does not mean that every rational repair of a particular map has
nonintegral coefficients; particular maps may already respect the
action. It means that the universal tensor repair cannot be defined
integrally.

\subsection{Reconstructive exercises}\label{reconstructive-exercises}

\begin{enumerate}
\def\labelenumi{\arabic{enumi}.}
\tightlist
\item
  Starting only from (A1), prove the repair lemma and repair both a
  section and a retraction. Explain precisely where the surrounding maps
  must be \(A\)-linear.
\item
  Derive (A9), (A10), and (A11) from the divided difference, without
  using eigenvectors. For (A11), verify directly that the repaired map
  intertwines the two cubic operators.
\item
  Prove that a commutative algebra has at most one separability tensor.
  Identify the tensor for \(k^d\) and explain why it deletes maps
  between distinct direct factors.
\item
  Check that (A12) is independent of the representative of \(\bar f\).
  Use it to prove the necessity in (A13). Why does this argument not
  prove that every nonseparable algebra has a nonsplit sequence of its
  own left modules over the original field?
\item
  Prove that (A14) is nonsplit, then identify the failed hypothesis in
  an attempted application of the repair for \(S=k\).
\item
  Prove the augmentation converse to Maschke. Compare it with the
  denominator-free matrix-unit construction (A17).
\item
  Explain why the congruence \(a+b\equiv a-b\pmod2\) is the entire
  integral obstruction for \(\mathbb Z[t]/(t^2-1)\). Relate it to the
  reduction of the derivative \(2t\) modulo \(2\).
\end{enumerate}

\textbf{Answer landmarks.} For 1, use the tensor identity to move the
scalar from the first slot to the input slot, then apply (A3). For 2,
invert the derivative modulo the polynomial before multiplying the
divided difference; the inverses are \(y/2\), \(-y/2\), and \(y/6\),
respectively. For 3, \(ee'=e'=e\), and for \(k^d\) the tensor is
\(\sum_j e_j\otimes e_j\). For 4, the ambiguity is \(m'h+mh'\); the
inseparable field \(\mathbb F_p(s)[t]/(t^p-s)\) disproves the proposed
stronger conclusion. For 5, \(t\) kills both ends but not the middle.
For 6, the regular invariant vectors all have augmentation zero in
modular characteristic, whereas the product in (A17) is \(I\) in every
characteristic. For 7, a pair \((u,v)\) reconstructs \(a=(u+v)/2\),
\(b=(u-v)/2\).

\textbf{Sources and scope.} These are standard algebraic mechanisms,
presented here with explicit reconstruction and boundary examples; no
novelty claim is intended. Marcelo Aguiar's original paper
\href{https://pi.math.cornell.edu/~maguiar/strongly.pdf}{\emph{A note on
strongly separable algebras}, §§1--2} states the separability-tensor
identities and includes the group-algebra and matrix-algebra examples.
The paper's additional notion of \emph{strong} separability is not
required for the repair lemma. The
\href{https://stacks.math.columbia.edu/tag/00U3}{Stacks Project, Lemma
10.143.4} identifies finite étale algebras over a field as finite
products of finite separable field extensions; its proof supplies a
broader setting for the polynomial criterion. The arguments above prove
the formulas used here without requiring that general theory. The
source-comparison guide separately locates the directly inspected
textbook proofs of polynomial decomposition and Maschke's theorem.

\section{Integral spectral gluing: the lattice behind Fourier
coordinates}\label{integral-spectral-gluing-the-lattice-behind-fourier-coordinates}

\textbf{Status and scope.} This is a local continuation of Chapter 9 at
baseline
\href{https://github.com/Drsunflowerherbert/Codex-for-Math/tree/a4a28796bdd63c5cbddbfba698c95d34728bd94c}{a4a2879}.
The CRT, resultant, cyclotomic, and Fourier facts below are classical.
The contribution of this chapter is their organization into explicit
reconstruction certificates, together with small examples separating
conditions that are easily conflated. No literature novelty, search
improvement, universal existence theorem, or new Hadamard obstruction is
asserted. The proofs are supplied here. The
\href{PROOF-COMPARISON-ARITHMETIC.md}{arithmetic proof comparison}
identifies the inspected textbook arguments and the additional integral
gluing proved in this chapter.

Chapter 9 identified Fourier diagonalization with CRT over a splitting
field. This chapter starts at the point where that description ceases to
remember enough: the original coefficients must be integers, and often
signs. An isomorphism of rational algebras need not identify their
integral lattices. Even an integral inverse need not have sign
coefficients. The central conclusion is Theorem 6A: once the full
spectral norm equations hold, sign reconstruction at every cyclic order
is exactly membership in a specified affine integral lattice. The
prime-order formulas and prime-power gluing explain how that lattice is
tested.

Throughout, \(R=\mathbb Z[t]\), and \[
 \mathcal A_m=\mathbb Z[z]/(z^m-1),\qquad
 X=\sum_{j=0}^{m-1}x_jz^j.
\] The involution \(X^*=X(z^{-1})\) is defined modulo \(z^m-1\). We use
the Fourier convention \(\lambda_r=X(\zeta_m^r)\); hence inverse Fourier
coefficients use \(\zeta_m^{-rj}\). The augmentation is
\(\epsilon(X)=X(1)=\sum_jx_j\), and \(1\) denotes the constant
polynomial, not the all-ones polynomial \(T=1+z+\cdots+z^{m-1}\).

The dependency chain is short:

\begin{longtable}[]{@{}
  >{\raggedright\arraybackslash}p{(\columnwidth - 4\tabcolsep) * \real{0.3333}}
  >{\raggedright\arraybackslash}p{(\columnwidth - 4\tabcolsep) * \real{0.3333}}
  >{\raggedright\arraybackslash}p{(\columnwidth - 4\tabcolsep) * \real{0.3333}}@{}}
\toprule\noalign{}
\begin{minipage}[b]{\linewidth}\raggedright
Input
\end{minipage} & \begin{minipage}[b]{\linewidth}\raggedright
Reconstructive output
\end{minipage} & \begin{minipage}[b]{\linewidth}\raggedright
What remains to check
\end{minipage} \\
\midrule\noalign{}
\endhead
\bottomrule\noalign{}
\endlastfoot
Monic factors coprime over \(\mathbb Q\) & A rational CRT inverse and an
exact finite integral defect & Membership in the integral image \\
An integral cyclotomic tuple & Rational coefficients with an explicit
divisibility certificate & The coefficient lattice congruences \\
A tuple in that lattice & A unique integral row & Symmetry,
normalization, and signs \\
Integral odd coefficients and exact total energy & Every coefficient is
a sign & Any additional symmetry or normalization required by the
construction \\
Coherent four-row spectral Gram equations & The full group-algebra Gram
identity & Parity or another condition forcing signs \\
\end{longtable}

\subsection{The integral CRT defect is a finite
module}\label{the-integral-crt-defect-is-a-finite-module}

\textbf{Theorem 1.} Let \(f,g\in\mathbb Z[t]\) be monic of positive
degrees \(r,s\), and assume \(\gcd_{\mathbb Q[t]}(f,g)=1\). There is an
exact sequence of \(R\)-modules \[
 0\longrightarrow R/(fg)
 \xrightarrow{\alpha}R/(f)\oplus R/(g)
 \xrightarrow{\delta}R/(f,g)\longrightarrow0,                 \tag{1}
\] where \[
 \alpha(h)=(h\bmod f,h\bmod g),\qquad
 \delta(u,v)=u-v\bmod(f,g).
\] The last map is an additive module map; it is generally not a ring
map. The first map is a ring injection, and its image has index \[
 [R/(f)\oplus R/(g):\alpha(R/(fg))]
       =|\operatorname{Res}(f,g)|.                           \tag{2}
\] In particular, rational coprimality gives an integral CRT isomorphism
exactly when the resultant is \(1\) or \(-1\).

\textbf{Proof of exactness.} Monic division makes all three principal
quotients free abelian groups in their usual power bases. If
\(h\in(f)\cap(g)\), write \(h=fu\) with \(u\in R\). Rational coprimality
implies \(g\mid u\) in \(\mathbb Q[t]\). Divide \(u\) by the monic
polynomial \(g\) in \(R\); uniqueness of the rational remainder shows
that the integral remainder is zero. Thus \(h\in(fg)\), proving
\((f)\cap(g)=(fg)\) and the injectivity of \(\alpha\).

The difference map is well defined and surjective. If \(u-v=af+bg\),
then \[
 h=u-af=v+bg
\] maps to \((u\bmod f,v\bmod g)\). This proves exactness in the middle.
The reverse inclusion follows directly from the definitions. \(\square\)

The nontrivial part of (2) is a determinant computation over the
integers, not an application of field CRT.

\textbf{Determinant proof of the index formula.} Put \(L_f=R/(f)\), with
basis \(1,t,\ldots,t^{r-1}\), and let \[
 M_g:L_f\longrightarrow L_f,\qquad h\longmapsto gh.
\] Then \[
 \operatorname{coker}M_g\cong R/(f,g).                       \tag{3}
\] Over \(\mathbb Q\), the map \(M_g\) is invertible by Bézout's
identity, so its integral matrix has a nonzero determinant. An integral
square matrix \(M\) with nonzero determinant has cokernel of order
\(|\det M|\): integral row and column operations reduce it to Smith
normal form, whose nonzero diagonal entries multiply to that determinant
in absolute value.

For monic \(f\), the standard resultant convention gives \[
 \det M_g=\operatorname{Res}(f,g)
         =\prod_{\alpha:f(\alpha)=0}g(\alpha),               \tag{4}
\] with roots counted with multiplicity. To see the determinant
identity, extend scalars to an algebraic closure. Multiplication by
\(t\) on \(L_f\) has companion matrix with characteristic polynomial
\(f\). Triangularize that matrix; applying \(g\) gives diagonal entries
\(g(\alpha)\), including repeated roots.

This is also the Sylvester determinant computation. Consider the
integral map \[
 S:R_{<s}\oplus R_{<r}\longrightarrow R_{<r+s},
       \qquad (a,b)\longmapsto af+bg.
\] Use monic division by \(f\) to identify its target with
\(R_{<s}\oplus R_{<r}\), by \(qf+h\leftrightarrow(q,h)\). This is an
isomorphism of integral lattices. In these coordinates the matrix of
\(S\) has block form \[
 \begin{pmatrix} I_s&Q\\0&M_g\end{pmatrix}.
\] Thus its determinant is \(\det M_g\), up to the sign from the chosen
ordering of Sylvester columns and target basis. Consequently (3), (4),
and exactness of (1) give (2).

This argument also proves the claimed equivalence: the defect vanishes
exactly when its order is one, which is exactly \((f,g)=R\). \(\square\)

Two cautions make the theorem more useful.

\begin{itemize}
\tightlist
\item
  Rational coprimality is weaker than integral comaximality. For
  \(f=t-1,\ g=t+1\), the defect is \(\mathbb Z/2\), and the image in
  \(\mathbb Z^2\) is \(\{(a,c):a\equiv c\pmod2\}\).
\item
  The resultant gives the \textbf{order} of the defect, not its complete
  abelian-group structure. For \(f=t^2-1,\ g=t^2+1\), the defect is \[
    R/(2,t^2-1)\cong\mathbb F_2[t]/((t-1)^2),
  \] whose additive group is \((\mathbb Z/2)^2\), although the resultant
  is \(4\). Smith normal form gives the additional structure.
\end{itemize}

Monicity performs two jobs: it gives finite free power-basis lattices,
and it makes rational divisibility descend to integral divisibility.
Primitivity is enough for the latter job by Gauss's lemma; it does not
make the principal quotient a finite free abelian group. Without even
primitivity, injectivity can fail: \(f=2t\) and \(g=2(t-1)\) are coprime
over \(\mathbb Q\), but \(2t(t-1)\) belongs to \((f)\cap(g)\) and not to
\((fg)\).

\textbf{Corollary 2: several factors.} If the monic polynomials
\(f_1,\ldots,f_k\) are pairwise coprime over \(\mathbb Q\), the
evaluation map \[
 R/(f_1\cdots f_k)\longrightarrow\bigoplus_i R/(f_i)
\] is injective and has index \[
 \prod_{i<j}|\operatorname{Res}(f_i,f_j)|.                  \tag{5}
\]

\textbf{Proof.} Apply Theorem 1 first to \(F=f_1\cdots f_{k-1}\) and
\(f_k\), and then refine the \(R/(F)\) coordinate. Indices of nested
finite-index lattices multiply. The product-over-roots formula (4) gives
\(\operatorname{Res}(F,f_k)=\prod_{i<k}\operatorname{Res}(f_i,f_k)\).
Induction proves the assertion. \(\square\)

Formula (5) does \textbf{not} say that checking the original factor
pairs separately characterizes the image. Section 5 gives a small
counterexample.

\subsection{Prime cyclic order is a single fiber
product}\label{prime-cyclic-order-is-a-single-fiber-product}

Let \(p\) be prime, \(\zeta=\zeta_p\), and \[
 K=\mathbb Q(\zeta),\qquad O=\mathbb Z[\zeta],\qquad
 \pi=1-\zeta.
\] We use the concrete order \(O\); no theorem identifying it with the
full ring of integers is needed here. The polynomial
\(\Phi_p(t)=1+t+\cdots+t^{p-1}\) is irreducible over \(\mathbb Q\):
after replacing \(t\) by \(t+1\), Eisenstein's criterion at \(p\)
applies. Hence \([K:\mathbb Q]=p-1\).

Reduction \(\zeta\mapsto1\) identifies \[
 O/(\pi)\cong\mathbb Z/(\Phi_p(1))\cong\mathbb F_p.          \tag{6}
\] Write \(\rho:O\to\mathbb F_p\) for this reduction.

\textbf{Theorem 3: exact prime gluing and reconstruction.} The map \[
 \mathcal A_p\longrightarrow\mathbb Z\oplus O,\qquad
 X\longmapsto(a,b)=(X(1),X(\zeta))
\] has image \[
 \mathbb Z\times_{\mathbb F_p}O
       =\{(a,b):\rho(b)=a\bmod p\},                         \tag{7}
\] of index \(p\). Its unique rational inverse has coefficients \[
 \boxed{\displaystyle
 x_j=\frac{a+\operatorname{Tr}_{K/\mathbb Q}
                  (b\zeta^{-j})}{p}},\qquad 0\le j<p.      \tag{8}
\] For \(a\in\mathbb Z,\ b\in O\), all the coefficients in (8) are
integers if and only if the single congruence in (7) holds.

\textbf{Proof of the image statement.} Apply Theorem 1 to \(t-1\) and
\(\Phi_p(t)\). Their common quotient is \(\mathbb F_p\), and their
resultant is \(\Phi_p(1)=p\). There is also a direct construction. Write
\[
 b=B(\zeta),\qquad B\in\mathbb Z[t],\quad\deg B\le p-2,
 \qquad d=B(1).
 \] If \(a\equiv d\pmod p\), set \(q=(a-d)/p\). Then \[
 X(t)=B(t)+q\Phi_p(t)                                    \tag{9}
\] has integral coefficients, degree less than \(p\), and prescribed
evaluations \(a,b\). Conversely any integral polynomial takes the same
residue at \(1\) and \(\zeta\) in the common quotient (6).

For arbitrary \(a\in\mathbb Q,\ b\in K\), the same construction with a
rational representative \(B\) proves existence and uniqueness of the
rational inverse.

\textbf{Proof of the trace formula.} The \(p-1\) embeddings are
\(\sigma_r(\zeta)=\zeta^r\), \(1\le r<p\), so \[
 \operatorname{Tr}(\zeta^\ell)=
 \begin{cases}
 p-1,&p\mid\ell,\\
 -1,&p\nmid\ell.
 \end{cases}                                             \tag{10}
\] For the rational inverse \(b=\sum_kx_k\zeta^k\), this gives \[
 \operatorname{Tr}(b\zeta^{-j})
       =(p-1)x_j-\sum_{k\ne j}x_k=px_j-a.
 \] This proves (8). Moreover, for every \(y\in O\), \[
 \operatorname{Tr}(y)\equiv-\rho(y)\pmod p,                \tag{11}
\] as follows by expanding \(y\) in its integral power basis and using
(10). Since \(\rho(\zeta^{-j})=1\), all \(p\) numerators in (8) have the
same residue \(a-\rho(b)\) modulo \(p\). They are all divisible by \(p\)
exactly when (7) holds. \(\square\)

The prime \(2\) is included without modification. Here \(\zeta=-1\),
\(K=\mathbb Q\), and \(O=\mathbb Z\), so \[
 a\equiv b\pmod2,\qquad
 x_0=(a+b)/2,\quad x_1=(a-b)/2.                            \tag{12}
\] There is no extra factor of two in the field trace at this endpoint.

\textbf{A tuple that does not glue.} At \(p=3\), the pair
\((a,b)=(1,0)\) consists of algebraic integers and gives real symmetric
Fourier coordinates \((1,0,0)\). Its inverse is \((1/3,1/3,1/3)\),
because the congruence in (7) fails. Thus neither algebraic integrality
of each spectral coordinate nor reality of all coordinates is enough.

\subsection{Galois coherence, symmetry, normalization, and
signs}\label{galois-coherence-symmetry-normalization-and-signs}

One field element \(b\) encodes every nontrivial frequency: \[
 \lambda_0=a,\qquad\lambda_r=\sigma_r(b)\quad(1\le r<p).
                                                               \tag{13}
\] Those coordinates cannot be chosen independently. In particular, if
one chooses real values at conjugate frequency pairs without preserving
(13), the resulting coefficients may be real but irrational. At \(p=5\),
the proposed tuple \[
 (\lambda_0,\lambda_1,\lambda_2,\lambda_3,\lambda_4)
       =(0,0,1,1,0)
\] is real and invariant under \(r\mapsto-r\), and each listed value is
an algebraic integer. It cannot be the spectrum of a rational row:
\(\lambda_1=0\) would force every Galois conjugate of that field
coordinate to be zero.

Suppose from now on that \(a\in\mathbb Z,\ b\in O\) satisfy (7), so
\(X\) is already an integral row. The following are exact tests: \[
 \begin{array}{rcl}
 X=X^* &\Longleftrightarrow& b=\overline b,\\[2mm]
 x_0=1&\Longleftrightarrow&a+\operatorname{Tr}(b)=p,\\[2mm]
 x_j\in\{-1,1\}\text{ for all }j
 &\Longleftrightarrow&
 a+\operatorname{Tr}(b\zeta^{-j})\in\{-p,p\}
       \text{ for all }j.                                \tag{14}
 \end{array}
\] For the first equivalence, the evaluations of \(X^*\) are
\((a,\bar b)\), and rational CRT is injective. The other two statements
follow from (8). Since cyclotomic Galois automorphisms commute with
complex conjugation, \(b=\bar b\) also says that every \(\sigma_r(b)\)
is real.

The trace in (8) is the full trace from \(K\). If \(p\) is odd and
\(b=\bar b\), one may instead use \(K^+=\mathbb Q(\zeta+\zeta^{-1})\):
\[
 \operatorname{Tr}_{K/\mathbb Q}(b\zeta^{-j})
 =\operatorname{Tr}_{K^+/\mathbb Q}
       \bigl(b(\zeta^j+\zeta^{-j})\bigr).                  \tag{15}
\] Indeed the relative trace from \(K\) to \(K^+\) sums a term and its
complex conjugate. The expression \(b\zeta^{-j}\) itself need not be
real; replacing its trace indiscriminately by twice a real-subfield
trace is invalid. Formula (15) does not apply at \(p=2\).

There is a useful substitute for checking all the sign conditions in
(14), provided parity and energy are checked exactly.

\textbf{Lemma 4: spectral parity for odd prime order.} For odd \(p\), an
integrally glued pair \((a,b)\) has every coefficient \(x_j\) odd if and
only if \[
                 a\equiv1\pmod2,\qquad b\in2O.            \tag{16}
\]

\textbf{Proof.} If all coefficients are odd, write \(X=T-2N\) with
\(N\in\mathcal A_p\). Since \(T(1)=p\) and \(T(\zeta)=0\), (16) follows.
Conversely assume (16), and put \[
 c=(p-a)/2\in\mathbb Z,\qquad d=-b/2\in O.
 \] Because \(2\) is invertible modulo \(p\) and \(\rho(b)=a\),
\(\rho(d)=-a/2=(p-a)/2=c\) in \(\mathbb F_p\). Theorem 3 therefore glues
\((c,d)\) to an integral \(N\). The evaluations of \(T-2N\) equal those
of \(X\); injectivity gives \(X=T-2N\), so all its coefficients are odd.
\(\square\)

At \(p=2\), the exact replacement is \[
 a+b\equiv a-b\equiv2\pmod4.                              \tag{17}
\] Merely requiring \(a,b\) even fails: \((a,b)=(2,2)\) glues to
\((2,0)\).

\textbf{Lemma 5: parity and the minimum energy force signs.} If
\(X_1,\ldots,X_q\in\mathcal A_m\) have odd coefficients and \[
                 \sum_{\ell,j}x_{\ell j}^2=qm,            \tag{18}
\] then every coefficient of every row is \(1\) or \(-1\).

\textbf{Proof.} There are \(qm\) odd integers in the sum; each has
square at least one. Equality in their summed lower bound forces every
square to be one. \(\square\)

For one prime-order row, exact Parseval reads \[
 a^2+\operatorname{Tr}(b\bar b)=p\sum_jx_j^2.              \tag{19}
\] To prove it, expand \(\sum_{r=0}^{p-1}|X(\zeta^r)|^2\); the geometric
sum kills each term \(x_jx_k\) with \(j\ne k\), and leaves
\(p\sum_jx_j^2\). The nontrivial frequencies give the trace term. This
is the positive quadratic trace form, not the multiplicative field norm
\(\operatorname{N}_{K/\mathbb Q}(b)\).

Consequently, at an odd prime, an integrally glued pair represents a
sign row exactly when \[
 a\text{ is odd},\qquad b\in2O,\qquad
 a^2+\operatorname{Tr}(b\bar b)=p^2.                      \tag{20}
\] Necessity follows from signs, and sufficiency follows from Lemmas
4--5. To represent a \textbf{normalized symmetric} sign row, add
\(b=\bar b\) and \(a+\operatorname{Tr}(b)=p\). These are finite exact
certificates, not an assertion that a tuple with prescribed extra Gram
constraints exists.

\textbf{Energy and every individual spectral bound still do not
suffice.} At \(p=11\), take \[
 x_0=1,\qquad x_{\pm1}=2,\qquad x_{\pm2}=-1,
 \qquad x_j=0\text{ otherwise}.                           \tag{21}
\] This row is integral, normalized, and symmetric. Its augmentation is
\(a=3\), so even the necessary normalized sign-row congruence
\(a\equiv p\pmod4\) holds. Its energy is exactly
\(1+2\cdot4+2\cdot1=11\). Every Fourier value has the form \[
 1+4\cos\theta-2\cos2\theta
      =3+4c-4c^2,\qquad -1\le c\le1.
 \] The concave quadratic has maximum \(4\) at \(c=1/2\) and minimum
\(-5\) at \(c=-1\). Thus every Galois conjugate satisfies \[
 |\lambda_r|^2\le25<44,
 \] which is stronger than the individual bound required of a row in a
four-row order-eleven Gram identity. The pair \[
 b=1+2(\zeta+\zeta^{-1})-(\zeta^2+\zeta^{-2}),\quad a=3
 \] glues integrally, and (19) has the exact sign-row energy \(121\).
Nevertheless the row is not a sign row. Its failure is visible in
parity, not in the energy or these individual spectral inequalities. No
completion of this row is being asserted.

\subsection{Full inverse Fourier divisibility for composite
order}\label{full-inverse-fourier-divisibility-for-composite-order}

For general \(m\), take one coordinate \[
 b_d\in O_d=\mathbb Z[\zeta_d]\qquad(d\mid m),
 \quad O_1=\mathbb Z.
 \] Here the standard cyclotomic factorization and irreducibility give
\[
 \mathbb Q[z]/(z^m-1)\cong
       \prod_{d\mid m}\mathbb Q(\zeta_d).
 \] The new issue is the integral image of this familiar rational
decomposition.

\textbf{Proposition 6: an integral inverse Fourier certificate.} Define
\[
 C_j=\sum_{d\mid m}
       \operatorname{Tr}_{\mathbb Q(\zeta_d)/\mathbb Q}
                 (b_d\zeta_d^{-j}),\qquad 0\le j<m.        \tag{22}
\] Every \(C_j\) is an integer. The unique rational inverse tuple has
coefficients \(x_j=C_j/m\), and the tuple comes from \(\mathcal A_m\) if
and only if \[
                         C_j\equiv0\pmod m
                         \quad(0\le j<m).                 \tag{23}
\]

\textbf{Proof.} Choose \(\zeta_d=\zeta_m^{m/d}\). For a frequency \(r\)
whose character has order \(d\), write \(r=(m/d)u\) with \(u\) a unit
modulo \(d\), and prescribe the coherent value
\(\lambda_r=\sigma_u(b_d)\). The zero frequency is \(b_1\). Partitioning
the ordinary inverse Fourier sum by character order gives exactly (22).
Orthogonality of the \(m\) characters proves that the resulting \(x_j\)
have the prescribed evaluations. They are rational because (22) is a sum
of traces.

Each trace is an integer: multiplication by \(b_d\zeta_d^{-j}\)
preserves the free integral power-basis lattice \(O_d\), so its matrix
has integer trace, equal to the field trace. The coefficient criterion
(23) now follows. \(\square\)

The certificate can contain redundant congruences; its advantage is that
it is exact and applies uniformly. It does not replace the task of
finding a tuple satisfying sign, symmetry, or Gram constraints. If signs
are required, one can check \(C_j\in\{-m,m\}\) directly. Alternatively,
after (23), check that every \(C_j/m\) is odd and use the exact energy
certificate of Lemma 5.

For odd \(m\), parity itself can be tested componentwise: an already
integrally glued tuple represents \(X\equiv T\pmod2\) if and only if its
components agree with those of \(T\) modulo \(2O_d\). To justify the
converse, reduce \[
 z^m-1=\prod_{d\mid m}\Phi_d(z)
 \] modulo \(2\). Its derivative is \(z^{m-1}\), relatively prime to the
polynomial, so the product is squarefree. Its displayed factors are
therefore pairwise coprime in \(\mathbb F_2[z]\), where ordinary CRT is
injective. This argument fails when \(m\) is even, precisely as the
\(p=2\) example already shows.

The following lattice formulation includes even order and avoids making
parity look like a collection of unrelated special cases. Keep the
cyclotomic factors \textbf{labelled by \(d\)} even when two underlying
number fields happen to be isomorphic. Set \[
 M_m=\bigoplus_{d\mid m}O_d,\qquad
 L_m=\operatorname{image}(\mathcal A_m\longrightarrow M_m),
 \qquad J_m=(m,0,\ldots,0).
 \] Thus \(J_m\) is the image of \(T\).

\textbf{Theorem 6A: the affine lattice and energy certificate, all
orders.} Let \(m\ge1\), and let four tuples \(b_\ell=(b_{\ell,d})\)
belong to \(M_m\). Assume the exact identities \[
 \sum_{\ell=1}^4b_{\ell,d}\overline{b_{\ell,d}}=4m
       \quad\text{in }\mathbb Q(\zeta_d),\qquad d\mid m.   \tag{23a}
\] Then the tuples reconstruct four sign rows with summed cyclic Gram
\(4m\,1\) if and only if \[
                         b_\ell\in J_m+2L_m
                         \quad(1\le\ell\le4).             \tag{23b}
\] Symmetry holds exactly when every component is fixed by complex
conjugation. Normalization \(x_{\ell0}=1\) holds exactly when \[
             \sum_{d\mid m}\operatorname{Tr}(b_{\ell,d})=m.
                                                               \tag{23c}
\] These last two conditions are additional requirements.

\textbf{Proof.} If \(b_\ell=J_m+2c_\ell\) with \(c_\ell\in L_m\), the
corresponding row is \(T+2Y_\ell\) with integral \(Y_\ell\). Its
coefficients are therefore odd integers. Apply every cyclotomic
embedding to (23a); all Fourier norm equations follow. Fourier
injectivity gives the full Gram identity. Its constant coefficient is
total energy \(4m\), so Lemma 5 forces signs. Conversely every sign row
has the form \(T+2Y_\ell\) with integral \(Y_\ell\), establishing (23b).
The symmetry test follows from the involution and injectivity; (23c) is
Proposition 6 at \(j=0\). This also covers \(m=1\), where there is only
the rational coordinate. \(\square\)

This certificate requires no separate assumption of integral gluing:
membership in the affine lattice already includes it. For practical
separation of the two tests, observe that \[
                         mM_m\subseteq L_m.               \tag{23d}
\] Indeed Proposition 6 reconstructs \(mb\) with coefficients \(C_j\),
which are integers. Thus \(M_m/L_m\) is annihilated by \(m\). For odd
\(m\), it has no nonzero two-torsion, and consequently \[
 b\in J_m+2L_m
 \quad\Longleftrightarrow\quad
 b\in L_m\ \text{ and }\ b-J_m\in2M_m.                    \tag{23e}
\] For the reverse implication, put \(c=(b-J_m)/2\in M_m\). Both \(2c\)
and \(mc\) lie in \(L_m\); an integer combination of \(2\) and \(m\)
equal to one gives \(c\in L_m\). At odd order the component conditions
in (23e) are exactly: the augmentation coordinate is odd, and every
\(d>1\) coordinate belongs to \(2O_d\).

At even order one must retain the actual half-shift test
\((b-J_m)/2\in L_m\). For \(m=4\), take four copies of the integral row
\(2\cdot1=(2,0,0,0)\). Every row has energy \(4\), every spectral
coordinate is \(2\), and (23a) holds. Each tuple lies in \(L_4\) and
satisfies \(b-J_4\in2M_4\). Nevertheless its half-shift reconstructs
\((1,-1,-1,-1)/2\), so (23b) fails and the rows are not signs. This
example does not impose normalization: the constant coefficient is
\(2\), showing again why (23c) is a separate condition.

Integral gluing cannot be dropped at odd order either. At \(m=3\) take
augmentation coordinates \((3,1,1,1)\) and nontrivial coordinates
\((0,2,2,2)\). These are odd/even as required, satisfy both exact
spectral Gram identities, and give individual rational row energies
\(3\). The first row reconstructs \((1,1,1)\); each of the other three
reconstructs \((5,-1,-1)/3\). Their failure is precisely the prime
gluing congruence \(1\not\equiv2\pmod3\).

If a procedure generates coefficient rows as signs from the outset,
their spectra automatically satisfy the integral and parity conditions.
In that setting the certificate explains existing structure; it is not a
new filter or a claimed speedup. No implementation claim about the
earlier search scripts is needed for this conditional observation.

\subsection{Why separate pairwise compatibility loses a condition at
order
four}\label{why-separate-pairwise-compatibility-loses-a-condition-at-order-four}

Factor \[
 z^4-1=(z-1)(z+1)(z^2+1).
 \] Write the three rational coordinates as \[
 a=X(1),\qquad c=X(-1),\qquad b=X(i)=u+iv.
 \] The inverse is \[
 \begin{aligned}
 x_0&=(a+c+2u)/4,&x_1&=(a-c+2v)/4,\\
 x_2&=(a+c-2u)/4,&x_3&=(a-c-2v)/4.
 \end{aligned}                                           \tag{24}
\] It follows that the exact integral image is \[
 \boxed{a+c\equiv2u\pmod4,\qquad
        a-c\equiv2v\pmod4.}                              \tag{25}
\] Indeed these congruences make the two displayed plus numerators
divisible by four; subtracting \(4u\) or \(4v\) proves the same for the
minus numerators. Necessity is immediate.

The evaluation matrix in coordinates \((a,c,u,v)\) is \[
 E_4=\begin{pmatrix}
 1&1&1&1\\1&-1&1&-1\\1&0&-1&0\\0&1&0&-1
 \end{pmatrix},\qquad \det E_4=-8.
 \] Thus the integral image has index \(8\). This agrees with (5), since
the three pairwise resultants each have absolute value \(2\).

Nevertheless, the common quotient for \textbf{each original factor pair}
is only \(\mathbb F_2\). The pairwise compatibility conditions are \[
                  a\equiv c\equiv u+v\pmod2.              \tag{26}
\] There are two independent mod-two equations, so this
pairwise-compatible lattice has index \(4\), not \(8\). For example, \[
          (a,c,b)=(0,0,1+i)
 \] passes (26), but its inverse is \((1,1,-1,-1)/2\). Even the real
tuple \((2,0,0)\) passes (26) and reconstructs \((1,1,1,1)/2\).

The missing information has an explicit explanation. Glue \(a,c\) first,
requiring \(a\equiv c\pmod2\). Their integral polynomial modulo
\(z^2-1\) is \(A+Bz\), where \[
 A=(a+c)/2,\qquad B=(a-c)/2.
 \] The next gluing step is between \(z^2-1\) and \(z^2+1\). Its common
quotient is \[
 \mathbb Z[z]/(z^2-1,z^2+1)
       \cong\mathbb F_2[z]/((z-1)^2),                     \tag{27}
\] and the necessary and sufficient condition is \[
                 b\equiv A+Bi\pmod{2\mathbb Z[i]}.
 \] Its two coefficient conditions are exactly (25). Replacing the
common quotient (27) by its residue field \(\mathbb F_2\) discards a
nilpotent direction. Thus the failure of the naive pairwise test is
already visible in a ring of four elements.

Order nine gives another exact scale for the lattice issue: \[
 z^9-1=\Phi_1\Phi_3\Phi_9,\qquad
 |\operatorname{Res}(\Phi_1,\Phi_3)|=3,\quad
 |\operatorname{Res}(\Phi_1,\Phi_9)|=3,\quad
 |\operatorname{Res}(\Phi_3,\Phi_9)|=9.
 \] For the last equality, \(\Phi_9(\zeta_3)=3\), so its multiplication
map on the rank-two lattice \(O_3\) is \(3I\). The integral cyclotomic
image therefore has index \(81\). This calculation is distinct from the
order-nine compression fixture in Chapter 9: an index of \(81\) neither
proves nor disproves any sign-lifting claim about that fixture.

\subsubsection{Prime-power gluing records a truncated Taylor
expansion}\label{prime-power-gluing-records-a-truncated-taylor-expansion}

There is an elementary general version of (27). Let \(p\) be prime,
\(a\ge1\), \(q=p^{a-1}\), and set \[
 f=z^q-1,\qquad
 g=\Phi_{p^a}(z)=1+z^q+\cdots+z^{(p-1)q}.
 \] Since \(g\equiv p\pmod f\), \[
 (f,g)=(f,p),\qquad
 R/(f,g)\cong\mathbb F_p[z]/((z-1)^q)
              \cong\mathbb F_p[u]/(u^q).                 \tag{27a}
\] Thus a lower-order integral row \(A\in\mathcal A_q\) and an upper
coordinate \(b=B(\zeta_{p^a})\in O_{p^a}\) glue exactly when their
images in this ring agree. This is a direct application of Theorem 1, so
it is both necessary and sufficient. The defect of this stage has order
\(p^q\).

To give all its conditions without factorial denominators, define the
Hasse derivative \[
 D^{[k]}\!\left(\sum_jc_jz^j\right)(1)
       =\sum_{j\ge k}\binom jk c_j.
 \] The binomial identity gives the exact expansion \[
 C(1+u)=\sum_{k\ge0}D^{[k]}C(1)\,u^k.
 \] Therefore the gluing conditions are precisely \[
 D^{[k]}A(1)\equiv D^{[k]}B(1)\pmod p,
                 \qquad 0\le k<q.                       \tag{27b}
\] They are independent of the chosen representatives: modulo \(p\),
\(f=(z-1)^q\) and \(g=(z-1)^{(p-1)q}\), so adding a multiple of either
defining relation does not change the relevant truncated expansion.

The word ``jet'' here means only this finite list of coefficients of
\(C(1+u)\) modulo \(u^q\). It adds no constraints beyond equality in
(27a). On polynomials of degree less than \(q\), substitution \(z=1+u\)
is an invertible triangular integer basis change, with inverse
\(u=z-1\). Thus the jet congruences and the coefficient congruences in
the common quotient are equivalent descriptions.

At order \(9\), the final stage has \(q=3\), so it remembers the
constant, linear, and quadratic coefficients in \(\mathbb F_3[u]/u^3\).
Taking \(A=0\) and \(b=(\zeta_9-1)^2\) passes the constant and linear
tests but fails the quadratic test. At order \(8\), taking \(A=0\) and
\(b=(\zeta_8-1)^2\) shows why ordinary derivatives are unsafe in
characteristic two: every positive ordinary derivative of \((z-1)^2\)
vanishes modulo two, as does its value at \(1\), but its second Hasse
derivative is \(1\).

Iterating the stage indices from \(1\) to \(p^a\) gives \[
 [M_{p^a}:L_{p^a}]
       =p^{1+p+\cdots+p^{a-1}}.                          \tag{27c}
\] For order \(9\), the full evaluation lattice has Smith invariants \[
                    1,1,1,1,1,1,3,3,9.
 \] For an explicit reduction, write a degree-less-than-nine polynomial
uniquely as \(B+C\Phi_9\), with \(\deg B<6\), \(\deg C<3\). The
\(\Phi_9\) coordinate is \(B\); since \(\Phi_9\) equals \(3\) in both
lower factors, subtracting the lower evaluations of \(B\) leaves the
block \(3E_3\), where \(E_3\) is the order-three integral evaluation
matrix. Its Smith invariants are \(1,1,3\) (an entry and a two-by-two
minor are units, and its determinant has magnitude \(3\)). The full
order-nine matrix is therefore equivalent over integral unimodular basis
changes to \(I_6\oplus3E_3\), proving the displayed invariants.

\subsection{What the joint Gram equations actually
enforce}\label{what-the-joint-gram-equations-actually-enforce}

For integral or real coefficient rows, Fourier evaluation gives \[
 \sum_{\ell=1}^{q}X_\ell X_\ell^*=qm\,1
 \quad\Longleftrightarrow\quad
 \sum_{\ell=1}^{q}|\lambda_{\ell,r}|^2=qm
       \quad\text{at every frequency }r.                 \tag{28}
\] The reverse implication uses injectivity of the full complex Fourier
map. It requires a coherent spectrum for each row. Independent
feasibility intervals at different frequencies do not supply such rows.

For \(q=4\) at prime order, exact algebraic coordinates make (28)
particularly economical: \[
 \sum_{\ell=1}^4a_\ell^2=4p,\qquad
 \sum_{\ell=1}^4b_\ell\bar b_\ell=4p
                    \quad\text{in }K.                   \tag{29}
\] Applying every Galois embedding to the second identity gives every
nonzero-frequency equation. For symmetric rows one may replace
\(b_\ell\bar b_\ell\) by \(b_\ell^2\). One equality proved \textbf{as an
algebraic equality} in a faithful embedding is enough; testing a
numerical approximation or only an inequality at one embedding is not
the same thing.

\textbf{Corollary 7: an exact prime-order four-row sign certificate.}
Let \(p\) be odd, and choose four pairs
\((a_\ell,b_\ell)\in\mathbb Z\oplus O\). They reconstruct a normalized
symmetric Williamson quadruple if and only if all of the following hold:

\begin{enumerate}
\def\labelenumi{\arabic{enumi}.}
\tightlist
\item
  \(\rho(b_\ell)=a_\ell\bmod p\) for each row.
\item
  \(b_\ell=\bar b_\ell\) and \(a_\ell+\operatorname{Tr}(b_\ell)=p\) for
  each row.
\item
  Every \(a_\ell\) is odd and every \(b_\ell\) belongs to \(2O\).
\item
  The two exact Gram identities in (29) hold.
\end{enumerate}

\textbf{Proof.} Necessity follows by evaluation of a normalized
symmetric sign quadruple. Conversely, condition 1 reconstructs integral
rows, condition 2 gives symmetry and normalization, and condition 3
gives odd coefficients by Lemma 4. Condition 4 gives (28), whose
constant coefficient is total energy \(4p\). Lemma 5 makes every entry a
sign. The resulting group-algebra identity is exactly the Williamson
condition from Chapter 9. \(\square\)

This is an equivalent exact formulation, not a relaxation and not an
existence proof. Its value is that each obstruction has an identified
location: a lattice congruence, a linear condition, parity, or a coupled
quadratic identity.

\textbf{Even the full Gram identity does not force signs without a lower
bound on coefficient size.} In \(\mathcal A_3\), put \[
 X_1=1-2z-2z^2,\qquad X_2=X_3=X_4=1.                     \tag{30}
\] All four rows are integral, symmetric, and normalized by \(x_0=1\).
Writing \(S=z+z^2\), the relation \(S^2=2+S\) gives \[
 X_1^2=1-4S+4(2+S)=9,\qquad
 \sum_{\ell=1}^4X_\ell X_\ell^*=12.
 \] Their prime spectral coordinates are \[
 (a_1,a_2,a_3,a_4)=(-3,1,1,1),\qquad
 (b_1,b_2,b_3,b_4)=(3,1,1,1).
 \] They satisfy integral gluing, normalization, symmetry, and both
equations in (29), exactly. Yet the coefficient rows are \((1,-2,-2)\)
and three copies of \((1,0,0)\). Their separate energies are
\(9,1,1,1\); the example does not claim to enforce energy \(3\) on each
row. Condition 3 of Corollary 7 fails. This example and (21) expose
different missing constraints.

More generally, integrality, odd coefficients, and the full identity
(28) force signs at \textbf{any} order by Lemma 5. Oddness can be
weakened to the requirement that all integral coefficients be nonzero,
but nonzeroness has a less convenient spectral description.

\textbf{The nonsymmetric boundary.} Without symmetry, the exact norm
equations produce complementary circulant sign rows. They do not by
themselves justify the standard Williamson block array displayed in
Chapter 9. For that array its block \((1,2)\) in the Gram matrix is \[
                  -AB^{\mathsf T}+BA^{\mathsf T}
                  -CD^{\mathsf T}+DC^{\mathsf T}.
 \] Commutation of circulants does not make each term cancel its
partner: one sufficient extra condition is pairwise amicability
\(UV^{\mathsf T}=VU^{\mathsf T}\). Symmetric circulants satisfy this
automatically.

For a concrete check, at length three take first rows \[
 A=(1,1,1),\quad B=C=(1,1,-1),\quad D=(1,-1,1).
 \] The sum of their cyclic Gram matrices is \(12I_3\). With the Chapter
9 convention \(P_{ij}=1\) when \(j=i+1\), the standard block array
nevertheless has Gram block \((1,2)\) \[
 \begin{pmatrix}
 0&-4&4\\4&0&-4\\-4&4&0
 \end{pmatrix}\ne0.
 \] Indeed \(D=B^{\mathsf T}\), while \(A\) commutes amicably with every
circulant; the offending block is \((B^{\mathsf T})^2-B^2\). Thus
Theorem 6A gives complementarity without symmetry, whereas Corollary 7
deliberately retains symmetry to give a Williamson quadruple and the
stated block construction.

\subsection{Negative supports retain the discrete
information}\label{negative-supports-retain-the-discrete-information}

For actual sign rows define \[
 N_\ell=\sum_{x_{\ell j}=-1}z^j,\qquad
 r_\ell=\epsilon(N_\ell),\qquad R=\sum_{\ell=1}^4r_\ell.
 \] Then \(X_\ell=T-2N_\ell\). The essential assertion is that the
coefficients of \(N_\ell\) lie in \(\{0,1\}\); it is \textbf{not} the
group-algebra equation \(N_\ell^2=N_\ell\). Coefficientwise
multiplication and cyclic convolution are different.

Using \(T^2=mT\) and \(TN_\ell=N_\ell T=r_\ell T\), expansion gives \[
 \sum_{\ell=1}^4X_\ell X_\ell^*=4m\,1
 \quad\Longleftrightarrow\quad
 \sum_{\ell=1}^4N_\ell N_\ell^*+(m-R)T=m\,1.              \tag{31}
\] No symmetry is required for this equivalence. At a nonzero shift, the
equation says \[
 \sum_{\ell=1}^4
    |\operatorname{supp}(N_\ell)\cap
          (\operatorname{supp}(N_\ell)+s)|=R-m.            \tag{32}
\] At shift zero the equation holds automatically because the \(N_\ell\)
are indicators. Thus (31) is a joint system of cyclic difference counts,
not just four independent row conditions.

For odd \(m\), impose symmetry and normalization \(x_{\ell0}=1\). Each
negative support is a union of nonzero pairs \(\{j,-j\}\); therefore
every \(r_\ell\) is even and \(R\) is even. The row sums satisfy \[
 \epsilon(X_\ell)=m-2r_\ell\equiv m\pmod4,\qquad
                  2-m\le\epsilon(X_\ell)\le m.            \tag{33}
\] These are necessary conditions. They do not characterize either
individual sign rows in spectral space or a joint solution.

Symmetry lets us replace \(N_\ell N_\ell^*\) by \(N_\ell^2\). Reduce
(31) modulo two. Since \(m\) and \(m-R\) are odd, \[
             \sum_\ell N_\ell^2=T+1
                 \quad\text{in }\mathbb F_2[C_m].         \tag{34}
\] In characteristic two, squaring sends \(\sum n_jz^j\) to
\(\sum n_jz^{2j}\). Because \(m\) is odd, doubling permutes the residue
classes, so this map is an automorphism fixing \(T\) and \(1\). Applying
its inverse to (34) yields \[
                        \sum_\ell N_\ell=T+1.            \tag{35}
\] At every nonzero position an odd number of the four entries are
negative. Equivalently, \[
                x_{1j}x_{2j}x_{3j}x_{4j}=-1
                         \qquad(j\ne0).                  \tag{36}
\] This recovers the odd normalized symmetric Williamson product
condition with every ring and divisibility step visible.

The converse is absent: reducing (31) modulo two throws away its integer
difference counts. The order-nine fixture in Chapter 9, and its separate
continuation analysis, illustrates this loss. The present argument does
not need or re-run that enumeration.

At prime order the support coordinates are \[
 \epsilon(N)=(p-a)/2,\qquad N(\zeta)=-b/2.                 \tag{37}
\] Lemma 4 explains exactly when these data glue to an \textbf{integral}
\(N\). To obtain an actual negative support, its coefficients must also
lie in \(\{0,1\}\). A useful energy reformulation makes the last issue
explicit: if \(X=T-2N\) has integral \(N\), then \[
 \sum_jx_j^2=m
 \quad\Longleftrightarrow\quad
 \sum_jn_j(n_j-1)=0.
 \] Every summand on the right is nonnegative for integral \(n_j\);
equality is therefore equivalent to \(n_j\in\{0,1\}\) for all \(j\). The
same proof works after summing over several rows.

\subsection{Hypotheses, failed replacements, and bounded
conclusions}\label{hypotheses-failed-replacements-and-bounded-conclusions}

\begin{longtable}[]{@{}
  >{\raggedright\arraybackslash}p{(\columnwidth - 4\tabcolsep) * \real{0.3333}}
  >{\raggedright\arraybackslash}p{(\columnwidth - 4\tabcolsep) * \real{0.3333}}
  >{\raggedright\arraybackslash}p{(\columnwidth - 4\tabcolsep) * \real{0.3333}}@{}}
\toprule\noalign{}
\begin{minipage}[b]{\linewidth}\raggedright
Hypothesis or exact requirement
\end{minipage} & \begin{minipage}[b]{\linewidth}\raggedright
Tempting replacement
\end{minipage} & \begin{minipage}[b]{\linewidth}\raggedright
Counterexample or precise failure
\end{minipage} \\
\midrule\noalign{}
\endhead
\bottomrule\noalign{}
\endlastfoot
Integral comaximality & Coprimality over \(\mathbb Q\) & \(t-1,t+1\)
have defect \(\mathbb Z/2\) \\
Monic/primitive control of divisibility & Arbitrary rationally coprime
integral polynomials & \(2t,2(t-1)\): the product ideal is smaller than
the intersection \\
Full Galois coherence & Independently chosen real algebraic-integer
eigenvalues & The order-five tuple \((0,0,1,1,0)\) in Section 3 \\
Integral inverse Fourier congruences & Each cyclotomic coordinate is an
algebraic integer & \((a,b)=(1,0)\) at \(p=3\) gives coefficients
\(1/3\) \\
Simultaneous multi-factor gluing & All original pairwise tests &
\((a,c,b)=(0,0,1+i)\) at order four \\
Parity plus exact energy & Exact energy, symmetry, normalization, and
separate spectral upper bounds & The order-eleven row (21) \\
Parity plus full Gram & Full Gram, integrality, symmetry, and
normalization & The order-three quadruple (30) \\
The affine lattice \(J_m+2L_m\) at even order & Integral gluing and
componentwise congruence modulo \(2\) & Four copies of \(2\cdot1\) at
order four satisfy every norm equation but fail the half-shift gluing
test \\
Odd \(p\) in the componentwise parity lemma & The same mod-two
conditions at \(p=2\) & \((a,b)=(2,2)\) gives \((2,0)\) \\
Symmetry or sufficient amicability for the stated Williamson array &
Summed Gram equations for nonsymmetric circulant sign rows & The
explicit order-three array in Section 6 has a nonzero off-diagonal Gram
block \\
Odd \(m\) in the normalized product theorem & Arbitrary \(m\) & At
\(m=2\), rows \((1,1),(1,1),(1,-1),(1,-1)\) satisfy the full Gram
identity but have product \(+1\) at position \(1\) \\
Normalization in the stated product theorem & Symmetry alone & At
\(m=3\), start with \((1,1,1)\) and three \((1,-1,-1)\); negate the
first entire row. Gram and symmetry survive, while the nonzero-position
products become \(+1\) \\
\end{longtable}

None of these examples rules out a Hadamard matrix at any order. They
identify failures of specified implications and specified
representations. Conversely, Corollary 7 is sufficient precisely because
it restores the lost integral and discrete conditions.

\subsection{Exercises with answers}\label{exercises-with-answers}

\textbf{Exercise 1.} For \(f=t-1,\ g=t^2+t+1\), describe the image of
\(R/(fg)\) in \(\mathbb Z\oplus\mathbb Z[\zeta_3]\). Reconstruct
\((a,b)=(1,1+\zeta_3)\), and decide whether it is integral.

\textbf{Answer.} Write \(b=u+v\zeta_3\). The condition is
\(a\equiv u+v\pmod3\). Here \(u+v=2\), so the pair does not glue.
Formula (9) gives \[
 X=(1+t)-\tfrac13(1+t+t^2)
   =\tfrac23+\tfrac23t-\tfrac13t^2.
 \] This example also shows why taking a small integral power-basis
representative for \(b\) is not enough.

\textbf{Exercise 2.} At \(p=5\), reconstruct the pair \[
 a=1,\qquad b=-2(\zeta+\zeta^{-1}).
 \] Check symmetry, normalization, and signs by both direct coefficients
and the spectral tests.

\textbf{Answer.} The row is \((1,-1,1,1,-1)\), since \(X=T-2(z+z^4)\).
Its sum is \(1\), its trace \(\operatorname{Tr}(b)=4\), and
\(\rho(b)=-4=1\) in \(\mathbb F_5\). Also \(b=\bar b\), \(b\in2O\), and
\[
 \operatorname{Tr}(b^2)
 =4\operatorname{Tr}(\zeta^2+2+\zeta^{-2})
 =4(-1+8-1)=24.
 \] Thus \(a^2+\operatorname{Tr}(b^2)=25\), as required by (20).

\textbf{Exercise 3.} Explain why a four-row prime spectral certificate
cannot check only \(\sum_\ell|b_\ell|^2\le4p\) in one complex embedding,
even when the \(b_\ell\) are individually integral.

\textbf{Answer.} The proposed inequality imposes neither equality nor
the augmentation equation, integral gluing, parity, normalization, or
coherent constraints at the other embeddings. The exact identity
\(\sum b_\ell\bar b_\ell=4p\) in \(K\) transports to all embeddings; a
numerical upper bound in one embedding does not. The examples (21) and
(30) show that even much stronger subsets of the exact conditions can
fail to force signs.

\textbf{Exercise 4.} Prove that the order-four defect has additive group
\(\mathbb Z/4\oplus\mathbb Z/2\), not \((\mathbb Z/2)^3\).

\textbf{Answer.} The cokernel of \(E_4\) has order \(8\). Its first
Smith invariant is \(1\), since it has an entry \(1\); the second is
\(1\), since a \(2\times2\) minor from rows \(1,3\) and columns \(2,3\)
has determinant \(-1\). Every \(3\times3\) minor is even because \(E_4\)
has rank \(2\) modulo \(2\), and a \(3\times3\) minor (rows \(1,2,3\),
columns \(1,2,3\)) equals \(4\), while the minor using rows \(1,3,4\),
columns \(1,2,3\), equals \(2\). Thus the third determinantal divisor is
\(2\). The Smith invariants are \(1,1,2,4\).

\textbf{Exercise 5.} Why does the positivity argument in Lemma 5 require
integrality or an equivalent lower bound? Give a real symmetric
normalized row of odd length with energy equal to its length that is not
a sign row.

\textbf{Answer.} At length \(5\), take \((1,\sqrt2,0,0,\sqrt2)\). It is
symmetric and normalized, and its energy is \(1+2+2=5\), but its entries
are not signs. Without the integer odd/nonzero condition, an entry can
have square strictly between zero and one, and the minimum-energy proof
does not apply. The example itself also has zero entries and entries of
magnitude greater than one.

\textbf{Exercise 6.} Show directly that (35) is weaker than (31),
without using a search or the order-nine fixture.

\textbf{Answer.} At odd length \(m=5\), take \(N_1=T-1\) and
\(N_2=N_3=N_4=0\). These are normalized symmetric negative supports, and
their sum modulo two is \(T+1\), so (35) and (36) hold. The
corresponding row sums are \((-3,5,5,5)\), whose squared sum is \(84\),
not \(4m=20\). Thus the frequency-zero consequence of (31) already
fails.

\subsection{Verification and source
scope}\label{verification-and-source-scope}

The accompanying \href{scripts/check_integral_crt.py}{checker} uses
exact integer and rational arithmetic. All twelve check groups passed
with SymPy 1.14.0: 129 prime trace reconstructions (including \(p=2\));
150 bounded coefficient parity cases; all component-basis trace inverses
at orders \(1,2,3,4,5,6,8,9,10,11,12\); 81 order-four congruence cases;
66 prime-power gluing probes; Smith invariants at orders four and nine;
and the fixed energy, Gram, affine-lattice, and nonsymmetric-array
counterexamples and exercise calculations involving explicit rows. The
checks are consistency evidence for this exposition; the general proofs
above do not depend on testing finitely many cases.

The strongest formulation here is the all-order affine-lattice
certificate in Theorem 6A, with the normalized symmetric prime
specialization in Corollary 7. These reformulate classical identities
and elementary parity/energy reasoning; no novelty claim is made.
Independent mathematical reviews checked the all-order affine theorem,
the odd-order torsion argument, prime trace formula, even-order
boundary, order-four gluing, prime-power jets, order-nine Smith
reduction, and the nonsymmetric array counterexample. A separate
computation using Ramanujan sums verified both products of the
trace-inverse and evaluation matrices at eleven fixed orders. See the
\href{PHASE2-REVIEW.md}{continuation review} for the evidence boundary.

For the broader classical integral CRT theory, see Kamalakshya Mahatab
and Kannappan Sampath, \href{https://arxiv.org/abs/1401.7696}{``Chinese
Remainder Theorem for Cyclotomic Polynomials in Z{[}X{]}''},
\emph{Journal of Algebra} 435 (2015), 223--262. Their paper studies the
full cyclotomic cokernel and its Smith structure. The elementary special
cases and reconstruction certificates used here are proved explicitly
above; no claim is made that the inspected textbook passages already
contain them.

\section{Compression, orthogonal ideals, and a common-lift
obstruction}\label{compression-orthogonal-ideals-and-a-common-lift-obstruction}

The earlier Hadamard chapter proves that compression preserves
correlation equations. Here we determine exactly what it retains, count
the individual sign lifts, and explain the fixed core-nine failure by
one polynomial identity. These are three different tasks: lifting each
row, satisfying a quotient equation, and lifting all rows so that their
full equations hold. Their separation is useful even when every
underlying algebra is semisimple.

\subsection{The quotient has an explicit orthogonal
splitting}\label{the-quotient-has-an-explicit-orthogonal-splitting}

Write \(m=dr\) with \(d,r\) positive integers, and work over the reals:

\[A=\mathbb R[z]/(z^m-1),\qquad
B=\mathbb R[w]/(w^d-1),\qquad q(z)=w.\]

The involution is \(z^*=z^{-1}\). Put

\[e=\frac1r\sum_{\ell=0}^{r-1}z^{\ell d}.\]

Each element of the subgroup generated by \(z^d\) occurs \(r\) times in
the square of the numerator, so \(e^2=e\). Negation permutes that
subgroup, so \(e^*=e\), and \(q(e)=1\). Also \(ez^d=e\). It follows that

\[\iota:B\longrightarrow eA,\qquad \iota(w^s)=ez^s\]

is well defined and is inverse to \(q|_{eA}\). The identity of \(eA\) is
\(e\), not \(1_A\). Since \(e(1-e)=0\), multiplication gives the algebra
decomposition

\[A=eA\oplus(1-e)A,\qquad \ker q=(1-e)A.\]

It is also orthogonal for the coefficient inner product
\(\langle X,Y\rangle=[1](XY^*)\). Indeed,
\(\langle eX,(1-e)Y\rangle=[1](e(1-e)XY^*)=0\).

If \(X=\sum_jx_jz^j\) and \(y_s=\sum_{j\equiv s\pmod d}x_j\), direct
coefficient collection gives

\[(eX)_j=\frac{y_{j\bmod d}}r,\qquad
\|eX\|^2=\frac1r\sum_{s=0}^{d-1}y_s^2.\]

The retained component \(eX\) is the mean on each fiber; compression
records \(r\) times those means. The remaining coefficients record
deviations from them; they are not arbitrary once signs and symmetry are
imposed.

For four rows put \(L_i=eX_i\) and \(R_i=(1-e)X_i\). Cross products
vanish, so the full Gram identity \(\sum_iX_iX_i^*=4m1_A\) is equivalent
to two independent algebra equations:

\[\sum_iL_iL_i^*=4me,\qquad
\sum_iR_iR_i^*=4m(1-e).\]

The first is precisely \(\sum_iq(X_i)q(X_i)^*=4m\) in \(B\). There is no
extra factor of \(r\) in the right side. Taking the coefficient of the
identity gives total energies \(4d\) and \(4(m-d)\) respectively. For
example, the first follows from \(\sum_{i,s}y_{i,s}^2=4m\) and the
factor \(1/r\) above.

The algebraic decomposition holds for all positive \(d,r\). Oddness,
normalization, and signs enter the following discrete problem.

\subsection{The exact count of normalized symmetric row
lifts}\label{the-exact-count-of-normalized-symmetric-row-lifts}

Assume now that \(m,d,r\) are odd. A normalized symmetric sign row
satisfies

\[x_j\in\{-1,1\},\qquad x_{-j}=x_j,\qquad x_0=1.\]

An integer row \(y\in\mathbb Z^d\) has such a lift if and only if:

\begin{enumerate}
\def\labelenumi{\arabic{enumi}.}
\tightlist
\item
  \(y_s=y_{-s}\) for all \(s\);
\item
  for \(s\ne0\), \(|y_s|\le r\) and \(y_s\equiv r\pmod2\);
\item
  \(y_0=r-4a\) for an integer \(0\le a\le(r-1)/2\).
\end{enumerate}

When these conditions hold, the number of lifts is

\[
\binom{(r-1)/2}{(r-y_0)/4}
\prod_{s=1}^{(d-1)/2}\binom r{(r-y_s)/2}.
\]

\textbf{Proof.} Negation exchanges the two nonzero fibers \(F_s\) and
\(F_{-s}\). Choose \((r-y_s)/2\) negative positions among the \(r\)
positions of \(F_s\); reflection then fixes the entire opposite fiber.
Different pairs of fibers are independent. The zero fiber contains the
fixed coordinate \(0\), with value \(1\), and \((r-1)/2\) opposite
pairs. Making \(a\) pairs negative changes its all-positive sum \(r\) by
\(4a\). This proves both necessity and sufficiency, as well as the
product count.

Notice that the smallest possible zero-fiber sum is \(2-r\), not \(-r\).
For \(d=1\) the product is empty; for \(r=1\) the conditions say exactly
that \(y\) itself is the unique normalized symmetric sign row.

This proof gives explicit reconstruction. Choose the negative-position
subsets on representative fibers, reflect them, and independently choose
the negative pairs on the zero fiber. It does not solve the equations
that couple the four rows.

\subsection{What remains after subtracting the fiber
means}\label{what-remains-after-subtracting-the-fiber-means}

On a nonzero representative fiber, write \(x=\mathbf1-2u\), where
\(u\in\{0,1\}^r\) and \(\sum u_j=b=(r-y_s)/2\). Then

\[x-\frac{y_s}{r}\mathbf1
 =-2\left(u-\frac br\mathbf1\right).\]

These are the translated and rescaled vertices of the hypersimplex

\[\Delta(r,b)=\{u\in[0,1]^r:\ \sum u_j=b\}.\]

The opposite fiber is determined by reflection, not chosen
independently. The zero fiber has one fixed coordinate and duplicated
opposite pairs. Its choices are an affine image of the vertices of
\(\Delta((r-1)/2,a)\); treating it as an unrestricted
\(\Delta(r,(r-y_0)/2)\) would count forbidden rows. Explicitly, if \(u\)
records the \(a\) negative pairs, the residual coordinate at zero is
\(4a/r\), and both coordinates of pair \(j\) are \(4a/r-2u_j\). Use
\(\Delta(0,0)=\{()\}\) when \(r=1\).

Consequently the unresolved equation \(\sum_iR_iR_i^*=4m(1-e)\) is
imposed on a finite product of specified vertex sets. Real positivity
alone does not enforce those sets. Nor does the existence of an
individual lift for every compressed row guarantee a simultaneous
solution.

There is a particularly transparent real relaxation. Suppose real
compressed rows \(Y_i\in B\) already satisfy \(\sum_iY_iY_i^*=4m\). Set
\(L_i=\iota(Y_i)\), choose

\[R_1=2\sqrt m(1-e),\qquad R_2=R_3=R_4=0,\]

and put \(X_i=L_i+R_i\). Then \(q(X_i)=Y_i\) and \(\sum_iX_iX_i^*=4m\),
because \(R_1R_1^*=4m(1-e)\). If the \(Y_i\) are symmetric, the \(X_i\)
are symmetric too. This proves existence of an unrestricted real lift,
not of normalized sign rows. It explains why an obstruction can
disappear completely after replacing a discrete coefficient problem by a
real algebra problem.

\subsection{A proof of the fixed core-nine
obstruction}\label{a-proof-of-the-fixed-core-nine-obstruction}

Consider the four compressed rows, with labels fixed in this order,

\[
Y_A=(-1,-1,-1),\quad Y_B=(-1,1,1),\quad
Y_C=(-1,3,3),\quad Y_D=(3,-1,-1).
\]

They satisfy the compressed Gram equation with target \(36\). The
lift-count formula gives \(3,3,1,3\) individual lifts, hence \(27\)
labeled four-row lifts. Since the compressed rows are distinct, there
are also \(27\) unordered four-row multisets; if only their compressed
multiset is prescribed and row order matters, there are
\(24\cdot27=648\) ordered tuples. These are different counting
conventions.

A normalized symmetric row of length nine is determined by positions
\(1,2,3,4\), and its compression is

\[y_0=1+2x_3,\qquad y_1=y_2=x_1+x_2+x_4.\]

Let \(P=\{1,2,4\}\). On \(P\), rows \(A,D\) each have one positive
entry, row \(B\) has one negative entry, and row \(C\) is entirely
positive. At position \(3\) the column is \((-1,-1,-1,1)^{\mathsf T}\).
Write \(p,q\) for the positive exceptions of \(A,D\), and \(k\) for the
negative exception of \(B\).

The odd Williamson product theorem, proved in the earlier chapter,
requires the product of every nonzero column to be \(-1\). If \(p=q\),
the product of \(A,D\) is positive everywhere on \(P\), which would
require \(B\) to be negative everywhere there. If \(p\ne q\), that
product is negative at \(p,q\) and positive at the remaining position.
Thus the product condition holds exactly when \(p,q,k\) are the three
different elements of \(P\). It rejects \(21\) lifts and leaves six.

All six fail for the same algebraic reason. In \(\mathbb Z[z]/(z^9-1)\)
put \(S_j=z^j+z^{-j}\). Let \(\mathbf1\in\mathbb Z^4\) and
\(v_i=\mathbf1-2e_i\), with labels \(i=A,B,C,D\). Then

\[v_i\cdot v_j=4[i=j],\qquad \mathbf1\cdot v_i=2.\]

For the surviving pattern, the columns at \(p,q,k,3\) are respectively
\(v_D,v_A,-v_C,-v_D\). Hence the vector of four row polynomials is

\[\mathbf X=\mathbf1+v_DS_p+v_AS_q-v_CS_k-v_DS_3.\]

Since each \(X_i^*=X_i\) by symmetry, taking its dot product with itself
gives the relevant Gram polynomial:

\[
\sum_iX_i^2
=4+4(S_p+S_q-S_k-S_3)
 +4\sum_{j=1}^4S_j^2-8S_pS_3.
\]

Doubling permutes the four nonzero opposite-index classes modulo nine,
so

\[\sum_{j=1}^4S_j^2=8+S_p+S_q+S_k+S_3.\]

For \(p\in\{1,2,4\}\), the two classes represented by \(3-p,3+p\) are
exactly the other two members \(q,k\) of \(P\). Therefore
\(S_pS_3=S_q+S_k\), and substitution gives the precise defect

\[\boxed{\ \sum_iX_i^2-36=8(S_p-S_k)\ne0.\ }\]

The supports of \(S_p\) and \(S_k\) are disjoint. The summed
autocorrelation is \(8\) at shifts \(\pm p\), \(-8\) at shifts
\(\pm k\), and zero at all other nonzero shifts. Compression kills this
defect: \(q(S_p)=q(S_k)=w+w^{-1}\). Thus the remaining obstruction lies
exactly in the complementary ideal discarded by the quotient.

This proves that none of the \(27\) lifts in the specified normalized
symmetric family works. The \href{scripts/check_lifting.py}{exact
checker} independently reconstructs the rows and their correlations,
including all six polynomial defects. It is a bounded verification of
this proof and fixture, not a proof about all Hadamard matrices of order
\(36\) or a new enumeration campaign. The odd-product theorem is
classical; see Curtis Bright,
\href{https://arxiv.org/abs/1711.07056}{``A New Form of Williamson's
Product Theorem,'' Theorem 1}. No literature-priority claim is made for
the explanatory calculation.

\subsection{What the obstruction suggests, and what it does not
establish}\label{what-the-obstruction-suggests-and-what-it-does-not-establish}

Compression gives exact necessary equations and an explicit real
splitting. Integral Fourier reconstruction gives exact coefficient
lattice conditions. The affine-lattice sign theorem in the preceding
chapter combines those conditions with parity and energy. Each is a
complete statement about its specified data. None proves that feasible
partial data have a common completion.

For an algorithm, the remaining question is operational: can the
constraints be represented and propagated more cheaply than retaining
all original sign variables? Rewriting a condition as an exact lattice
membership test does not answer that question. A useful next experiment
would compare two implementations on the same fixed fixtures while
recording time, memory, and rejected states. That experiment is outside
this foundation study; no improvement in search complexity is claimed.

\subsection{Reconstruction exercises with
solutions}\label{reconstruction-exercises-with-solutions}

\textbf{1. Recover the scale.} For \(m=15,d=5\), a compressed four-row
solution has total coefficient energy \(60\). Find the total energies of
its retained and complementary components.

\textbf{Solution.} Here \(r=3\). Fiber averaging divides the retained
energy by three, giving \(20=4d\); the complementary target is
\(60-20=40=4(m-d)\). The quotient Gram target remains \(60\).

\textbf{2. Detect a forbidden zero fiber.} Can a normalized symmetric
sign row with fiber size five have compressed zero coordinate \(-5\),
\(-3\), or \(1\)?

\textbf{Solution.} The possible zero sums are \(5-4a\) for \(a=0,1,2\),
namely \(5,1,-3\). Thus \(-5\) is forbidden even though it has the
correct ordinary sign-sum parity and bound.

\textbf{3. Explain a false converse.} Every compressed row in the
core-nine fixture has a sign lift and the compressed Gram equation
holds. Which step fails in inferring a Williamson quadruple?

\textbf{Solution.} Individual choices must also solve a common equation
in \((1-e)A\). The product theorem leaves six choices, each with nonzero
defect \(8(S_p-S_k)\) in that ideal. Neither individual feasibility nor
the quotient equation couples those choices adequately.

\textbf{4. Separate algebra and arithmetic.} Why does the real lift
above not contradict the no-sign-lift proof?

\textbf{Solution.} It imposes neither the finite fiber vertex sets nor
the fixed coefficient \(x_0=1\). The real group algebra splits by an
idempotent containing \(1/r\); the normalized integral sign domain is
not the product of unrestricted real subsets in its two ideals.

\section{Selective comparison: what can be reorganized, what must be
retained}\label{selective-comparison-what-can-be-reorganized-what-must-be-retained}

The continuation adds two proof-level supplements:
\href{PROOF-COMPARISON-OPERATORS.md}{operators and multilinear
constructions} and \href{PROOF-COMPARISON-ARITHMETIC.md}{arithmetic and
representations}. They broaden the actual proof reading behind this
initial overview and give precise substitution boundaries, dependency
risks, and stopping tests. In particular, Xi supplies two Jordan
constructions and both Maschke proofs; the supplements distinguish those
methods from the graduate texts' routes.

\subsection{How to read the
comparison}\label{how-to-read-the-comparison}

The criterion is the user's intended use: rebuilding algebra around
reusable principles, then returning to linear algebra and the
Hadamard/Williamson study. ``Covered'' has several meanings. A theorem
may be stated, proved through a different construction, left as an
exercise, or illustrated geometrically. These are distinct
contributions. A chapter list can establish where to look; it cannot
establish that two books are interchangeable.

Evidence levels and edition/page conversions are in the source ledger.
The comparisons below are selective, and pedagogical judgments are
identified as judgments based on the inspected material. No global
difficulty ranking or exercise-by-exercise replacement audit is claimed.

\subsection{Xi I versus the graduate algebra
books}\label{xi-i-versus-the-graduate-algebra-books}

Xi I starts with equations and low-order determinants, develops sets and
maps, then connects vector rows/columns, rank, and linear maps before
abstract group/ring/field language. Its Chapter 4 does more than provide
a determinant formula: oriented-volume motivation and multilinearity
appear on printed pp.~75--76 (PDF pp.~81--82), and §4.4 gives multiple
characterizations on p.~93 (PDF p.~99). The surrounding exercises
include minor identities. That material has structural value even when
its elementary calculations are already familiar.

Hungerford III and Grillet III give a broader ring-theoretic
organization of polynomial division and factorization; their later
module chapters explain why a PID is the right setting for one-operator
structure. The present study therefore reorganizes the \textbf{use} of
Xi I's polynomial arithmetic into Bézout, Chinese remainder, and
annihilator arguments. It does not infer that Xi I's derivations,
examples, or root-location exercises are redundant.

Xi I §7.4 starts with symmetry in roots and elementary symmetric
polynomials (p.~177, PDF p.~183); §7.6 makes common-factor detection
into a resultant matrix (p.~190, PDF p.~196). These are useful links
from coefficient calculations to invariant expressions and elimination.
Grillet IV.7's resultant/discriminant location is verified from contents
(p.~176, PDF p.~188), but its entire proof and exercise treatment was
not compared against Xi's. No preference about replacement follows from
that heading alone.

For a returning reader who can prove quotient factorization,
rank--nullity, determinant multilinearity, Euclidean division, and
Bézout, much of Xi I can be used as a reference rather than reread
sequentially. That is a competence criterion for selective study, not a
judgment that introductory material lacks content.

\subsection{Xi II and two routes to Jordan
form}\label{xi-ii-and-two-routes-to-jordan-form}

There is substantial, verified theorem overlap. Xi II Theorem 2.51
(printed pp.~72--73, PDF pp.~78--79) proves the coprime-annihilator
split using Bézout. Theorems 2.52--2.53 build the generalized
eigenspaces. Theorems 2.56--2.57 then establish Jordan form under the
actual hypothesis that the characteristic polynomial splits in the base
field. Existence is proved by quotient induction for nilpotent
operators; uniqueness comes from dimensions of images of powers
(pp.~75--77, PDF pp.~81--83). These proof pages were read, not inferred
from the section title.

Grillet VIII.6 proves the PID structure theorem (pp.~336--342, PDF
pp.~348--354) and VIII.7 makes \(V\) a \(k[t]\)-module, proves its
annihilator has a monic generator, and derives primary cyclic blocks and
Jordan bases (pp.~342--346, PDF pp.~354--358). Its sampled proof gives a
real alternative construction, with rational canonical forms and general
coefficient rings naturally close at hand. Hungerford's suggested
linear-algebra course explicitly takes torsion PID modules before
VII.1--5 (printed p.~xviii, PDF p.~15); his IV.6 and VII.4--5 locations
support the same broad organization. A detailed proof-by-proof audit of
Hungerford VII was not made.

For the user's purpose, the module route is an efficient
\textbf{structural spine}: one remembers a module and its annihilator,
then recovers the canonical forms. Xi's route retains pedagogical value:
it shows what the projections do inside the actual vector space and
constructs chains through quotients. A reader who knows only the
abstract theorem should still reconstruct one chain and the
kernel-growth count. The two proofs illuminate different failure cases.

\subsection{Xi II tensor constructions versus Grillet and Mac
Lane}\label{xi-ii-tensor-constructions-versus-grillet-and-mac-lane}

Xi II Definition 6.11 constructs tensor products from a basis criterion,
proves uniqueness, and \textbf{proves the universal property} in Theorem
6.15 (pp.~273--275, PDF pp.~279--281). It is inaccurate to say that this
book has only coordinates or no universal-property treatment. Its Hom
identification, scalar extension, functorial maps, Kronecker products,
contraction as trace, and contraction as operator composition are
explicit examples (pp.~274--283, PDF pp.~280--289).

Grillet XI.5 constructs tensor products by generators and relations for
a right module and a left module, proving the balanced universal
property (pp.~434--437, PDF pp.~446--449). XI.6--8 adds right exactness,
dual modules, and flatness. This is a genuine extension of scope:
vector-space bases do not exist for arbitrary modules, and tensoring can
destroy injections. The simplest \(\mathbb Z/2\) counterexample explains
why that generality matters.

Mac Lane III.1--2 (pp.~55--62, PDF pp.~68--75) identifies universal
arrows, universal elements, and representability, and proves Yoneda.
IV.1 (pp.~79--80, PDF pp.~92--93) formulates adjunction, including the
Hom--tensor example. Its contribution is a common recognition language,
with careful naturality. It does not replace tensor calculations or tell
the reader which finite-dimensionality hypotheses can be forgotten.

There is one confirmed scope error: the remark after Xi II Theorem 6.19
(p.~278, PDF p.~284) extends its identification of \textbf{all} mixed
multilinear forms to infinite-dimensional spaces. The
identity-operator/finite-rank counterexample in the Hom chapter
disproves that extension for the preceding definition (p.~263, PDF
p.~269). This local correction does not invalidate the
finite-dimensional construction, its universality theorem, or the
chapter's examples. The definition and remark were checked as images.

\subsection{Exterior algebra: a real change in construction and
scope}\label{exterior-algebra-a-real-change-in-construction-and-scope}

Xi II §6.4 begins by explicitly assuming characteristic zero and models
exterior powers inside antisymmetric tensors (p.~292, PDF p.~298). Its
exterior product uses the alternation operator. This model naturally
continues the chapter's tensor calculations and geometric language.

Grillet XIII.2--4 builds tensor, symmetric, and exterior algebras by
universal properties and quotient ideals over a commutative ring; the
exterior relation is \(v\otimes v=0\) (pp.~518--527, PDF pp.~530--539).
The basis and alternating-map arguments were sampled in the actual proof
pages. This quotient construction is better suited to retaining
characteristic \(2\) and avoiding factorial denominators. In
characteristic zero, the models are connected by normalized alternation;
normalizations must be translated.

For determinants and minors, Xi I's characterizations and Xi II's
geometric motivation remain helpful. The notes choose the quotient
exterior construction so the rank and Cauchy--Binet arguments work in
arbitrary field characteristic. That is a decision about the desired
scope, not an assertion that Xi's stated characteristic-zero treatment
is incomplete on its own terms.

\subsection{Metric and geometric material is not superseded by module
theory}\label{metric-and-geometric-material-is-not-superseded-by-module-theory}

Xi II Chapter 3 develops real and Hermitian inner products,
self-adjoint, isometric and normal operators, complexification,
orthogonal expansion, projections, least squares, and
orthogonal-polynomial examples. The normal spectral theorem, Theorem
3.62 (pp.~126--128, PDF pp.~132--134), was read with its proof. This
supplies an orthonormal conclusion under a positive metric; Jordan form
supplies a different conclusion under arbitrary basis changes.

The inspected metric exercises include dual-basis transformations and
norm comparisons, providing practice that is not supplied by merely
knowing the module structure theorem. Xi's Chapters 4--5 add affine and
Euclidean geometry, transformation groups, indefinite spaces, quadrics,
and projective geometry. Only their contents and selected contextual
passages were audited here, so no claim that a graduate algebra book
covers all their proofs or exercises is warranted. Their relevance is
already concrete in the Hadamard direction: positive real Gram matrices
and exact sign constraints are central, while similarity alone discards
the metric problem.

\subsection{Xi III, Grillet, and Hall: symmetry with a terminology
trap}\label{xi-iii-grillet-and-hall-symmetry-with-a-terminology-trap}

Xi III has an actual averaging proof for finite complex representations
(Theorem 3.20, pp.~109--111, PDF pp.~115--117), followed by examples and
remarks about compact groups. Its §3.7 gives dual and tensor
representations (pp.~152--162, PDF pp.~158--168). Its module definitions
in §4.4 (pp.~186--188, PDF pp.~192--194) connect groups, algebras and
actions. The rendered pages show substantive construction coverage.

Grillet IX.7 proves the group-algebra correspondence and Maschke under
the field-characteristic condition (pp.~380--384, PDF pp.~392--396). It
initially calls a representation ``irreducible'' when it is
\textbf{indecomposable} (p.~381, PDF p.~393), explicitly making that
correspondence on p.~383. These notes, Xi and Hall use irreducible to
mean a nonzero representation with no proper nonzero invariant subspace.
Outside a semisimple setting, Grillet IX.7.1 must therefore not be read
as an unconditional complete reducibility theorem in the latter
terminology. This is a convention difference verified against a rendered
page, not an OCR mistake or a false theorem under his definition.

Hall §§4.3--4.7 proves unitary splitting, Schur's lemma, and the
\(\mathfrak{sl}_2\) classification (pp.~90--101, PDF pp.~103--114). Its
unipotent representation of \(\mathbb R\) demonstrates failure of
complete reducibility even over \(\mathbb C\). Xi III's §4.6 starts from
Lie module conditions and explicitly points outside that section for
enveloping algebra and PBW theory (pp.~215--216, PDF pp.~221--222).
Hall's later PBW, root-system and highest-weight chapters are verified
destinations from contents; their full proofs were not compared against
Xi's later examples or appendix. Thus ``Hall develops a systematic
further theory'' is supported, while ``all Xi representation examples
are replaced'' is not established.

For Hadamard/Williamson, the immediately useful representation theory is
finite cyclic characters and simultaneous splitting of the shift
operator. Full compact Lie theory and \(\mathfrak{sl}_2\) are enriching
branches that can be deferred for that project.

\subsection{Homological and geometric graduate books: bridges, not
substitutes}\label{homological-and-geometric-graduate-books-bridges-not-substitutes}

Hilton--Stammbach III.1--4 carefully distinguishes extension equivalence
from middle-module isomorphism and computes cyclic Ext groups
(pp.~84--99, PDF pp.~95--110). Its introduction explains why extensions
precede general derived functors. This is the right added construction
when invariant splitting fails. Over the field alone, positive Ext/Tor
vanish; complex homology remains meaningful. No need for advanced
homological machinery has been demonstrated in the sign-lifting
constraints of the Hadamard study.

Hartshorne II.5 and II.8 develop local module data and differential
universal properties (pp.~108--111, 172--174; PDF pp.~123--126,
187--189). They require sheaves, localization and local rings beyond
one-vector-space arguments. His ``A'' references make clear that some
algebra is cited rather than rederived. Lang I.9 and III.1 add
projectivity over Dedekind rings and integral trace duality (pp.~29, 57;
PDF pp.~43, 71). They require integral and local arithmetic information
that scalar extension forgets.

Hirsch IV.2.2 proves splitting for vector bundles on a paracompact base
using partitions of unity (p.~95, PDF p.~105). This is an important
extension of the linear splitting discussion, while his preface
deliberately chooses not to use tensors or differential forms (printed
p.~vii, PDF p.~8). An algebraic prerequisite ladder that makes him
depend on those subjects would misrepresent the actual book.

The practical decision is selective: use Xi for the concrete, metric,
and geometric language; choose Grillet or Hungerford for the module
spine; use Mac Lane to sharpen universal reasoning and Hilton--Stammbach
when a real exactness obstruction appears. Enter Hall, Lang, Hartshorne,
and Hirsch along their particular missing bridges. This is a branching
reading plan, not eleven sequential complete readings.

\section{Proof-level comparison: operators and multilinear
constructions}\label{proof-level-comparison-operators-and-multilinear-constructions}

This comparison supports a smaller route through the specified topics,
not replacement of any whole book. It is based on newly read contiguous
proof passages in the supplied editions: Xi Nanhua, \emph{Basic
Algebra}, volume II (Science Press, 2018); Thomas W. Hungerford,
\emph{Algebra} (the supplied Springer reprint of the 1974 text); and
Pierre Antoine Grillet, \emph{Abstract Algebra}, second edition (2007).
Page references below give printed pages first and one-based physical
PDF positions second. The offsets at these locations are respectively
+6, +19, and +12.

\textbf{P} means the passage itself was read, including its argument;
\textbf{V} means a rendered page was inspected; \textbf{T} means only a
location or declared scope was checked. An image request returning a
pointer or a warning does not earn V. The new proof readings below are
P. The exact glyphs in Xi's quotient-chain correction and
normalized-antisymmetrization displays remain \textbf{unverified
visually} because the available reader returned extracted text and image
pointers or explicit image-unavailable warnings; the prescribed download
helper was unavailable. The reconstructed arguments are written
independently in consistent notation and checked mathematically; they
are not transcriptions.

\subsection{Three existence proofs, one closely related uniqueness
mechanism}\label{three-existence-proofs-one-closely-related-uniqueness-mechanism}

\begin{longtable}[]{@{}
  >{\raggedright\arraybackslash}p{(\columnwidth - 4\tabcolsep) * \real{0.3333}}
  >{\raggedright\arraybackslash}p{(\columnwidth - 4\tabcolsep) * \real{0.3333}}
  >{\raggedright\arraybackslash}p{(\columnwidth - 4\tabcolsep) * \real{0.3333}}@{}}
\toprule\noalign{}
\begin{minipage}[b]{\linewidth}\raggedright
Source and inspected passage
\end{minipage} & \begin{minipage}[b]{\linewidth}\raggedright
Hypotheses and actual construction
\end{minipage} & \begin{minipage}[b]{\linewidth}\raggedright
What it supplies and what it costs
\end{minipage} \\
\midrule\noalign{}
\endhead
\bottomrule\noalign{}
\endlastfoot
Xi II §2.4, Theorem 2.44, pp.~68--69; PDF 74--75, P & Finite-dimensional
vector space over a field. Proves Cayley--Hamilton by dimension
induction. A proper invariant subspace and its quotient give two
annihilators whose product kills the whole space. If no proper nonzero
invariant subspace exists, every nonzero vector is cyclic; its
successive images form a basis, and the companion relation gives the
characteristic polynomial. & A linear-algebra proof that precedes Jordan
form. Requires quotient operators, the characteristic-polynomial product
for a block triangular matrix, and polynomial division, without a PID
module classification theorem. \\
Xi II §2.5, Theorems 2.51--2.53, pp.~72--74; PDF 78--80, P & A
polynomial \(fg\) annihilates \(T\), with \(\gcd(f,g)=1\). Evaluating a
Bézout identity splits the space as \(\ker f(T)\oplus\ker g(T)\). The
proof also proves that the other factor acts invertibly on each kernel.
Applying this to a split characteristic polynomial gives generalized
eigenspaces and their dimensions. & The essential coprime-primary
mechanism already appears in operator language. The two-factor lemma
itself does not need eigenvalues or algebraic closure. The subsequent
Jordan theorem needs splitting over the ground field. \\
Xi II §2.5, Theorems 2.56--2.57, pp.~75--77; PDF 81--83, P & Reduces to
nilpotence on each generalized eigenspace. Quotients by a line in the
kernel, uses induction to obtain chains downstairs, then changes their
lifted generators so that at most one chain absorbs the kernel line. & A
direct existence proof over any field for a nilpotent operator; split
characteristic polynomial suffices for general Jordan form. It teaches
the actual obstruction in lifting a quotient basis. \\
Hungerford IV.6, Lemma 6.8 and Theorem 6.9, pp.~222--223; PDF 241--242,
P & For a module killed by \(p^n\), chooses an element of order \(p^n\).
An elementary correction produces a nonzero cyclic submodule disjoint
from it whenever the module is not cyclic. A maximal disjoint submodule,
obtained using Zorn's lemma, is shown to be a complement by applying the
same argument in the quotient. Finitely generated primary modules then
decompose by induction on their number of generators. & A maximal-order
cyclic splitting argument over every PID, not just \(K[t]\). The
splitting lemma itself permits modules that are not finitely generated,
provided their \(p\)-power exponent is bounded. That extra scope should
not be erased by calling it merely another Jordan proof. \\
Grillet VIII.6, Theorems 6.1--6.3, pp.~336--339; PDF 348--351, P & First
proves submodules of free PID modules free, with a Zorn argument in the
arbitrary-rank statement. For finite free \(F\) and \(M\subseteq F\),
maximizes an ideal \(\mu(M)\) among images under maps \(F\to R\). Bézout
combinations force divisibility of all coordinates of a suitable
\(m\in M\); writing \(m=ae\) gives \(\mu(e)=1\). This splits both \(F\)
and \(M\), and rank induction yields a relative basis \(a_i e_i\), with
\(a_i\mid a_{i+1}\). Taking a free presentation gives invariant factors.
& The relative-basis theorem supplies integral presentations and
arbitrary-field rational forms in one route. It needs free modules,
quotients, principal ideals, the ascending chain condition for ideals of
a PID, and rank. For the finite route, the arbitrary-rank Zorn theorem
can be replaced by its finite-rank induction. \\
\end{longtable}

The two PID treatments also use Bézout to obtain primary decomposition:
Hungerford IV.6, Theorem 6.7, p.~222/PDF 241; Grillet VIII.6,
Proposition 6.5, p.~340/PDF 352. Their statements concern arbitrary
torsion modules over a PID. They factor an annihilator of each
individual element; they do not assume one polynomial annihilates the
entire module. Xi's finite-dimensional application has a common
annihilator available from Cayley--Hamilton. These are the same coprime
splitting mechanism at different scopes.

\subsubsection{The quotient-lifting step worth
retaining}\label{the-quotient-lifting-step-worth-retaining}

Here is a reconstruction of Xi's existence mechanism, with the delicate
exponent made explicit. Let \(N\) be nilpotent and choose
\(0\ne w\in\ker N\). In \(V/Kw\), take chains with lifted top vectors
\(e_j\) and lengths \(m_j\). Their defects satisfy \[
N^{m_j}e_j=\alpha_jw.
\] If every \(\alpha_j=0\), the lifted chains together with \(w\)
already give a chain basis. Otherwise choose \(r\) with \(\alpha_r\ne0\)
and \(m_r\) maximal among indices having nonzero defect. Keep \(e_r\),
extend its chain by one place, and, for \(j\ne r\) with
\(\alpha_j\ne0\), replace \[
e_j\quad\hbox{by}\quad
e'_j=e_j-\frac{\alpha_j}{\alpha_r}N^{m_r-m_j}e_r.
\] The exponent is nonnegative because of the choice of \(r\), and
\(N^{m_j}e'_j=0\). These are invertible elementary changes among the old
chain vectors; \(w=\alpha_r^{-1}N^{m_r}e_r\) is now in the extended
chain. This explains both why a quotient decomposition need not lift
unchanged and how the correction repairs it. It is a useful
finite-dimensional precursor to extension problems.

This construction and Hungerford's maximal-order splitting should not be
conflated. Xi starts from an arbitrary one-dimensional submodule and
repairs the lifted chains. Hungerford first chooses a cyclic submodule
of maximal possible order and proves that this chosen submodule splits.
The false assertion that \textbf{every} cyclic invariant subspace has an
invariant complement is used by neither proof.

\subsubsection{A second Xi construction missed by a theorem-only
comparison}\label{a-second-xi-construction-missed-by-a-theorem-only-comparison}

Xi continues beyond the existence theorem. On pp.~79--81/PDF 85--87, P,
he constructs chains from \[
W_k=\operatorname{im}N^{k-1}\cap\ker N.
\] One starts with a basis of the deepest nonzero \(W_k\), extends it
through the descending filtration, and chooses preimages under the
appropriate powers of \(N\). Applying high powers of \(N\) to a putative
dependence eliminates successive coefficients. The number of vectors is
justified by \[
\dim W_k=\dim\ker N^k-\dim\ker N^{k-1},
\] using the surjection \(N^{k-1}:\ker N^k\to W_k\), and then
telescoping. This supplies a practical basis-finding procedure and an
independence proof. The PID theorem's existence statement alone does not
replace that skill.

\subsubsection{Exactly how uniqueness
agrees}\label{exactly-how-uniqueness-agrees}

Xi pp.~77--78/PDF 83--84, P, counts blocks by differences of image
dimensions. Grillet VIII.6, Lemma 6.7 and the completion of Theorem 6.3,
pp.~340--342/PDF 352--354, P, first counts primary cyclic summands using
\(M/pM\) over \(R/(p)\), then applies induction to \(pM\). Hungerford
IV.6, Lemmas 6.10--6.11 and Theorem 6.12, pp.~224--225/PDF 243--244, P,
explains how to adapt the earlier abelian-group uniqueness argument.
That referenced argument was also read: II.2, Theorem 2.6,
pp.~78--81/PDF 97--100, P. It detects the number of summands through
\(p\)-torsion, then multiplies by a power of \(p\) to contradict the
first discrepancy between exponent lists. Over a general PID,
residue-field dimension replaces finite cardinality.

An economical common formulation, reconstructed from these arguments, is
this. For a finite primary module \(M\simeq\bigoplus_jR/(p^{a_j})\), set
\[
d_k=\dim_{R/(p)}p^{k-1}M/p^kM.
\] Each summand contributes one dimension precisely when \(a_j\ge k\).
Thus \(d_k-d_{k+1}\) is the number of summands of exact exponent \(k\).
For \(R=K[t]\) and \(p=t-\lambda\), this becomes Xi's rank formula \[
b_k=\operatorname{rank}(T-\lambda I)^{k-1}
-2\operatorname{rank}(T-\lambda I)^k
+\operatorname{rank}(T-\lambda I)^{k+1}.
\] Other primary components contribute equal ranks at all three stages
and cancel. Therefore three full uniqueness proofs are unnecessary for
theorem coverage; one proof of the residue-field layer count, followed
by this translation, supplies the common invariant. For an irreducible
polynomial of degree greater than one, \(K\)-dimensions of the layers
must be divided by \(\deg p\); pretending that every irreducible is an
eigenvalue loses the rational case.

\subsection{The replacement map has a genuine
boundary}\label{the-replacement-map-has-a-genuine-boundary}

Grillet VIII.7, pp.~342--346/PDF 354--358, P, constructs the
\(K[X]\)-action, identifies invariant subspaces as submodules, obtains a
nonzero annihilator from finite dimensionality of
\(\operatorname{End}_K(V)\), and then applies the elementary-divisor
theorem. Polynomial division proves the cyclic basis; replacing \(T\) by
\(T-\lambda I\) converts it into a Jordan chain. His displayed main
Jordan theorem assumes an algebraically closed field. Its proof also
works when the relevant minimal polynomial splits, but that weaker
hypothesis is a mathematical sharpening of the displayed statement.

Hungerford VII.4, pp.~355--365/PDF 374--384, P, carries this translation
through both invariant factors and elementary divisors. It proves
rational canonical form, primary rational form, and the split Jordan
specialization. Proposition 4.9, pp.~363--364/PDF 382--383, is an
additional computational bridge: it presents the operator module by
\(tI-A\), proves exactness by reducing polynomial vectors to constant
vectors, and identifies the nonunit Smith entries as the invariant
factors. This is more than merely renaming Jordan blocks.

\begin{longtable}[]{@{}
  >{\raggedright\arraybackslash}p{(\columnwidth - 4\tabcolsep) * \real{0.3333}}
  >{\raggedright\arraybackslash}p{(\columnwidth - 4\tabcolsep) * \real{0.3333}}
  >{\raggedright\arraybackslash}p{(\columnwidth - 4\tabcolsep) * \real{0.3333}}@{}}
\toprule\noalign{}
\begin{minipage}[b]{\linewidth}\raggedright
Learning goal
\end{minipage} & \begin{minipage}[b]{\linewidth}\raggedright
Coverage that may be compressed after mastery
\end{minipage} & \begin{minipage}[b]{\linewidth}\raggedright
Material still needed
\end{minipage} \\
\midrule\noalign{}
\endhead
\bottomrule\noalign{}
\endlastfoot
Primary decomposition when an annihilator is known & Xi 2.51 and the
finite \(K[t]\)-specializations of the two PID primary-decomposition
proofs & One explicit Bézout projector calculation; the distinction
between coprime and noncoprime factors; the elementwise formulation for
arbitrary torsion modules if that scope is wanted \\
Jordan existence over a splitting field & One of the two existence
routes suffices for the theorem & Xi's lift-correction argument is a
valuable alternate proof; his filtration method should remain a
basis-construction exercise \\
Jordan uniqueness and block recovery & The three filtration proofs can
be replaced by one common layer argument and translation & Actual
rank/nullity calculations; constructing a conjugating basis; field and
splitting hypotheses \\
Rational form over an arbitrary field & Xi's split Jordan proof is
insufficient & A PID cyclic-decomposition theorem plus polynomial
division; a nonsplit example; the \(tI-A\) presentation if effective
invariant-factor computation is required \\
Finitely generated modules over a PID & Reading both complete existence
proofs is unnecessary for the same classification theorem & Choose one
proof; retain Hungerford's bounded-primary splitting lemma as an
optional stronger result and Grillet's relative-basis theorem when
integral presentations are needed \\
Cayley--Hamilton & Once independently proved, repeated derivations can
be shortened & Preserve a coherent dependency order: Xi proves it
\textbf{before} Jordan; Hungerford derives it from rational form in
VII.5, p.~367/PDF 386; Grillet derives it after Jordan and passes to an
algebraic closure, pp.~345--346/PDF 357--358 \\
\end{longtable}

In particular, a route that uses Xi's characteristic-polynomial
splitting argument to prove Jordan and then uses Jordan as its only
proof of Cayley--Hamilton would be circular. One may instead follow Xi's
preceding quotient proof, or follow the PID route and establish torsion
through linear dependence among \(I,T,T^2,\ldots\).

The exercises are not interchangeable. Xi §2.5, pp.~83--84/PDF 89--90,
P, includes explicit Jordan bases, powers, matrix-root equations, the
polynomial semisimple/nilpotent decomposition, a commutator condition
forcing nilpotence, and the finite exponential/logarithm correspondence.
Hungerford VII.4 exercises, pp.~365--366/PDF 384--385, and VII.5
exercises, pp.~369--370/PDF 388--389, P, add invariance of invariant
factors under field extension, descent of similarity, rational-form
classification with nonsplit polynomials, a commutative-ring
Cayley--Hamilton exercise, and norm/trace connections. These are
separate destinations, even where some standard tasks overlap. This
reading does not claim that every exercise in the chapters has been
solved or that every application occurs in only one book.

\subsection{Tensor products: the universal property is already proved in
Xi}\label{tensor-products-the-universal-property-is-already-proved-in-xi}

Xi II §6.2, Definition 6.11 and Theorem 6.15, pp.~273--275/PDF 279--281,
P, defines a tensor product by a bilinear map whose pairs of basis
vectors form a basis. He constructs such a space, proves uniqueness by
the basis map, and proves the universal property by assigning the images
of those basis tensors. This applies to arbitrary vector spaces with
bases: the relevant expansions have finite support. The
finite-dimensional qualifications on the adjacent
\(\operatorname{Hom}\)-and-dual identifications should not be
transferred to the universal property itself.

Hungerford IV.5, pp.~207--211/PDF 226--230, P, and Grillet XI.5,
pp.~434--438/PDF 446--450, P, build the tensor product of a right and a
left module over an arbitrary ring. Both start with the free abelian
group on pairs and quotient by additivity and balancing relations. A
balanced map extends to the free group, kills the relation subgroup, and
therefore descends uniquely. Grillet motivates this with the same
vector-space basis calculation that Xi uses.

The additional content is the applicability to modules without bases,
noncommutative sidedness, and the subsequent exactness theory. It is
inaccurate to say that one needs a graduate text merely to obtain tensor
universality. Once the quotient construction and its factorization proof
have been learned in either graduate source, the other construction
proof is logically redundant. Xi's coordinate realization is still
useful as the explicit free case.

Hungerford's right-exactness proof, p.~210/PDF 229, P, deserves separate
retention. Given a surjection \(B\to C\), it constructs the inverse map
from \(D\otimes C\) into the quotient of \(D\otimes B\) by the preceding
image; independence of a chosen lift is exactly what the quotient
enforces. This teaches a use of universality that is not supplied merely
by the statement that basis tensors form a basis.

\subsection{Exterior algebra: related bases, different constructions and
scope}\label{exterior-algebra-related-bases-different-constructions-and-scope}

Xi II §6.4, pp.~292--299/PDF 298--305, P, assumes characteristic zero
and finite dimension at the section opening. He realizes exterior powers
inside tensor powers as alternating tensors and defines the product by
normalized antisymmetrization. The associativity proof, pp.~293--295/PDF
299--301, proves that antisymmetrizing within a block and then globally
agrees with global antisymmetrization. The ordered-wedge basis proof,
pp.~296--297/PDF 302--303, first eliminates repetitions and sorts
indices for spanning, then isolates the coefficient of each increasing
tensor word for independence. This is a complete basis proof, not just a
dimension assertion.

Grillet XIII.4, pp.~524--527/PDF 536--539, P, instead takes \[
\Lambda(M)=T(M)/(v\otimes v:v\in M)
\] over a commutative ring with identity. Universality follows by
killing the defining ideal. He identifies the degree-\(n\) relation
submodule with the span of tensors having repeated entries and proves
the universal property for alternating \(n\)-linear maps. For free
\(M\), pp.~525--526/PDF 537--538 prove the ordered-wedge basis by
describing relations among tensor basis words, grouping the terms caused
by repeated general vectors, and splitting off a free complement spanned
by increasing words.

Thus both sources prove spanning \textbf{and} independence. Xi's
independence comes from coefficients in an alternating subspace;
Grillet's comes from an explicit quotient and a complement to its
relation submodule. Grillet does not need division by \(n!\),
characteristic zero, or finite rank for the ordered-basis statement.

A concise bridge between them is an original deduction: when \(n!\) is
invertible, normalized antisymmetrization descends from the quotient
exterior power to the alternating subspace. Ordered wedge classes map to
normalized alternating basis tensors, so the map is an isomorphism in
the free finite-dimensional case. Xi's product is normalized
accordingly. This does not justify identifying the quotient with an
alternating subspace in arbitrary characteristic. In characteristic two,
the equation \(u\otimes v=-v\otimes u\) alone does not impose
\(u\otimes u=0\); the quotient construction must retain the latter
relation.

For a general algebra route, Grillet's quotient construction can be the
single main proof. Retain Xi's alternating realization as a
characteristic-zero model and as practice in coefficient calculations.
Both sources connect the top exterior power with determinants; the proof
of the ordered basis does not by itself replace exercises with minors,
decomposability, or geometric interpretation. No absence claim about
exterior algebra in Hungerford is made: a literal search found no match,
which is not proof of noncoverage.

\subsection{A smaller coherent route and stopping
tests}\label{a-smaller-coherent-route-and-stopping-tests}

For the intended bridge from undergraduate linear algebra to modules,
use this order:

\begin{enumerate}
\def\labelenumi{\arabic{enumi}.}
\tightlist
\item
  Learn invariant subspaces, quotients, polynomial evaluation, and
  Bézout. Prove the two-factor kernel decomposition and compute its
  projectors. If using Xi's route, include the pp.~68--69
  Cayley--Hamilton proof before applying the characteristic polynomial.
\item
  Work through Xi's quotient-lifting Jordan construction once, then
  prove the common layer-count uniqueness formula. Construct at least
  one Jordan basis by the \(W_k\) filtration. Further repetitions of
  Jordan existence and uniqueness may then become short translations.
\item
  Learn the operator/\(K[t]\)-module dictionary. For broader
  classification, use \textbf{one} finite PID structure proof: Grillet's
  relative-basis proof is the natural choice when integer lattices,
  presentations, and rational form are downstream goals. Translate the
  result to operators with Hungerford VII.4, including the \(tI-A\)
  presentation. Hungerford's complete alternative IV.6 proof becomes
  optional proof comparison; its stronger bounded-primary lemma remains
  available when that generality is needed.
\item
  Use one balanced-tensor quotient construction and its universal
  property. Derive the free-basis model, induced maps, and a
  right-exactness example. Then use the exterior quotient and
  ordered-basis proof, followed by the characteristic-zero alternation
  comparison.
\end{enumerate}

These are mastery gates, not estimated reading times:

\begin{itemize}
\tightlist
\item
  \textbf{Splitting gate:} for coprime \(f,g\), construct both
  projectors from a Bézout identity and check their images, kernels,
  sum, and idempotence. Explain why a vector-space complement need not
  be \(T\)-invariant.
\item
  \textbf{Chain gate:} reconstruct the correction \(e'_j\) without
  assuming a quotient basis automatically lifts to chains. From
  nullities distinguish nilpotent block sizes \((4,2,2)\) and
  \((4,3,1)\), although dimension, characteristic polynomial, minimal
  polynomial, and eigenspace dimension agree. Their second nullities are
  respectively \(6\) and \(5\).
\item
  \textbf{Rational gate:} treat \(K[t]/((t^2+1)^2)\) over
  \(K=\mathbb R\), then after extending to \(\mathbb C\). Explain which
  invariant factors remain the same and which elementary divisors split.
  Produce a companion matrix and a valid module presentation; merely
  writing a complex Jordan matrix does not pass.
\item
  \textbf{Tensor gate:} define a map out of a tensor product by checking
  balancing and uniqueness; compute \((R/I)\otimes_R M\simeq M/IM\).
  Explain the failure of injectivity after tensoring
  \(2\mathbb Z\hookrightarrow\mathbb Z\) with \(\mathbb Z/2\mathbb Z\).
\item
  \textbf{Exterior gate:} prove ordered wedges independent, not just
  spanning. Explain the role of \(n!\) in Xi's model, and construct
  \(\Lambda^2(K^2)\) in characteristic two using the repeated-vector
  relation.
\end{itemize}

A reader who passes these gates can compress repeated theorem coverage
in the inspected portions. The route does not replace metric/spectral
geometry, affine or projective geometry, tensor contraction practice,
the remaining exercise collections, general ring and module theory, or
homological algebra. Its economy comes from choosing one proof per
common structural theorem and keeping the alternate constructions that
teach a distinct method.

\subsection{Evidence boundary}\label{evidence-boundary}

The newly inspected operator passages are Xi II PDF 74--90, Grillet PDF
348--358, Hungerford PDF 237--245 and 374--390, with the referenced
Hungerford uniqueness passage PDF 97--100. The new multilinear passages
are Xi II PDF 279--281 and 298--305; Grillet PDF 446--450 and 536--539;
Hungerford PDF 226--230. The comparison does not promote these ranges
into an audit of the surrounding books. Source files and retrieval
identifiers are private and are not included here.

\section{A selective reconstruction of Xi I and III: proof routes,
prerequisites, and stopping
tests}\label{a-selective-reconstruction-of-xi-i-and-iii-proof-routes-prerequisites-and-stopping-tests}

This comparison asks a local question: which already familiar parts of
Xi can a returning undergraduate reconstruct from a smaller collection
of proofs, and what would that reconstruction omit? The answer is
conditional on competence. It is not a replacement claim about any
complete book.

The broadened source reading changes two tempting conclusions. Xi III
does \textbf{not} stop at the complex inner-product proof of complete
reducibility: it also gives the general averaged-projection proof of
Maschke. Xi III does \textbf{not} introduce modules only as terminology:
its module section proves the submodule theorem over a PID and uses a
determinant argument to characterize integral elements. A graduate text
is useful where it supplies a missing proof or a reusable
generalization; its higher placement in a curriculum does not make every
local proof shorter.

Evidence follows the existing ledger: \textbf{P} = actual
statement/proof text read; \textbf{V/P} = rendered source page and proof
inspected; \textbf{T} = contents/location only. Physical PDF pages are
one-based. The checked body offsets remain Xi I/III +6, Grillet +12,
Hungerford +19, and Lang +14. The supplied Xi III scan has 361 physical
pages, although the Library text reader reports only 150. Local
rendering recovered the relevant pages; source bytes and rendered
evidence remain private.

\subsection{The resultant: Xi's matrix construction already supplies an
organizing
principle}\label{the-resultant-xis-matrix-construction-already-supplies-an-organizing-principle}

\textbf{Xi I §7.6, printed 190--193 / PDF 196--199, V/P.} Definition
7.76 starts with the Sylvester determinant and permits zero leading
coefficients in the displayed, fixed-degree coefficient lists.
Proposition 7.77 does not merely assert that a resultant detects a
common factor. Its proof first converts the question into a
bounded-degree relation

\[
f u+g v=0,\qquad \deg u<m,\quad \deg v<n,
\]

where the displayed bounds for \(f,g\) are \(n,m\). Equating
coefficients gives a square linear system whose determinant is the
resultant, up to the transpose convention visible in the proof. Euclid's
lemma supplies the reverse implication when \(f,g\) are coprime. Thus
the proof already connects polynomial divisibility, a kernel, and
determinant singularity.

There is a real hypothesis issue here, rather than a cosmetic
convention. With the nominal coefficient lengths fixed, Xi's criterion
includes the alternative that \textbf{both displayed leading
coefficients vanish}. For instance, if \(f=x+1\) and \(g=x+2\) are both
padded to nominal degree two over \(\mathbb Q\), the resulting
determinant is zero even though they are coprime. A reconstruction that
silently replaces nominal degrees by actual degrees changes the
statement.

Proposition 7.78, printed 191--192 / PDF 197--198, then obtains the
root-product formula by replacing the roots of \(f\) by independent
indeterminates and studying the polynomial \(\operatorname{Res}(F,g-y)\)
in \(y\). Its roots are the values \(g(x_i)\); specialization gives the
desired identity. The proof uses polynomial arithmetic in a
rational-function field, not a prior theory of norms or resultants.
Proposition 7.79 derives the discriminant identity. The exercises at
printed 193 / PDF 199 include multiplicativity, the discriminant of a
product, cyclotomic-looking examples, and elimination in two variables.
Reading only the opening matrix definition misses these connections.

\textbf{Grillet IV.7, printed 176--179 / PDF 188--191, P.} The direction
is reversed. Grillet defines the resultant by the product of differences
between roots in an algebraic closure. Proposition 7.2 then proves that
it equals the Sylvester determinant, assuming the leading coefficients
are nonzero. The proof introduces universal coefficient and root
variables over \(\mathbb Z\), shows that each root difference divides
the specialized determinant, controls the total degree, and determines
the remaining integer constant. The root-product definition makes the
vanishing criterion immediate; the longer work is proving that the
invariant is a coefficient polynomial.

These are overlapping theorems with different constructions. Xi offers
the more immediate route from a concrete linear system to elimination.
Grillet offers a universal polynomial identity that survives
specialization. Neither route supplies the other perspective merely by
stating the same resultant formula. Hungerford was searched for this
topic but no exact counterpart was established; that is an
access/coverage limit, not evidence that the book contains no relevant
material.

\textbf{A compact reconstruction.} Define

\[
S_{f,g}:k[x]_{<m}\oplus k[x]_{<n}\longrightarrow k[x]_{<m+n},
\qquad (u,v)\longmapsto fu+gv
\]

for nonzero polynomials of actual degrees \(n,m\). Identify its matrix
with the Sylvester matrix, with a recorded ordering/sign convention.
Prove \(\ker S_{f,g}\neq0\) exactly when \(\gcd(f,g)\neq1\). Then derive
the product formula using either Xi's generic-root argument or Grillet's
universal identity. This is a reorganization of inspected proofs, not a
newly discovered structure absent from Xi.

The following are useful stopping tests before skipping these pages:

\begin{enumerate}
\def\labelenumi{\arabic{enumi}.}
\tightlist
\item
  Derive the matrix from the map rather than memorizing its rows, and
  explain why the domain and codomain have the same dimension.
\item
  Produce the padded-degree counterexample above and state the corrected
  criterion.
\item
  Explain why \(x-1\) and \(x+1\) are comaximal in \(\mathbb Q[x]\) but
  not in \(\mathbb Z[x]\). Their resultant is a nonzero nonunit over
  \(\mathbb Z\); evaluation at \(1\) and \(-1\) sends
  \(\mathbb Z[x]/(x^2-1)\) to the pairs of integers with equal parity,
  not to all of \(\mathbb Z^2\).
\item
  Recover the discriminant identity, including its sign, and compute one
  cubic directly.
\end{enumerate}

\textbf{Narrow, visually verified erratum.} Xi I Example 7.80, printed
193 / PDF 199, displays \(-2p^3-27q^2\) for the discriminant of
\(x^3+px+q\). The correct expression is \(-4p^3-27q^2\). The special
case \(x^3+x\), with roots \(0,i,-i\), gives discriminant \(-4\),
whereas the displayed expression gives \(-2\). The general cubic
exercise on printed 190 / PDF 196 has the correct coefficient
\(-4a_0a_2^3\). This is a local example error; it does not invalidate
the inspected determinant criterion or product-formula argument.

\subsection{Finite fields: the short route must include construction and
subfields}\label{finite-fields-the-short-route-must-include-construction-and-subfields}

\subsubsection{What the inspected books actually
prove}\label{what-the-inspected-books-actually-prove}

\textbf{Xi III §1.4 Theorem 1.21, printed 34 / PDF 40, V/P.} The finite
multiplicative subgroup theorem uses the maximum order of an element and
the polynomial root bound. The step asserting that every element is a
root of \(X^r-1\), with \(r\) the maximum order, uses the preceding
finite-abelian-group structure, or can be supplied independently by the
elementary least-common-multiple lemma below. It is not valid for
arbitrary finite groups merely because \(r\) is a maximum.

\textbf{Xi III §5.6, printed 260--266 / PDF 266--272, V/P.} Theorem 5.40
recalls multiplicative cyclicity. Theorems 5.41--5.43 establish the
prime-power order, uniqueness, and existence. Existence takes a
splitting field of \(X^{p^n}-X\), uses derivative \(-1\) for distinct
roots, and verifies that the roots form a subfield. Proposition 5.44 at
printed 261 / PDF 267 proves the subfield criterion directly: necessity
uses the tower dimension; sufficiency uses divisibility of \(X^{p^r}-X\)
into \(X^{p^n}-X\) when \(r\mid n\). Theorem 5.45 gives extensions of
any finite degree.

At printed 262 / PDF 268, Xi identifies the automorphism group using
Frobenius and the already established Galois theorem. It then obtains
irreducible polynomials in every degree. Printed 262--266 / PDF 268--272
develops Möbius inversion, cyclotomic polynomials, and the count of
monic irreducible polynomials over \(\mathbb F_q\). Thus ``replace §5.6
by a finite-field existence theorem'' would omit much of the section.
The exercises beginning on printed 266 / PDF 272 require minimal
polynomials and explicit factorizations, not just theorem recall. The
remaining exercise pages 267--269 were not audited.

\textbf{Grillet IV.1.6, printed 157 / PDF 169, P.} Cyclicity is proved
by decomposing a finite abelian group into its prime-primary subgroups.
In each subgroup, a maximal-order element and the root bound force the
whole subgroup to be cyclic; the relatively prime orders then combine.
This avoids needing the full invariant-factor classification at that
point, but it still relies on the earlier primary decomposition.

\textbf{Grillet V.1.3, printed 192--193 / PDF 204--205, P.} Existence
and uniqueness again use \(X^{p^n}-X\), derivative \(-1\), root-set
closure, and uniqueness of splitting fields. Subfield
existence/uniqueness, the divisibility criterion, and the intersection
formula are \textbf{exercises} on printed 193, not locally supplied
proofs. A ``Grillet replaces Xi §5.6'' claim must either solve these
exercises or admit that it has traded Xi's written proof for work
assigned to the reader.

\textbf{Hungerford V.5, printed 279--282 / PDF 298--301, P.} Theorem 5.3
proves cyclicity using the invariant-factor theorem for finite abelian
groups: all elements satisfy the polynomial determined by the largest
invariant factor, and the root bound rules out more than one nontrivial
factor. Proposition 5.6 and Corollary 5.7 provide the same
splitting-field construction and uniqueness. Frobenius is introduced in
Lemma 5.5; Proposition 5.10 identifies the cyclic Galois group using the
fundamental theorem of Galois theory. The construction of fields of
orders 9 and 25 is assigned as Exercise 5. Irreducible-polynomial
existence is also assigned as a short proof. A local shortcut is
possible, but reading these paragraphs as written imports both
finite-abelian classification and general Galois theory.

\subsubsection{A smaller undergraduate sequence, with its imports made
explicit}\label{a-smaller-undergraduate-sequence-with-its-imports-made-explicit}

For applications needing only finite fields, their cyclic multiplicative
groups, Frobenius, and subfields, the following sequence is sufficient:

\begin{enumerate}
\def\labelenumi{\arabic{enumi}.}
\tightlist
\item
  \textbf{Polynomial arithmetic and one-root adjunction.} Prove the
  field property of \(k[x]/(h)\) for irreducible \(h\), the degree/root
  bound, and the derivative criterion. Construct a splitting field of a
  fixed polynomial by finitely many successive adjunctions. Prove
  uniqueness of this finite splitting field by extending an isomorphism
  one root at a time. An algebraic closure, Zorn's lemma, and the full
  Galois correspondence are unnecessary for this finite task.
\item
  \textbf{Root-set field construction.} Carry out the proof for
  \(X^{p^n}-X\) completely. Do not start by choosing an irreducible
  polynomial of degree \(n\) unless its existence has already been
  proved; that would conceal the main existence problem.
\item
  \textbf{Cyclicity from element orders.} In a finite abelian group,
  elements of orders \(a,b\) yield an element of order
  \(\operatorname{lcm}(a,b)\): for each prime, retain from one of them a
  power with the largest required prime-power order, then multiply these
  commuting components of pairwise coprime orders. Iteration shows that
  the exponent is attained by an element. Applied to a finite
  multiplicative subgroup of a field, the root bound for \(X^e-1\) gives
  \(|G|\le e\), while the element of order \(e\) gives \(e\le |G|\). No
  classification theorem is needed.
\item
  \textbf{Subfields by polynomial divisibility.} Use Xi's direct proof
  of Proposition 5.44. Uniqueness follows because a subfield of order
  \(p^r\) must be exactly the roots of \(X^{p^r}-X\) inside the larger
  field. Intersections follow by the same root-set description or the
  Euclidean algorithm for these polynomials.
\item
  \textbf{Frobenius without full Galois theory.} In \(\mathbb F_{p^n}\),
  Frobenius has order \(n\): its \(n\)-th power is the identity and a
  smaller positive power cannot fix all \(p^n\) elements by the root
  bound. Choose a generator \(\alpha\) of the multiplicative group, so
  the field is \(\mathbb F_p(\alpha)\). Its minimal polynomial has
  degree \(n\); an automorphism is determined by its image of
  \(\alpha\), with at most \(n\) choices. The \(n\) Frobenius powers
  therefore exhaust the automorphisms.
\end{enumerate}

This sequence is genuinely smaller in prerequisites than importing a
full abstract-algebra chapter. It is not much smaller than the relevant
core pages of Xi once those earlier facts are available. It also does
not replace the Möbius-counting and cyclotomic portions.

\textbf{Stopping tests.} Construct
\(\mathbb F_9=\mathbb F_3[u]/(u^2+1)\); prove the defining polynomial
irreducible; show \(1+u\) has order eight; compute Frobenius on
\(a+bu\). Explain why \(\mathbb F_{p^6}\) has a unique subfield of order
\(p^2\) and none of order \(p^4\). For the counting portion,
independently derive

\[
q^m=\sum_{d\mid m}dN_d,\qquad
N_m=\frac1m\sum_{d\mid m}\mu(m/d)q^d,
\]

where \(N_d\) counts monic irreducible degree-\(d\) polynomials. A
reader who can only quote cyclicity and field uniqueness has not yet
passed these construction tests.

\subsection{Representations and modules: Xi already contains both
averaging
proofs}\label{representations-and-modules-xi-already-contains-both-averaging-proofs}

\textbf{Xi III §3.2, printed 109--115 / PDF 115--121, V/P.} Theorem 3.20
averages a positive Hermitian form to make a finite-dimensional complex
representation unitary. The accompanying remarks distinguish real
orthogonal representations, compact-group integration, and the practical
problem of finding a useful matrix change of basis. Example 3.22,
printed 111--113 / PDF 117--119, explicitly splits the permutation
representation of \(S_3\), writes the two-dimensional matrices, and
changes bases to geometric and eigenvector coordinates.

The crucial additional pages are printed 113--115 / PDF 119--121.
Theorem 3.24 states complete reducibility in characteristic zero or
characteristic not dividing \(|G|\), and gives \textbf{two proofs}:

\begin{itemize}
\item
  Over \(\mathbb C\), the averaged Hermitian form makes the orthogonal
  complement of an invariant subspace invariant.
\item
  Over any allowed field, start with an arbitrary vector-space
  projection \(\theta\) onto the invariant subspace \(U\), and average
  its conjugates:

  \[
  P(v)=|G|^{-1}\sum_{g\in G}g\,\theta(g^{-1}v).
  \]

  Xi proves \(P|_U=1\), constructs \(V=U\oplus\ker P\), and checks
  invariance of the kernel by reindexing. The text then distinguishes
  the nonuniqueness of individual summands from the later uniqueness of
  isotypic data.
\end{itemize}

Therefore a comparison based only on printed 109--111 would understate
Xi's algebraic proof content.

\textbf{Grillet IX.7, printed 381--386 / PDF 393--398, P.} The group
algebra is constructed from the group basis and its universal extension
property. Propositions 7.2--7.3 prove the correspondence between
representations and modules, including equivalence/isomorphism and
invariant-subspace/submodule compatibility. Theorem 7.4, printed
383--384 / PDF 395--396, averages a projection just as Xi does, with an
equivalent placement of inverses. Its additional organization is to
conclude semisimplicity of \(k[G]\), and then use the preceding theory
of semisimple rings for classification consequences.

There are two different reading tasks here. The local
representation/module dictionary and projection proof require only basic
linear algebra, finite sums, module definitions, and invertibility of
\(|G|\). The subsequent matrix-algebra decomposition, counts of
irreducibles, and center arguments invoke the developed semisimple-ring
theory. They do not arrive for free from the one-page averaging proof.

Terminology must also be translated: Grillet's definition on printed 381
uses ``irreducible'' for what standard module language calls
\textbf{indecomposable}. The text explicitly identifies this with
indecomposability on printed 383. This does not contradict its results,
but it can produce false comparisons with Xi's
absence-of-proper-invariant-subspaces definition unless the
semisimplicity hypothesis has already been imposed.

\textbf{Hungerford IX.5, Proposition 5.8, printed 454--455 / PDF
473--474, P.} This is a materially different proof, labeled a sketch.
The regular representation realizes elements of the group algebra as
matrices. An element of the Jacobson radical can be normalized to have
identity coefficient one. The radical is nilpotent because the algebra
is Artinian, so its matrix has trace zero; the regular-representation
calculation instead makes that trace \(|G|\), giving a contradiction
when \(|G|\neq0\) in the field. For the converse, the sum of all group
elements generates a nonzero square-zero ideal when the characteristic
divides the group order.

This is an efficient synthesis \textbf{after} radicals and Artinian ring
theory. It is not the shorter elementary proof for a reader who has not
learned those theories. The following matrix-algebra decomposition also
invokes preceding structure theorems.

\subsubsection{What the module chapter adds---and where a gap
remains}\label{what-the-module-chapter-addsand-where-a-gap-remains}

\textbf{Xi III §4.4, printed 186--196 / PDF 192--202, V/P.} Besides
module definitions, quotients, kernels, annihilators and Schur's lemma,
the section develops finite free modules and their mapping property.
Proposition 4.46 discusses uniqueness of rank, including the distinction
between domains, commutative rings, and general noncommutative rings.
Theorem 4.47, printed 193 / PDF 199, proves that a submodule of a
finite-rank free module over a PID is free, by projecting onto a
rank-one quotient and lifting a basis. This is a genuine construction.

At printed 194 / PDF 200, the general cyclic decomposition with nested
annihilators is stated, with references to the previously developed
abelian-group and \(k[x]\)-operator cases. A complete local proof of the
general invariant-factor theorem is not supplied there. The exercises at
printed 196 / PDF 202 then explicitly ask the reader to recover Jordan
form from the \(\mathbb C[x]\)-module theorem and to generalize earlier
abelian-group results to arbitrary PIDs.

This is a defensible location for a focused supplement from Grillet
VIII.6, whose fuller proof is documented in the existing source ledger,
or from another fully proved PID treatment. It does \textbf{not} justify
saying that the module viewpoint itself is missing from Xi. Nor should
one advertise Xi III's brief statement plus exercise as an already
completed proof of the general classification.

\subsubsection{A compact representation sequence and its
limits}\label{a-compact-representation-sequence-and-its-limits}

For cyclic Fourier methods and finite-group invariant complements, a
small sequence suffices:

\begin{enumerate}
\def\labelenumi{\arabic{enumi}.}
\tightlist
\item
  Define representations, intertwiners, invariant subspaces, and the
  group algebra; prove the representation/module dictionary by extending
  the action linearly.
\item
  Prove the averaged-projection lemma directly---Xi printed 114--115 or
  Grillet printed 383--384 both work.
\item
  Prove Schur's lemma from kernels and images, adding the algebraically
  closed scalar conclusion only when required.
\item
  For a cyclic group, identify \(k[C_n]\cong k[x]/(x^n-1)\). Over a
  splitting field with characteristic not dividing \(n\), derive the
  eigenspace projectors by polynomial CRT or the finite geometric-series
  identity. Over \(\mathbb Q\), retain irreducible cyclotomic blocks
  rather than claiming all constituents are one-dimensional.
\item
  Work one nonabelian example, such as Xi's \(S_3\) permutation-module
  split, to test what cyclic intuition fails to capture.
\end{enumerate}

The first four steps need neither Jacobson radicals nor
Wedderburn--Artin theory, compact integration, Lie algebras, or a full
character table theory. They also do not replace those topics, Xi's
rotation-group and \(SU(2)\) examples, or the representation appendix.
Hall's previously sampled unitary and Lie-theoretic pages remain a
separate extension; no new Hall proof audit was made for this
comparison.

\textbf{Stopping tests.} Prove that the averaged map is a projection and
commutes with the action; identify exactly where \(|G|^{-1}\) is used.
In characteristic \(p\), use \(C_p\) acting by \(I+N\) on a
two-dimensional space, with \(N\neq0\) and \(N^2=0\), to show an
invariant line need not have an invariant complement. Over
\(\mathbb C\), distinguish the existence of a common invariant Hermitian
form from separate diagonalizability of each group element. Finally
split the \(S_3\) permutation representation and exhibit an actual
intertwining matrix. These test construction and failure cases, not just
theorem names.

\subsection{Integrality and cyclotomic fields: retain the elementary
bridge before
Lang}\label{integrality-and-cyclotomic-fields-retain-the-elementary-bridge-before-lang}

\textbf{Xi III §4.4, Lemma 4.49 and Theorem 4.50, printed 194--195 / PDF
200--201, V/P.} The text proves that an element \(t\) of a domain is
integral over \(\mathbb Z\) exactly when \(\mathbb Z[t]\) is finitely
generated as a \(\mathbb Z\)-module. The reverse direction writes
multiplication by \(t\) on a finite generating family and uses the
vanishing of the determinant of \(tI-A\), whose polynomial is monic. It
then deduces closure of integral elements under addition and
multiplication using finite generation and the PID submodule result.
Character values, being sums of roots of unity, are connected back to
algebraic integers.

This already gives a concrete bridge from matrix algebra and modules to
integrality. A reconstructed version can phrase the determinant step as
the adjugate identity; no Dedekind domain or valuation is needed. It is
especially useful before passing from rational spectral data back to
integer coefficients.

\textbf{Xi III §§5.6--5.7, printed 264--265 and 270--272 / PDF 270--271
and 276--278, V/P.} Example 5.47 constructs cyclotomic polynomials
through the factorization of \(x^n-1\), proves integral coefficients by
monic division and induction, and postpones general irreducibility.
Theorem 5.50 later proves irreducibility over \(\mathbb Q\) by Gauss's
lemma and reduction modulo a prime \(p\nmid n\): if a primitive root's
\(p\)-th power belonged to a different factor, the reductions of the
factors would share a factor, contradicting squarefreeness of \(x^n-1\)
modulo \(p\). Iterating over prime factors of an exponent reaches all
primitive roots. This is an undergraduate polynomial argument with
explicit prerequisites.

\textbf{Lang, Algebraic Number Theory IV.1, printed 71--75 / PDF 85--89,
P.} The inspected route is different. It first studies prime-power roots
of unity through \(\pi=1-\zeta\), units of the form
\((1-\zeta^i)/(1-\zeta)\), and the product expressing \(p\). Valuations
and ramification-index bounds then establish the degree and total
ramification. The proof that \(\mathbb Z[\zeta]\) is the full ring of
integers uses localization and discriminant comparisons from earlier
chapters; the general case uses composita and a discriminant criterion.

Lang therefore supplies arithmetic structure that the elementary
cyclotomic irreducibility theorem does not establish. But sending an
undergraduate to these pages merely to obtain
\([\mathbb Q(\zeta_n):\mathbb Q]=\varphi(n)\) imports unnecessary
valuation, ramification, and discriminant theory. Xi's inspected proof
is the smaller route for that purpose.

Neither inspected passage was verified as a ready-made proof of the full
integral CRT image for a chosen factorization of \(x^n-1\). Such gluing
conditions require their own argument over \(\mathbb Z\). Rational
factorization and the fact that each component is an algebraic integer
do not prove that an arbitrary tuple of component values comes from one
integer polynomial.

\textbf{Stopping tests.} Reconstruct the determinant proof of the
integrality criterion; derive the elementary cyclotomic irreducibility
argument and pinpoint why \(p\nmid n\) matters; then give the parity
obstruction for \(x^2-1\) over \(\mathbb Z\). Enter Lang when the actual
next question asks about the ring of integers, prime decomposition,
ramification, or discriminants---not just because cyclotomic notation
has appeared.

\subsection{The local substitution
decision}\label{the-local-substitution-decision}

\begin{longtable}[]{@{}
  >{\raggedright\arraybackslash}p{(\columnwidth - 4\tabcolsep) * \real{0.3333}}
  >{\raggedright\arraybackslash}p{(\columnwidth - 4\tabcolsep) * \real{0.3333}}
  >{\raggedright\arraybackslash}p{(\columnwidth - 4\tabcolsep) * \real{0.3333}}@{}}
\toprule\noalign{}
\begin{minipage}[b]{\linewidth}\raggedright
Selected goal
\end{minipage} & \begin{minipage}[b]{\linewidth}\raggedright
Smaller sequence justified by the inspected evidence
\end{minipage} & \begin{minipage}[b]{\linewidth}\raggedright
What it does not replace
\end{minipage} \\
\midrule\noalign{}
\endhead
\bottomrule\noalign{}
\endlastfoot
Common-factor detection and elimination & Xi's bounded-degree kernel
argument, followed by one product-formula proof & Degree-degeneration
care, computational elimination practice, root-location material \\
Finite-field existence, cyclicity, subfields, Frobenius & Finite
splitting-field construction; root-set field; lcm/root-bound cyclicity;
direct subfields; automorphism counting & Möbius inversion, irreducible
counts, explicit arithmetic, full Galois correspondence \\
Finite-group invariant complements & Basic representation/module
dictionary plus Xi's or Grillet's averaged projection & Character
theory, matrix-algebra classification, compact/Lie theory, explicit
examples not worked \\
General PID module structure & Xi's definitions and submodule theorem
plus a complete invariant-factor proof from the chosen graduate source &
A proof cannot be replaced by Xi's later exercise requesting it \\
Elementary integrality and cyclotomic degree & Xi's finite-module
determinant criterion plus Gauss/reduction-mod-\(p\) irreducibility &
Integral CRT gluing, ring-of-integers theorems, ramification, valuations
and discriminants \\
\end{longtable}

The practical conclusion is to keep one proof of each reusable
mechanism, solve at least one construction and one failure-case test,
and use Xi's omitted computations as a diagnostic bank. Grillet or
Hungerford can organize the next layer, but only after their local
prerequisites are distinguished from the much larger theories
surrounding the desired result.

\subsection{Inspected-page ledger for this
supplement}\label{inspected-page-ledger-for-this-supplement}

\begin{longtable}[]{@{}
  >{\raggedright\arraybackslash}p{(\columnwidth - 4\tabcolsep) * \real{0.3333}}
  >{\raggedright\arraybackslash}p{(\columnwidth - 4\tabcolsep) * \real{0.3333}}
  >{\raggedright\arraybackslash}p{(\columnwidth - 4\tabcolsep) * \real{0.3333}}@{}}
\toprule\noalign{}
\begin{minipage}[b]{\linewidth}\raggedright
Source
\end{minipage} & \begin{minipage}[b]{\linewidth}\raggedright
New/current inspection
\end{minipage} & \begin{minipage}[b]{\linewidth}\raggedright
Scope and limitation
\end{minipage} \\
\midrule\noalign{}
\endhead
\bottomrule\noalign{}
\endlastfoot
Xi I & printed 190--193 / PDF 196--199, V/P & Entire resultant proof
sequence 7.76--7.80 and visible exercises; narrow cubic-example error
confirmed on image \\
Xi III & printed 34--35 / PDF 40--41, V/P & Cyclicity theorem and
surrounding finite-abelian context; earlier classification proof not
re-audited here \\
Xi III & printed 109--115 / PDF 115--121, V/P & Hermitian averaging,
explicit \(S_3\) example, both Maschke proofs, following remarks \\
Xi III & printed 186--196 / PDF 192--202, V/P & Module definitions
through PID submodules, structure statement, integrality proofs,
exercises \\
Xi III & printed 260--266 / PDF 266--272, V/P & Finite fields,
subfields, Frobenius, Möbius inversion, cyclotomic construction and
irreducible counts; later §5.6 exercise pages not audited \\
Xi III & printed 270--272 / PDF 276--278, V/P & Roots-of-unity theorem,
cyclotomic irreducibility proof including its final lines, exercises \\
Grillet & printed 157 / PDF 169, P & Cyclic finite multiplicative
subgroup proof \\
Grillet & printed 176--179 / PDF 188--191, P & Resultant definition,
universal determinant proof, discriminant identity \\
Grillet & printed 192--193 / PDF 204--205, P & Finite-field proof and
exact exercise status of subfield results \\
Grillet & printed 197--200 / PDF 209--212, P & General Galois framework
sampled to identify prerequisites; not used as a shorter proof
recommendation \\
Grillet & printed 381--386 / PDF 393--398, P & Group algebra,
representation/module dictionary, Maschke, adjacent classification and
cyclic-group exercises \\
Hungerford & printed 279--282 / PDF 298--301, P & Finite-field proofs,
explicit imports and exercises \\
Hungerford & printed 453--456 / PDF 472--475, P & Maschke radical/trace
proof and surrounding algebra-structure context \\
Lang & printed 71--75 / PDF 85--89, P & Cyclotomic degree, ramification
and integer-ring proof dependencies; formulas whose extraction was
uncertain were not transcribed \\
\end{longtable}

Public prose contains original paraphrases and small mathematical
reconstructions only. Source scans, raw excerpts, private Library
identifiers and access metadata are outside the repository. No commit or
push was performed.

\section{Source ledger and page
conventions}\label{source-ledger-and-page-conventions}

\subsection{Evidence levels}\label{evidence-levels}

The expanded edition adds contiguous proof readings documented in
\href{PROOF-COMPARISON-OPERATORS.md}{the operator supplement} and
\href{PROOF-COMPARISON-ARITHMETIC.md}{the arithmetic supplement}. Their
ending ledgers give the new page ranges and distinguish extracted proof
reading from rendered-page verification. The original selected readings
below remain part of the baseline evidence, not a claim to have read
every intervening page.

\textbf{P} means the specified text, theorem, construction, or proof was
inspected in this task. \textbf{V} means a rendered source page was
inspected, necessary for scans or a suspected extraction problem.
\textbf{T} means a contents or introductory description was inspected;
it establishes a location and stated scope, not that every theorem or
exercise was audited. \textbf{C} means a theorem is cited as a standard
result and is not fully reproved or source-audited here. The comparison
uses these distinctions rather than treating an index hit as proof
coverage.

PDF page numbers are \textbf{one-based physical page positions in the
supplied file}; printed pages are the numbers visible on the book pages.
The offsets below apply to the ordinary Arabic-numbered body at the
checked locations, not to Roman frontmatter. A section range is a
reading destination; a list of inspected pages is the actual evidence,
and should not be expanded into a full-book audit.

\subsection{Editions and physical
files}\label{editions-and-physical-files}

\begin{longtable}[]{@{}
  >{\raggedright\arraybackslash}p{(\columnwidth - 6\tabcolsep) * \real{0.2308}}
  >{\raggedright\arraybackslash}p{(\columnwidth - 6\tabcolsep) * \real{0.2308}}
  >{\raggedleft\arraybackslash}p{(\columnwidth - 6\tabcolsep) * \real{0.3077}}
  >{\raggedright\arraybackslash}p{(\columnwidth - 6\tabcolsep) * \real{0.2308}}@{}}
\toprule\noalign{}
\begin{minipage}[b]{\linewidth}\raggedright
Key
\end{minipage} & \begin{minipage}[b]{\linewidth}\raggedright
Supplied source
\end{minipage} & \begin{minipage}[b]{\linewidth}\raggedleft
File pages
\end{minipage} & \begin{minipage}[b]{\linewidth}\raggedright
Body page conversion
\end{minipage} \\
\midrule\noalign{}
\endhead
\bottomrule\noalign{}
\endlastfoot
Xi I & \CJKtext{席南华}, \CJKtext{基础代数} \CJKtext{第一卷}, Science
Press, 2016 & 207, local metadata & PDF = printed + 6 \\
Xi II & \CJKtext{席南华}, \CJKtext{基础代数} \CJKtext{第二卷}, Science
Press, 2018 & 334, local metadata & PDF = printed + 6 \\
Xi III & \CJKtext{席南华}, \CJKtext{基础代数} \CJKtext{第三卷}, Science
Press, 2021, second printing visible & 361, local metadata & PDF =
printed + 6 \\
Hungerford & Thomas W. Hungerford, Algebra, Springer edition/reprint of
the 1974 text; 1980 Springer preface and 2003 printing note & 523, local
metadata & PDF = printed + 19 \\
Grillet & Pierre Antoine Grillet, Abstract Algebra, second edition, 2007
& 684, local metadata & PDF = printed + 12 \\
Mac Lane & Saunders Mac Lane, Categories for the Working Mathematician,
second edition & 331, local metadata & PDF = printed + 13 \\
Hilton--Stammbach & Peter J. Hilton and Urs Stammbach, A Course in
Homological Algebra, second edition & 376, local metadata & PDF =
printed + 11 \\
Hall & Brian C. Hall, Lie Groups, Lie Algebras, and Representations,
second edition, 2015 & 452, local metadata & PDF = printed + 13 \\
Hartshorne & Robin Hartshorne, Algebraic Geometry, Springer & 511, local
metadata & PDF = printed + 15 \\
Lang & Serge Lang, Algebraic Number Theory, second edition; supplied
reprint identifies 1986 edition with later copyright information & 372,
local metadata & PDF = printed + 14 \\
Hirsch & Morris W. Hirsch, Differential Topology, 1976/reprint & 230
reported by Library reader; no local byte audit & PDF = printed + 10 at
checked body pages \\
\end{longtable}

Local SHA-256 fingerprints and byte sizes for the ten materialized files
are saved in \texttt{evidence/source-fingerprints.json}. They identify
which supplied editions underlie the references; they are not a
distribution of source bytes. The original Library filenames, which
sometimes have OCR-like punctuation, are not treated as authoritative
edition statements.

\subsection{Xi I: concrete foundations and determinant
characterizations}\label{xi-i-concrete-foundations-and-determinant-characterizations}

\begin{itemize}
\tightlist
\item
  \textbf{P/T:} title/edition and preface, PDF pp.~3--4; contents,
  pp.~5--6. The preface connects the course to Kostrikin's framework and
  names several other references, including Hungerford. This supports
  pedagogical context, not a claim that the books have identical proof
  coverage.
\item
  \textbf{P:} Chapter 4 opening, printed pp.~75--76 (PDF pp.~81--82):
  oriented volume, multilinearity, and vanishing on dependent vectors.
\item
  \textbf{P:} §4.4, printed p.~93 (PDF p.~99): several characterizations
  of the determinant, with the preceding exercises including
  Cauchy--Binet/Laplace identities. This is actual structural content,
  not merely a chapter title.
\item
  \textbf{P/T:} §6.3 opening, printed p.~133 (PDF p.~139): factorization
  and examples whose irreducibility changes with the coefficient ring.
\item
  \textbf{P/T:} §7.4 opening, printed p.~177 (PDF p.~183): elementary
  symmetric polynomials and invariance under ordering roots.
\item
  \textbf{P:} §7.6 opening, printed p.~190 (PDF p.~196): resultant
  matrix and criterion, with the leading-coefficient degeneracy
  explicitly visible. This was an opening-page sample in the baseline.
  The continuation inspected the full proof sequence on PDF pp.~196--199
  and its visible exercises; see the arithmetic supplement, including
  the narrow cubic-example erratum.
\item
  Quotients of sets (§2.4, printed p.~30, PDF p.~36) and matrix-map
  organization (Chapter 3, starting p.~40, PDF p.~46) were located and
  sampled.
\end{itemize}

\subsection{Xi II: the principal source
comparison}\label{xi-ii-the-principal-source-comparison}

\begin{itemize}
\tightlist
\item
  \textbf{P:} §1.4--1.6 samples, printed pp.~16, 19--24 (PDF pp.~22,
  25--30): independent subspaces, complements versus quotients, duals
  and evaluation.
\item
  \textbf{P:} §2.4, printed pp.~67--71 (PDF pp.~73--77), selected text:
  quotient operator and induced-map diagram; characteristic-polynomial
  factorization; Cayley--Hamilton statement and dual operator
  discussion.
\item
  \textbf{P:} §2.5, printed pp.~72--77 (PDF pp.~78--83), consecutive
  proof reading: Theorem 2.51 uses Bézout to split annihilated spaces;
  2.52--2.53 gives generalized eigenspaces; 2.56--2.57 gives Jordan
  form; the proof constructs chains by quotient induction and proves
  uniqueness by dimensions of images of powers.
\item
  \textbf{P:} §3.1 opening, printed p.~85 (PDF p.~91), and §3.3 opening,
  printed p.~109 (PDF p.~115), with examples/exercises from metric
  linear algebra. \textbf{P:} Theorem 3.62 and its normal-operator
  argument, printed pp.~126--128 (PDF pp.~132--134).
\item
  \textbf{T:} Chapters 4--5, printed pp.~161--262 (PDF pp.~167--268):
  affine and Euclidean geometry, transformations, convexity, indefinite
  spaces, quadrics, projective geometry. Their complete proofs were not
  audited; no blanket replacement claim about them is made.
\item
  \textbf{P/V:} tensor Definition 6.1, printed p.~263 (PDF p.~269),
  defines all multilinear scalar functions on the indicated product
  spaces.
\item
  \textbf{P:} tensor Definition 6.11 and Theorem 6.15, printed
  pp.~273--275 (PDF pp.~279--281): basis construction, uniqueness, and
  universal property \textbf{with a proof}. Subsequent maps,
  associativity, distributivity, scalar extension, and Kronecker
  formulas were read on PDF pp.~280--284.
\item
  \textbf{P:} §6.3 examples of contraction as trace/composition and
  structural tensors, printed pp.~280--283 (PDF pp.~286--289).
\item
  \textbf{P/V:} Theorem 6.19 and following remark, printed p.~278 (PDF
  p.~284). The finite-dimensional theorem is correct; the unrestricted
  infinite-dimensional extension fails for the definition on p.~263. See
  \texttt{VERIFICATION.md} for the explicit counterexample and its
  scope.
\item
  \textbf{P:} §6.4 opening and definition, printed pp.~292--293 (PDF
  pp.~298--299): characteristic-zero assumption and exterior product via
  alternation.
\end{itemize}

\subsection{Xi III: selected scanned pages, not unreliable
extraction}\label{xi-iii-selected-scanned-pages-not-unreliable-extraction}

The Library extraction's earlier 150-page report is not the physical
file length. PyMuPDF counted \textbf{361 pages}. The scan's extracted
text is extremely sparse. The following evidence comes from rendered
pages, not OCR guesses.

\begin{itemize}
\tightlist
\item
  \textbf{V/T:} contents, PDF pp.~5--6: group actions and structure,
  one-parameter groups (§2.5, p.~79), matrix Lie groups (§2.6, p.~85),
  representations (Chapter 3, p.~97), modules (§4.4, p.~186), algebras
  (§4.5, p.~197), \(\mathfrak{sl}_2\) (§4.6, p.~215), Galois theory
  (Chapter 5, p.~224), and the representations appendix (p.~317).
\item
  \textbf{V/P:} §3.2, printed pp.~109--111 (PDF pp.~115--117): complex
  Hermitian-form convention and Theorem 3.20's explicit finite-group
  averaging proof, with the compact-group integration comment
  distinguished from it.
\item
  \textbf{V/P:} §3.7, printed pp.~152--153 (PDF pp.~158--159): dual
  action with inverse, transpose formula, and tensor action.
\item
  \textbf{V/P:} §4.4, printed pp.~186--188 (PDF pp.~192--194): module
  motivation, unital left-module definition, examples, submodules,
  quotients and maps.
\item
  \textbf{V/P:} §4.6, printed pp.~215--216 (PDF pp.~221--222): Lie
  module conditions and discussion of the enveloping algebra. The page
  explicitly sends the reader elsewhere for the PBW theory; the notes do
  not claim that this section proves PBW.
\item
  \textbf{V/P:} §5.6, printed pp.~260--262 (PDF pp.~266--268): cyclic
  multiplicative group, finite-field existence and uniqueness, subfields
  and Frobenius. This supplies the finite-field prerequisites for the
  Paley bridge.
\item
  \textbf{V/P:} §5.8, printed pp.~272--273, 275--277 (PDF pp.~278--279,
  281--283): Noether/Hilbert 90 constructions, norm/trace examples and
  trace-pairing and coefficient-formula exercises. The trace pairing on
  p.~277 is an exercise, not a full proof at that location.
\item
  \textbf{V/P:} §5.11 opening, printed p.~286 (PDF p.~292): the
  normal-basis theorem and beginning of its interpolation construction.
  The rest of that proof, full Galois computations/inverse questions,
  and the representation appendix remain selective reading destinations,
  not audited proofs here.
\end{itemize}

\subsection{Graduate sources: where the extra structure is actually
built}\label{graduate-sources-where-the-extra-structure-is-actually-built}

\textbf{Grillet.} Title/preface and contents were inspected. The
detailed samples include VIII.6--7 (printed pp.~336--344, PDF
pp.~348--350, 354--356), X.1 and X.3 (pp.~393--395, 401--402; PDF
pp.~405--407, 413--414), XI.5--7 (pp.~434--437, 441--443, 447--448; PDF
pp.~446--449, 453--455, 459--460), XII.4--5 (pp.~487--489, 493--494; PDF
pp.~499--501, 505--506), and XIII.2--4 (pp.~518--526; PDF pp.~530--531,
533--538). These support the module, splitting, balanced tensor,
Ext/Tor, and free-algebra comparisons. The general PID proof's key basis
and invariant-factor arguments were read; not every ancillary exercise
was checked. IX.7, printed pp.~380--385 (PDF pp.~392--397), was
inspected for the group-algebra correspondence, Maschke proof, and its
nonstandard initial use of ``irreducible'' for indecomposable. The
p.~381 definition was checked as an image. IV.1.6 (printed p.~157, PDF
p.~169) was read for cyclic multiplicative subgroups, and V.1.3 (printed
pp.~192--193, PDF pp.~204--205) for the finite-field
existence/uniqueness proof and field-table exercises. The other IV--VI
destinations for norm/trace, orders/valuations and completions are
\textbf{T}, not a proof audit.

\textbf{Hungerford.} Prefaces, contents, and course suggestions were
read, especially printed p.~xviii (PDF p.~15), which explicitly places
torsion PID modules before VII.1--5 in the proposed linear-algebra
course. IV.6 begins at printed p.~218 (PDF p.~237); selected text on
torsion/cyclic modules and the structure statements was read at PDF
pp.~239--241. VII.4 starts at p.~355 (PDF p.~374), VII.5 at p.~366 (PDF
p.~385). The route and theorem locations are verified, but a
line-by-line comparison of all proofs in VII was not completed. IV.4
(p.~199, PDF p.~218), IV.5 (p.~207, PDF p.~226), are selective
destinations grounded in contents. V.5's finite-field proofs were then
read at printed pp.~279--281 (PDF pp.~298--300), including Theorem 5.3,
Proposition 5.6, Corollary 5.7, Frobenius, and the explicit exercise
constructing fields of orders 9 and 25.

\textbf{Mac Lane.} III.1 (printed pp.~55--58, PDF pp.~68--71), III.2
(pp.~60--62, PDF pp.~73--75), IV.1 (pp.~79--80, PDF pp.~92--93), and
IV.2's duality example (pp.~88--89, PDF pp.~101--102) were read. These
are actual definitions, examples, and proofs of universality,
representability/Yoneda, and adjunction. VIII's abelian-category
treatment and later monads/Kan extensions are \textbf{T/C} in this
study.

\textbf{Hilton--Stammbach.} Edition, contents, and introduction were
read; III.1--2 and III.4 were sampled in detail (printed pp.~84--86,
89--90, 97--98; PDF pp.~95--97, 100--101, 108--109). IV.1's definitions
and homology construction were read at pp.~117--118 (PDF pp.~128--129).
The introduction explicitly motivates the extension-first route before
general derived functors. The complete comparison theorem for
resolutions and the full general Ext classification are \textbf{C} here;
our one-step cyclic cases are proved directly.

\textbf{Hall.} The supplied book is the second edition. §§4.3--4.7 were
sampled at printed pp.~90--101 (PDF pp.~103--114), including the
unipotent nonsemisimple example, unitary splitting proof, Schur
argument, and \(\mathfrak{sl}_2\) classification proof. Later root
systems, PBW, highest-weight theory, and topology are \textbf{T}, not
compared proof by proof against Xi.

\textbf{Hartshorne.} Contents and introduction, Chapter II opening,
II.5, II.8, and Chapter III opening were read at PDF pp.~9--10, 14--15,
74--75, 123--126, 187--189, 216. In printed pages these include II.5 at
pp.~108--111, II.8 at pp.~172--174, and III's opening at p.~201. The
algebraic differential universal property and the rational-residue-field
cotangent calculation are visible at these locations. The ``A'' notation
distinguishes externally referenced theorems. Affine finite projectivity
and general cohomological theorems are cited rather than re-audited in
full.

\textbf{Lang.} Title/edition and contents were read. I.1 opening
(printed p.~3, PDF p.~17), I.9 Proposition 26 and its local-to-global
proof (p.~29, PDF p.~43), and III.1's complementary-module
definition/proposition (p.~57, PDF p.~71) were inspected. Full
valuation, completion, ramification and class field arguments remain
\textbf{T/C}. A route to them is identified without claiming that the
current notes replace them.

\textbf{Hirsch.} Located by title in Library. Preface/contents, PDF
pp.~8--10; Chapter IV.1--2 samples, PDF pp.~96--105 (printed
pp.~86--95). The preface's declared avoidance of tensors/forms and
IV.2.2's partition-of-unity splitting proof are \textbf{P}. Local
download via the Library helper returned HTTP 401; no local
physical-file metadata or fingerprint is claimed. The read-only Library
content route succeeded. The full transversality/Morse proofs are not
audited here.

\subsection{Existing Hadamard/Williamson
research}\label{existing-hadamardwilliamson-research}

The reference edition is pinned to repository commit
\texttt{2834fea31fb4ca7203531ca220e55a2fb7640119}, not a moving branch
tip:

\begin{itemize}
\tightlist
\item
  \href{https://github.com/Drsunflowerherbert/Codex-for-Math/blob/2834fea31fb4ca7203531ca220e55a2fb7640119/hadamard-study/README.md}{Study
  and original scope}.
\item
  \href{https://github.com/Drsunflowerherbert/Codex-for-Math/blob/2834fea31fb4ca7203531ca220e55a2fb7640119/hadamard-study/PROOFS.md}{Original
  constructions and exact sign conventions}.
\item
  \href{https://github.com/Drsunflowerherbert/Codex-for-Math/blob/2834fea31fb4ca7203531ca220e55a2fb7640119/hadamard-study/round2/PROOFS.md}{Second-round
  structural and spectral proofs}.
\item
  \href{https://github.com/Drsunflowerherbert/Codex-for-Math/blob/2834fea31fb4ca7203531ca220e55a2fb7640119/hadamard-study/round3/PROOFS.md}{Third-round
  proofs, especially §§1--4, 9, 11--13}.
\item
  \href{https://github.com/Drsunflowerherbert/Codex-for-Math/blob/2834fea31fb4ca7203531ca220e55a2fb7640119/hadamard-study/round3/LESSONS.md}{Third-round
  limitations}.
\item
  \href{https://github.com/Drsunflowerherbert/Codex-for-Math/blob/2834fea31fb4ca7203531ca220e55a2fb7640119/hadamard-study/round3/reviews/REVIEW.md}{Prior
  independent adversarial review}.
\end{itemize}

Those documents and the relevant exact fixture descriptions were read
from Git objects. The new foundation chapter gives its own derivations,
identifies where it extends the scope of the earlier prime Paley proof,
and retains the distinction between previous audit evidence and newly
executed finite checks. Numerical historical/status claims and whole
prior benchmark campaigns were not rerun for this foundational task.

\section{Verification record, corrections, and
limits}\label{verification-record-corrections-and-limits}

\subsection{What the evidence means}\label{what-the-evidence-means}

The expanded edition adds four chapters and two proof comparisons to the
corrected baseline. Separate AI agents reviewed the new chapter proofs,
and independent exact calculations checked their most delicate examples.
The \href{PHASE2-REVIEW.md}{continuation review} records the findings
and corrections. The additional runner \texttt{scripts/check\_phase2.py}
records 28 chapter check groups and one symbolic cubic-discriminant
check in \texttt{evidence/phase2-checks.json}, with hashes of the
executed scripts. These supplement the baseline's 18 groups; they are
not formal verification.

The notes contain ordinary written proofs and explicitly cited
structural theorems. They have undergone a separate self-review for
hypotheses, notation, and failure cases, plus exact computational spot
checks written independently of the earlier repository search
implementations. There is no claim of a separate human reviewer,
Lean/kernel verification, a full-book audit, or a new result about
universal Hadamard existence.

\texttt{evidence/exact-checks.json} records the actually executed finite
checks, software versions and their scope.
\texttt{evidence/pdf-build.json} and the LaTeX logs record compilation;
compilation verifies neither mathematical correctness nor source
fidelity. Source fidelity is supported by the specific inspected pages
listed in the source ledger, including image inspection where required.

\subsection{Central statements checked by written
derivation}\label{central-statements-checked-by-written-derivation}

\begin{longtable}[]{@{}
  >{\raggedright\arraybackslash}p{(\columnwidth - 2\tabcolsep) * \real{0.5000}}
  >{\raggedright\arraybackslash}p{(\columnwidth - 2\tabcolsep) * \real{0.5000}}@{}}
\toprule\noalign{}
\begin{minipage}[b]{\linewidth}\raggedright
Statement
\end{minipage} & \begin{minipage}[b]{\linewidth}\raggedright
Hypotheses and check
\end{minipage} \\
\midrule\noalign{}
\endhead
\bottomrule\noalign{}
\endlastfoot
Quotient factorization and first isomorphism theorem & Maps, domains,
kernels and representative independence are explicit; no complement is
assumed \\
Dual exactness and annihilator identities & Algebraic dual; finite
dimension for the double-dual isomorphism and dimension arguments; basis
extension for restriction surjectivity \\
Direct sums and splitting & Idempotent equations proved for modules;
vector-space splitting distinguished from module/equivariant splitting;
sections form an affine Hom space \\
Operator/module correspondence & Intertwiner equation checked;
finite-dimensional operator modules distinguished from all finitely
generated \(k[t]\)-modules \\
Canonical forms & Nonzero annihilator obtained without Cayley--Hamilton;
Euclidean Smith mechanism explained; arbitrary PID version cited;
splitting-field hypothesis retained for Jordan form \\
Polynomial projectors & Congruences, image/kernel identification, and
independence from Bézout representative checked; canonical primary
spaces distinguished from chosen cyclic summands \\
Cyclic Hom and commutant dimensions & Computed from the image of \(1\)
and divisibility; no unjustified extension to two-variable polynomial
rings \\
Hom--tensor formulas & Currying needs no finite dimension; the canonical
tensor map identifies its domain with finite-rank maps, and finite
dimension of either source or target suffices for all Hom;
noncommutative module sides specified \\
Symmetric/exterior constructions & Quotient universal properties;
alternation includes repeated-vector vanishing in characteristic \(2\);
averaging needs invertible factorial \\
Determinants and minors & Exterior basis proof precedes determinant
definition; rank test, Cauchy--Binet and characteristic coefficients
derived; determinant of a map between different spaces is a line map \\
Forms and spectral theorem & Congruence versus similarity; Hermitian
convention explicit; positive-definite complex finite-dimensional
hypotheses retained; real and indefinite failures exhibited \\
Ext/Tor and homology & One-step resolution requires a nonzero element of
a domain; explicit nonsplit \(k[t]\) extension; field Ext/Tor vanishing
separated from arbitrary complex homology \\
Representations & Maschke needs finite group and invertible order;
Schur's scalar result needs algebraic closure and finite dimension; dual
action has inverse; Lie tensor action has a sum \\
Trace pairing & Finite separable extension for nondegeneracy;
scalar-extension/CRT proof and purely inseparable counterexample
included \\
Bundles and sheaves & Smooth paracompact bundle splitting distinguished
from arbitrary pointwise choices and algebraic sheaf extensions; no
category-independent Ext vanishing claim \\
Hadamard bridge & Integer/real Gram equations; odd finite fields for
character sum; exact Paley signs; cyclic Fourier
splitting/characteristic hypotheses; compression and phase relaxation
retain only necessary implications \\
\end{longtable}

The proofs are inspectable in the chapters. The table is a review index,
not an independent formal certificate.

\subsection{Independently computed finite
checks}\label{independently-computed-finite-checks}

The script \texttt{scripts/check\_examples.py} reconstructs fixtures
directly and does not import the previous Hadamard project's search or
verification code. Its 18 check groups include:

\begin{itemize}
\tightlist
\item
  Polynomial CRT idempotents over \(\mathbb Q\) and characteristics
  \(2,3,5,7\); a conjugated mixed Jordan fixture checks the actual
  operator projections.
\item
  Kernel-growth block reconstruction, and the commutant dimension
  independently computed as the nullity of a \(100\times100\) commutator
  equation matrix. A real nonsplitting \(t^2+1\) summand preserves the
  kernel data while increasing the dimension from ten to twelve,
  exhibiting the missing-hypothesis failure excluded by the corrected
  exercise.
\item
  Sixteen cyclic-Hom dimension comparisons using direct intertwiner
  equations rather than the gcd formula being checked.
\item
  Exterior minors, rectangular Cauchy--Binet, and characteristic
  coefficients; Kronecker determinant/trace; adjoints for a
  nonorthonormal Hermitian form.
\item
  Explicit module-section failure over four small prime fields; a
  nonzero field-complex homology example with chain-homotopy equations;
  the \(\mathfrak{sl}_2\) polynomial commutators for highest weights
  zero through six.
\item
  Quadratic-character sums and both Paley block constructions for six
  small primes, plus an actual nonprime-field fixture
  \(\mathbb F_9=\mathbb F_3[u]/(u^2+1)\) producing order 20.
\item
  Williamson block signs at core three, the core-nine proper-quotient
  false positive, and all 27 possible normalized symmetric lifts of the
  \textbf{fixed} stronger compressed multiset from the prior audit. None
  is a Williamson lift. This is a bounded adversarial-fixture replay,
  not an expanded search campaign.
\item
  The core-seven positive spectral residual with no fourth Williamson
  row; all 16 two-color-bit equality inputs and the exact Boolean phase
  identity. The explicit edge-orbit correlation formula is also compared
  with direct integer correlations for all 1,364 four-color labelings at
  cores \(3,5,7,9,11\), using three sign fixtures per labeling (4,092
  assignments), including the parity-forced signs. This is a bounded
  identity check, not a search for or classification of solutions.
\end{itemize}

Finite checks can catch sign, transpose, coefficient and characteristic
errors in these examples. They do not prove the universal theorems or
classify all Hadamard matrices. The infinite-dimensional tensor
counterexample is proved by unbounded rank, not established by a finite
truncation test.

\subsection{Corrections incorporated in the local
continuation}\label{corrections-incorporated-in-the-local-continuation}

The continuation review supplied six focused corrections. Exercise 4 now
specifies \(\dim_kV=10\), and the chapter's parallel example identifies
only the \(\lambda\)-primary blocks until a total dimension is known.
Exercise 6, its answer, and the Hom/tensor summary use the exact
finite-rank image criterion; finite dimension of either space suffices.
Maschke's projection is explicitly an endomorphism of \(V\). The second
congruence example begins ``Similarly.'' Contraction is defined for
\(r\ge1\) and set to zero on \(\Lambda^0V\). Finally, \(E_t\) is defined
from the involution orbits, with the zero pair and fixed point removed
before passing to opposite-index classes. The original source findings
below are retained. That focused six-fix revision did not rematerialize
or independently reread the source books; the later substantive
expansion did inspect additional proof passages, as documented in the
two new source-comparison ledgers.

\subsection{Cited results, distinguished from complete proofs
here}\label{cited-results-distinguished-from-complete-proofs-here}

\begin{itemize}
\tightlist
\item
  The full finitely generated module structure theorem over an arbitrary
  PID (Grillet VIII.6; Hungerford IV.6). The notes give the Euclidean
  matrix construction and derive the operator consequences; they do not
  pretend every PID is Euclidean.
\item
  General derived-functor theory, independence of resolutions, and the
  complete Ext/extension correspondence (Hilton--Stammbach III--IV;
  Grillet XII). The cyclic class calculation and its extension
  interpretation are proved directly in the notes.
\item
  Exhaustion of finite-dimensional complex \(\mathfrak{sl}_2\)
  irreducibles (Hall §4.6). The polynomial models and their
  irreducibility are proved here; the exhaustion proof was read but is
  cited rather than reproduced fully.
\item
  Primitive-element theorem for finite separable extensions;
  finite-field existence; the general commutative-algebra/sheaf
  correspondence; finite projectivity on affine schemes;
  implicit-function, partition-of-unity and other geometric
  prerequisites. Their uses and sources are stated, without implying all
  are reproved in this study.
\item
  Prior study benchmarks, search completeness records, and historical
  matrix witnesses remain attributed to their pinned source edition.
  Only the small fixtures listed above were newly executed. No current
  worldwide open-problem status or smallest unresolved order is
  certified.
\end{itemize}

\subsection{Confirmed source
corrections}\label{confirmed-source-corrections}

\textbf{Xi I, Example 7.80, printed p.~193 / PDF p.~199.} The supplied
scan prints \(-2p^3-27q^2\) for the discriminant of \(x^3+px+q\). The
correct coefficient is \(-4\). Two independent image inspections
confirmed the displayed formula; the preceding general cubic exercise on
printed p.~190 / PDF p.~196 has the correct term. Direct symbolic
calculation gives \(\operatorname{Res}(f,f')=4p^3+27q^2\), and the monic
cubic discriminant is its negative. The specialization \(f=x^3+x\) gives
\(-4\) directly. This is a narrow example typo, not a defect in the
resultant criterion or its inspected proof.

\textbf{Xi II, Theorem 6.19's following remark, printed p.~278 / PDF
p.~284.} The remark extends the finite-dimensional identification
\(\mathcal T_q^p(V)\cong(V^*)^{\otimes p}\otimes V^{\otimes q}\) to
infinite-dimensional spaces. Definition 6.1 on printed p.~263 / PDF
p.~269 defines the left side as all multilinear functions on
\(V^p\times(V^*)^q\). Both pages were rendered and inspected.

For \(p=q=1\), evaluation \((v,\lambda)\mapsto\lambda(v)\) belongs to
that function space. If it were a finite tensor sum, there would be
\(\alpha_i\in V^*\) and \(w_i\in V\) such that

\[\lambda(v)=\sum_{i=1}^r\alpha_i(v)\lambda(w_i)
\qquad\text{for every }v,\lambda.\]

Functionals separate vectors, so \(v=\sum_i\alpha_i(v)w_i\) for every
\(v\). This puts an infinite-dimensional space in a finite-dimensional
span, a contradiction. The corrected statement retains
finite-dimensionality, or restricts the multilinear function space to
the image of finite tensor sums. The finite-dimensional theorem and the
tensor universal property are unaffected. This is one local confirmed
overstatement, not a blanket judgment of the text.

\subsection{Convention differences and missing hypotheses, not source
mistakes}\label{convention-differences-and-missing-hypotheses-not-source-mistakes}

\begin{itemize}
\tightlist
\item
  Grillet IX.7's initial ``irreducible'' means indecomposable (printed
  pp.~381, 383; PDF pp.~393, 395). It differs from the simple-module
  convention in these notes, Xi and Hall. The apparent unconditional
  complete-reducibility statement must be translated using his
  definition, not declared false.
\item
  Xi's Hermitian convention is linear in the first argument; these notes
  use conjugate-linearity in the first. Adjoint and coordinate formulas
  are translated explicitly.
\item
  Xi II's alternating-tensor model explicitly assumes characteristic
  zero. Its factorial denominators are valid there; moving the model to
  arbitrary characteristic without that assumption would be our error,
  not its theorem.
\item
  A Jordan form needs a splitting minimal polynomial; semisimple
  \(k[t]\)-modules need square-free irreducible factors and can lack
  eigenvectors over the base field. Commuting operators require
  individual diagonalizability before the simultaneous eigenbasis
  argument.
\item
  A positive spectral residual and individually feasible compressed rows
  need not have one common discrete lift. No linear splitting theorem
  supplies it.
\item
  Fourier eigenvalues over \(\mathbb C\) are exact algebraic values in
  the proof; floating-point approximations are not used to certify sign
  matrices.
\end{itemize}

\subsection{OCR issues and unverified
comparisons}\label{ocr-issues-and-unverified-comparisons}

Xi III's local PDF has 361 pages, whereas its earlier Library extraction
reported 150. Sparse extraction and corrupted symbols are source-access
limitations, not mathematical errors. The key cited pages were read as
images. Xi I--II also have corrupted extracted signs and indices, so
suspected substantive claims were checked against the intended notation
and, where material, images. The exact source fingerprints allow later
reinspection.

The following comparisons remain deliberately limited: the entire
exercises of any two books; Xi II Chapters 4--5 versus a putative
replacement; the whole Xi III appendix or Galois computations against
Hall/Lang; complete Hungerford VII proof alternatives; and the
destination books' full advanced prerequisite coverage. They are not
promoted from contents evidence into proven replacement claims. No
source book is called filler or wholly replaceable here.

\subsection{Remaining access
limitation}\label{remaining-access-limitation}

The exact failed action was local materialization of the Library item
titled \texttt{Differential\ Topology\ (Morris\ W.\ Hirsch).pdf} using
the provided download helper. It returned
\texttt{hosted\ apps\ tools/list\ request\ failed\ with\ HTTP\ status\ 401}.
This is a failed local-download action, not a successful or uncertain
transfer and not an automatic sandbox-approval denial. The same source's
read-only Library pages were successfully read afterward. Consequently
the Hirsch content needed here is sourced, but no physical local PDF
page count, byte hash, or full local-PDF audit is claimed.

There is no pending mathematical/source blocker to using the present
notes. Completing the optional unverified comparisons would be a new,
larger audit; recovering a local Hirsch copy is only necessary for that
sort of byte-level or full visual audit.

\end{document}
